% !TEX program = xelatex
\documentclass[11pt,a4paper,twoside]{report}

\usepackage{xcolor}
\usepackage[a4paper,margin=28mm]{geometry}
\usepackage{graphicx}
\usepackage{microtype}
\usepackage{setspace}
\usepackage{parskip}
\usepackage{booktabs}
\usepackage{tabularx}
\usepackage{longtable}
\usepackage{array}
\usepackage{multirow}
\usepackage{enumitem}
\usepackage{amssymb}
\usepackage{amsmath}
\usepackage{listings}
\usepackage[hyphens]{url}
\usepackage{tcolorbox}
\usepackage{tikz}
\usepackage[table]{colortbl}
\usetikzlibrary{arrows.meta,positioning,shapes.geometric,fit,calc,chains,
                decorations.pathmorphing,backgrounds}

\newtcolorbox{adrbox}[2][]{colback=blue!5,colframe=blue!75!black,fonttitle=\bfseries,title=ADR: #2,#1}
\newtcolorbox{hitlbox}[2][]{colback=green!5,colframe=green!50!black,fonttitle=\bfseries,title=HITL Checkpoint: #2,#1}
\newtcolorbox{gapbox}[2][]{colback=red!5,colframe=red!50!black,fonttitle=\bfseries,title=Gap: #2,#1}

% --- Semantic Hyperlinking --------------------------------------------------
\usepackage{xr-hyper}
\usepackage[hidelinks]{hyperref}

% External references to peer specifications
\externaldocument[tb:]{../tb/tb-review-spec}
\externaldocument[reg:]{../regulations/regulations-spec}
\externaldocument[ar:]{../arabic/arabic-ap-spec}
\externaldocument[sim:]{../simula/simula-training-spec}

% --- Fonts and listings -----------------------------------------------------
\usepackage{fontspec}
\setmainfont{Arial}
\usepackage[capitalize,noabbrev]{cleveref}

\definecolor{codebg}{gray}{0.96}
\definecolor{sapblue}{HTML}{0B5FA5}
\definecolor{sapgreen}{HTML}{00695C}
\definecolor{sapamber}{HTML}{B8360F}
\definecolor{sapgrey}{HTML}{546E7A}
\definecolor{sapmidnight}{HTML}{1C2B39}
\definecolor{sapgold}{HTML}{F0AB00}
\definecolor{sourcegreen}{HTML}{2E7D32}
\definecolor{reqblue}{HTML}{1565C0}
\definecolor{specamber}{HTML}{B8360F}
\definecolor{registrymidnight}{HTML}{1C2B39}
\definecolor{reqpurple}{HTML}{6A1B9A}
\definecolor{specblue}{HTML}{1565C0}
\definecolor{schemaorg}{HTML}{B8360F}

% --- Custom Column Types ----------------------------------------------------
\newcolumntype{L}[1]{>{\raggedright\arraybackslash}p{#1}}
\newcolumntype{C}[1]{>{\centering\arraybackslash}p{#1}}
\newcolumntype{R}[1]{>{\raggedleft\arraybackslash}p{#1}}

\lstdefinelanguage{jsonx}{
  basicstyle=\ttfamily\small,
  columns=fullflexible,
  breaklines=true,
  showstringspaces=false,
  comment=[l]{//},
  morestring=[b]",
  stringstyle=\color{sapgreen},
  keywordstyle=\color{sapblue},
}

\onehalfspacing
\setlist{topsep=2pt,itemsep=1pt,parsep=0pt}
\setlength{\emergencystretch}{3em}
\tolerance=1000
\hbadness=10000

% --- Global Command Overrides -----------------------------------------------
\newcommand{\corpus}{}
\newcommand{\tool}{}
\newcommand{\stateid}[1]{\textsc{\lowercase{#1}}}
\newcommand{\ar}[1]{#1}

\newcommand{\srcref}[1]{\texttt{#1}}
\newcommand{\reqref}[1]{\texttt{#1}}
\newcommand{\specref}[1]{\texttt{#1}}
\newcommand{\schemaref}[1]{\texttt{#1}}
\newcommand{\gapref}[1]{\texttt{#1}}

% --- Glossary Helper Commands -----------------------------------------------
\definecolor{pathcolor}{HTML}{004D40}
\newcommand{\schemapath}[1]{\texttt{\textcolor{pathcolor}{#1}}}
\newcommand{\enumval}[1]{\texttt{\textcolor{sapgreen}{#1}}}
\newcommand{\fieldref}[1]{\texttt{\textcolor{sapblue}{#1}}}

\begin{document}

\title{\Huge Trial Balance Review \\
       \textbf{Business Requirements \& Lineage} \\
       \large HITL Traceability Specification \\
       \small Ref: \texttt{TB-REQ-LIN-2026-v1.0}}
\author{SAP-OSS Financial Architecture Team}
\date{April 2026}

\maketitle

% --- Incorporated Executive Summary ---


\chapter*{Abstract}
\addcontentsline{toc}{chapter}{Abstract}

This specification defines the Human-in-the-Loop (HITL) Requirements
Traceability framework for Standard Chartered Bank's Trial Balance (TB)
review process. It provides the complete audit trail from source documents
(business cases, detailed operating instructions, BPMN process maps, and
workbooks) through machine-readable extraction, structured business
requirements, formal specifications, and validating JSON schemas.

The document serves as the \textbf{human validation checkpoint} for the
automated requirements extraction and traceability process, enabling
reviewers to verify the accuracy of automated extractions and establish
a clear audit trail from source to implementation.

Four orthogonal workstreams are traced in full:

\begin{itemize}
  \item a \textbf{source document inventory} cataloguing 6 source files
        totalling 137~MB, including business cases, DOI, BPMN process map,
        and HKG TB/PL review workbooks;
  \item an \textbf{extraction layer} of 7 machine-readable outputs
        (markdown, JSON, YAML, CSV) validated against YAML frontmatter
        schemas;
  \item a \textbf{requirements layer} of 27 business requirements
        categorised across executive summary, scope, process steps,
        AI augmentation, and BSS risk assessment;
  \item a \textbf{specification and schema layer} comprising 21
        specification elements and 6 JSON schemas with field-level
        traceability back to requirements.
\end{itemize}

\paragraph{Reading order.} \Cref{ch:overview} introduces the HITL framework
and traceability model; \cref{ch:source-documents} catalogues source
documents; \cref{ch:extraction-layer} covers machine-readable extraction;
\cref{ch:business-requirements} details requirement derivation;
\cref{ch:specification-mapping} maps specifications to requirements;
\cref{ch:schema-mapping} provides schema field traceability;
\cref{ch:bidirectional-traceability} presents the full forward and
backward matrices; \cref{ch:gap-analysis} identifies traceability gaps;
\cref{ch:human-review-signoff} captures review checkpoints and approvals.

\paragraph{Conventions.} Source document references appear as
\srcref{TB-BC-001}. Business requirement IDs appear as \reqref{TB-REQ-ES01}.
Specification elements appear as \specref{TB-STEP-003}. Schema references
appear as \schemaref{variance-record}. Gap references appear as
\gapref{G-001}. File paths are typeset in \texttt{monospace} and are
relative to the repository root.

\paragraph{Companion files.} The machine-readable companion to this
specification is \texttt{docs/tb/machine-readable/hitl-traceability-manifest.yaml},
which enables automated validation of all traceability links documented
in this specification.

\tableofcontents
\chapter*{Visual Traceability Map}
\addcontentsline{toc}{chapter}{Visual Traceability Map}
\label{ch:visual-map}

This single-page reference shows the complete traceability flow from each source document
through extracted content, business requirements, specifications, and implementing schemas.

% =============================================================================
\section*{End-to-End Traceability Overview}

\begin{figure}[h!]
\centering
\begin{tikzpicture}[
    node distance=0.4cm,
    source/.style={rectangle, rounded corners, draw=sourcegreen, fill=sourcegreen!20,
                   thick, minimum width=3.8cm, minimum height=0.7cm, font=\scriptsize\bfseries,
                   text=black},
    extract/.style={rectangle, rounded corners, draw=sourcegreen!70, fill=sourcegreen!10,
                    thick, minimum width=2.8cm, minimum height=0.6cm, font=\tiny},
    req/.style={rectangle, rounded corners, draw=reqpurple, fill=reqpurple!15,
                thick, minimum width=2.2cm, minimum height=0.5cm, font=\tiny},
    spec/.style={rectangle, rounded corners, draw=specblue, fill=specblue!15,
                 thick, minimum width=2.2cm, minimum height=0.5cm, font=\tiny},
    schema/.style={rectangle, rounded corners, draw=schemaorg, fill=schemaorg!20,
                   thick, minimum width=2.5cm, minimum height=0.7cm, font=\scriptsize\bfseries},
    arrow/.style={->, thick, >=stealth, sapgrey},
    label/.style={font=\scriptsize\bfseries, text=sapgrey}
]

% === Layer Labels ===
\node[label] at (-6.5, 4) {LAYER 1};
\node[label] at (-6.5, 2.5) {LAYER 2};
\node[label] at (-6.5, 0.5) {LAYER 3};
\node[label] at (-6.5, -1.5) {LAYER 4};
\node[label] at (-6.5, -3.5) {LAYER 5};

% === Source Documents (Layer 1) ===
\node[source] (src1) at (-3, 4) {TB-BC-001};
\node[source] (src2) at (0.5, 4) {TB-BC-002};
\node[source] (src3) at (4, 4) {TB-BPMN-001};

% === Extracted Files (Layer 2) ===
\node[extract] (ext1) at (-3, 2.5) {bc-tb.md};
\node[extract] (ext2) at (0.5, 2.5) {bc-bss.md};
\node[extract] (ext3) at (4, 2.5) {bpmn-glo.json};

% === Requirements (Layer 3) - grouped ===
\node[req] (req-es) at (-4.5, 0.5) {ES01-04 (4)};
\node[req] (req-sc) at (-2, 0.5) {SC01-06 (6)};
\node[req] (req-ai) at (0.5, 0.5) {AI01-06 (6)};
\node[req] (req-ps) at (3, 0.5) {PS01-08 (8)};
\node[req] (req-bss) at (5.5, 0.5) {BSS01-03 (3)};

% === Specifications (Layer 4) - grouped ===
\node[spec] (spec-step) at (-3.5, -1.5) {TB-STEP (8)};
\node[spec] (spec-dp) at (-0.5, -1.5) {DP-001-005 (5)};
\node[spec] (spec-ai) at (2.5, -1.5) {AI-* (6)};
\node[spec] (spec-ctrl) at (5.5, -1.5) {TB-CTRL (3)};

% === Schemas (Layer 5) ===
\node[schema] (sch1) at (-3, -3.5) {variance-record};
\node[schema] (sch2) at (0.5, -3.5) {commentary-draft};
\node[schema] (sch3) at (4, -3.5) {tb/pl-extract};

% === Arrows Layer 1→2 ===
\draw[arrow, sourcegreen] (src1) -- (ext1);
\draw[arrow, sourcegreen] (src2) -- (ext2);
\draw[arrow, sourcegreen] (src3) -- (ext3);

% === Arrows Layer 2→3 ===
\draw[arrow, sourcegreen!70] (ext1) -- (req-es);
\draw[arrow, sourcegreen!70] (ext1) -- (req-sc);
\draw[arrow, sourcegreen!70] (ext1) -- (req-ai);
\draw[arrow, sourcegreen!70] (ext2) -- (req-bss);
\draw[arrow, sourcegreen!70] (ext3) -- (req-ps);

% === Arrows Layer 3→4 ===
\draw[arrow, reqpurple] (req-ps) -- (spec-step);
\draw[arrow, reqpurple] (req-ps) -- (spec-dp);
\draw[arrow, reqpurple] (req-ai) -- (spec-ai);
\draw[arrow, reqpurple] (req-es) -- (spec-ctrl);
\draw[arrow, reqpurple] (req-bss) -- (spec-dp);

% === Arrows Layer 4→5 ===
\draw[arrow, specblue] (spec-step) -- (sch1);
\draw[arrow, specblue] (spec-step) -- (sch3);
\draw[arrow, specblue] (spec-dp) -- (sch1);
\draw[arrow, specblue] (spec-ai) -- (sch1);
\draw[arrow, specblue] (spec-ai) -- (sch2);

\end{tikzpicture}
\caption{Complete Five-Layer Traceability Flow}
\label{fig:visual-traceability-map}
\end{figure}

% =============================================================================
\section*{Traceability Quick Reference}

\begin{table}[h]
\centering
\small
\caption{Source Document to Schema Quick Reference}
\begin{tabularx}{\textwidth}{@{}L{2.2cm}L{3cm}L{2.5cm}L{2.5cm}L{3cm}@{}}
\toprule
\textbf{Source} & \textbf{Extracted} & \textbf{Requirements} & \textbf{Specs} & \textbf{Schemas} \\
\midrule
TB-BC-001 & bc-tb.md & ES01-04, SC01-06, AI01-06 & AI-*, TB-CTRL-* & variance-record, commentary-draft \\
TB-BC-002 & bc-bss.md & BSS-01-03 & DP-005 & variance-record \\
TB-BPMN-001 & bpmn-glo.json & PS01-08 & TB-STEP-*, DP-001-004 & variance-record, tb-extract \\
TB-DOI-001 & doi-tb.md & (supplementary) & -- & -- \\
TB-WB-001 & hkg-tb-schema.yaml & (data reference) & -- & tb-extract \\
TB-WB-002 & hkg-pl-schema.yaml & (data reference) & -- & pl-extract \\
\bottomrule
\end{tabularx}
\end{table}

% =============================================================================
\section*{Coverage at a Glance}

\begin{table}[h]
\centering
\begin{tabular}{@{}lcccl@{}}
\toprule
\textbf{Layer} & \textbf{Count} & \textbf{Traced} & \textbf{Coverage} & \textbf{Status} \\
\midrule
Source Documents & 6 & 6 & 100\% & \textcolor{sourcegreen}{\ding{51} Complete} \\
Extracted Files & 7 & 7 & 100\% & \textcolor{sourcegreen}{\ding{51} Complete} \\
Requirements & 27 & 27 & 100\% & \textcolor{sourcegreen}{\ding{51} Complete} \\
Specifications & 22 & 22 & 100\% & \textcolor{sourcegreen}{\ding{51} Complete} \\
Schemas & 6 & 6 & 100\% & \textcolor{sourcegreen}{\ding{51} Complete} \\
\midrule
\textbf{Gaps} & 2 & 2 & 100\% & \textcolor{sapamber}{\ding{51} Documented} \\
\bottomrule
\end{tabular}
\end{table}

\vspace{0.5cm}

\begin{tcolorbox}[colback=sourcegreen!5,colframe=sourcegreen!70,title=\textbf{Traceability Validation Status}]
\begin{center}
\Large\textbf{100\% BIDIRECTIONAL TRACEABILITY ACHIEVED}
\end{center}
\begin{itemize}[noitemsep]
  \item[\textcolor{sourcegreen}{\ding{51}}] Every source document section reaches at least one schema
  \item[\textcolor{sourcegreen}{\ding{51}}] Every schema field traces back to a source document
  \item[\textcolor{sourcegreen}{\ding{51}}] All 27 requirements have documented derivation evidence
  \item[\textcolor{sourcegreen}{\ding{51}}] All 22 specifications linked to requirements
  \item[\textcolor{sapamber}{\ding{51}}] 2 gaps documented with owners and resolution plans
\end{itemize}
\end{tcolorbox}
\chapter{Overview and HITL Framework}
\label{ch:overview}

This chapter introduces the Human-in-the-Loop (HITL) Requirements
Traceability framework, its purpose, and the five-layer traceability
model that underpins the Trial Balance review process documentation.

% =============================================================================
\section{Purpose Statement}
\label{sec:purpose}

The HITL Requirements Traceability specification serves as the
\textbf{human validation checkpoint} for the automated requirements
extraction and traceability process. It is designed to:

\begin{enumerate}
  \item Provide full transparency into how business requirements were
        derived from source documents
  \item Enable human reviewers to verify the accuracy of automated
        extractions
  \item Establish a clear audit trail from source to implementation
  \item Support regulatory compliance by demonstrating requirement
        provenance
  \item Facilitate change impact analysis when source documents are
        updated
\end{enumerate}

% =============================================================================
\section{Five-Layer Traceability Model}
\label{sec:five-layer-model}

The traceability framework is organised into five distinct layers,
each representing a transformation stage from raw source material
to validated implementation artefacts.

\begin{figure}[h]
\centering
\begin{tikzpicture}[
    node distance=0.5cm,
    layer/.style={rectangle, draw=sapgrey, thick, minimum width=12cm,
                  minimum height=1cm, font=\small\bfseries},
    arrow/.style={->, thick, >=stealth, sapgrey}
]
    \node[layer, fill=sourcegreen!15] (L1) {Layer 1: Source Documents (6 files, 137 MB)};
    \node[layer, fill=sourcegreen!10, below=0.3cm of L1] (L2) {Layer 2: Extracted Machine-Readable (7 files, 522 KB)};
    \node[layer, fill=reqpurple!15, below=0.3cm of L2] (L3) {Layer 3: Business Requirements (27 requirements)};
    \node[layer, fill=specblue!15, below=0.3cm of L3] (L4) {Layer 4: Specifications (21 elements)};
    \node[layer, fill=schemaorg!15, below=0.3cm of L4] (L5) {Layer 5: Schemas \& Validation (6 schemas)};
    
    \draw[arrow] (L1.south) -- (L2.north);
    \draw[arrow] (L2.south) -- (L3.north);
    \draw[arrow] (L3.south) -- (L4.north);
    \draw[arrow] (L4.south) -- (L5.north);
\end{tikzpicture}
\caption{Five-Layer Traceability Model with artefact counts}
\label{fig:five-layer-model}
\end{figure}

\subsection{Layer 1: Source Documents}

Original business documents in their native formats:
\begin{itemize}[noitemsep]
  \item Business case documents (\texttt{.doc})
  \item Detailed Operating Instructions (\texttt{.docx})
  \item BPMN process maps (\texttt{.pdf})
  \item Working spreadsheets (\texttt{.xlsm})
\end{itemize}

\subsection{Layer 2: Extracted Machine-Readable}

Automated extraction outputs in machine-processable formats:
\begin{itemize}[noitemsep]
  \item Markdown with YAML frontmatter (text content)
  \item JSON (structured BPMN elements)
  \item YAML (schemas and specifications)
  \item CSV (sample data)
\end{itemize}

\subsection{Layer 3: Business Requirements}

Structured requirements derived from source documents:
\begin{itemize}[noitemsep]
  \item Executive Summary requirements (4)
  \item Scope requirements (6)
  \item Process Step requirements (8)
  \item AI Augmentation requirements (6)
  \item BSS Risk Assessment requirements (3)
\end{itemize}

\subsection{Layer 4: Specifications}

Formal specification elements implementing requirements:
\begin{itemize}[noitemsep]
  \item Process Steps (8)
  \item Decision Points (5)
  \item AI Capabilities (6)
  \item Controls (3)
\end{itemize}

\subsection{Layer 5: Schemas \& Validation}

JSON Schemas validating specification implementation:
\begin{itemize}[noitemsep]
  \item \schemaref{variance-record} --- variance analysis records
  \item \schemaref{decision-point} --- decision point states
  \item \schemaref{entity-params} --- per-entity parameters
  \item \schemaref{tb-extract} --- trial balance extraction
  \item \schemaref{pl-extract} --- P\&L extraction
  \item \schemaref{commentary-draft} --- LLM commentary drafts
\end{itemize}

% =============================================================================
\section{Traceability Marker Reference}
\label{sec:marker-reference}

Throughout this specification, colour-coded markers indicate
traceability links across layers:

\begin{table}[h]
\centering
\caption{Traceability Marker Reference}
\label{tab:marker-reference}
\begin{tabularx}{\textwidth}{@{}C{3cm}L{5cm}L{6cm}@{}}
\toprule
\textbf{Marker} & \textbf{Meaning} & \textbf{Example} \\
\midrule
\srcref{XXX} & Source document reference & \srcref{TB-BC-001} = Trial Balance Business Case \\
\reqref{XXX} & Business requirement ID & \reqref{TB-REQ-ES01} = Executive Summary req \\
\specref{XXX} & Specification element & \specref{TB-STEP-003} = Process step \\
\schemaref{XXX} & Schema definition & \schemaref{variance-record} = JSON Schema \\
\gapref{XXX} & Identified gap & \gapref{G-001} = Traceability gap \\
\bottomrule
\end{tabularx}
\end{table}

% =============================================================================
\section{HITL Review Checkpoints}
\label{sec:review-checkpoints}

Six human review checkpoints must be validated before the traceability
chain is considered complete:

\begin{table}[h]
\centering
\caption{HITL Review Checkpoints}
\label{tab:review-checkpoints}
\begin{tabularx}{\textwidth}{@{}C{1cm}L{4cm}L{7cm}C{2cm}@{}}
\toprule
\textbf{ID} & \textbf{Checkpoint} & \textbf{Review Criteria} & \textbf{Status} \\
\midrule
H1 & Source Completeness & All source documents listed and extracted & Pending \\
H2 & Extraction Accuracy & Extracted text matches source content & Pending \\
H3 & Requirement Derivation & Requirements accurately reflect source intent & Pending \\
H4 & Traceability Links & All links are valid and bidirectional & Pending \\
H5 & Schema Alignment & Schemas support all requirement fields & Pending \\
H6 & Gap Analysis & No orphan requirements or missing mappings & Pending \\
\bottomrule
\end{tabularx}
\end{table}

\begin{hitlbox}{H1 -- Source Completeness}
Human reviewer must verify that all 6 source documents are correctly
catalogued with accurate file sizes, types, and metadata. Any missing
documents must be identified and the extraction process re-run.
\end{hitlbox}

% =============================================================================
\section{Scope and Coverage}
\label{sec:scope-coverage}

\subsection{In Scope}

\begin{itemize}[noitemsep]
  \item Trial Balance review process across 234 Legal Entities
  \item AI-augmented variance analysis and commentary generation
  \item Month-end close controls (M-3 to M+5 timeline)
  \item FORTM pack generation and submission
  \item BSS Risk Assessment integration points
\end{itemize}

\subsection{Out of Scope}

\begin{itemize}[noitemsep]
  \item Country IFRS reporting (separate process)
  \item SAP/PeopleSoft integration detail
  \item Real-time intraday controls
  \item Audit committee reporting (downstream consumer)
\end{itemize}

\subsection{Coverage Statistics}

\begin{table}[h]
\centering
\caption{Traceability Coverage Statistics}
\label{tab:coverage-statistics}
\begin{tabular}{@{}lrr@{}}
\toprule
\textbf{Metric} & \textbf{Count} & \textbf{Coverage} \\
\midrule
Source Documents & 6 & 100\% catalogued \\
Extracted Files & 7 & 100\% traced to source \\
Business Requirements & 27 & 100\% traced to source \\
Specifications & 21 & 100\% traced to requirements \\
Schemas & 6 & 100\% traced to specifications \\
Identified Gaps & 2 & Documented with owners \\
\bottomrule
\end{tabular}
\end{table}
\chapter{Source Document Inventory}
\label{ch:source-documents}

This chapter catalogues all source documents that form the foundation
of the Trial Balance review requirements. Each document is assigned a
unique source reference identifier for traceability purposes.

% =============================================================================
\section{Source Document Catalog}
\label{sec:source-catalog}

\begin{longtable}{@{}C{2cm}L{5cm}C{1.5cm}C{2cm}C{2.5cm}@{}}
\caption{Source Document Catalog}
\label{tab:source-catalog} \\
\toprule
\textbf{Source ID} & \textbf{Document Name} & \textbf{Type} & \textbf{Size} & \textbf{Extraction} \\
\midrule
\endfirsthead
\toprule
\textbf{Source ID} & \textbf{Document Name} & \textbf{Type} & \textbf{Size} & \textbf{Extraction} \\
\midrule
\endhead
\srcref{TB-BC-001} & Business case - Trial Balance (2).doc & .doc & 96 KB & Complete \\
\srcref{TB-BC-002} & Business case - BSS Risk Assessment (2).doc & .doc & 95 KB & Complete \\
\srcref{TB-DOI-001} & DOI for Trial Balance Process.docx & .docx & 6.3 MB & Complete \\
\srcref{TB-BPMN-001} & Operation of Month End Close Controls - Trial Balance Review - GLO (1).pdf & .pdf & 64 KB & Complete \\
\srcref{TB-WB-001} & HKG\_TB review Nov'25.xlsm & .xlsm & 73 MB & Schema Only \\
\srcref{TB-WB-002} & HKG\_PL review Nov'25.xlsm & .xlsm & 51 MB & Schema Only \\
\bottomrule
\end{longtable}

% =============================================================================
\section{Source Document Descriptions}
\label{sec:source-descriptions}

% -----------------------------------------------------------------------------
\subsection{\srcref{TB-BC-001}: Business Case - Trial Balance}
\label{sec:tb-bc-001}

\begin{table}[h]
\centering
\caption{\srcref{TB-BC-001} Metadata}
\begin{tabularx}{\textwidth}{@{}L{4cm}L{10cm}@{}}
\toprule
\textbf{Attribute} & \textbf{Value} \\
\midrule
Document Title & Business Case: Month End Controls -- Trial Balance \\
Filepath & \texttt{docs/tb/Business case - Trial Balance (2).doc} \\
Created By & 1498002 (S1 Badrinath) \\
Created Date & 2025-12-14 \\
Version & 1.0 \\
File Size & 97,792 bytes \\
Purpose & Define business justification for AI automation of TB review \\
\bottomrule
\end{tabularx}
\end{table}

\paragraph{Key Sections:}
\begin{itemize}[noitemsep]
  \item \textbf{Section 1: Executive Summary} --- TB review conducted M+1 to M+7 across 234 LEs; AI model for variance analysis and commentary generation
  \item \textbf{Section 2: Problem/Opportunity Statement} --- Current state analysis; 20-25 FTEs, 5 FTE savings target
  \item \textbf{Section 3: Project Description} --- Proposed AI solution using Studio AI (in-house)
  \item \textbf{Section 5: Financial Analysis} --- \textit{TBD --- Gap \gapref{G-001}}
  \item \textbf{Section 6: Risk Assessment} --- PSGL/S4 compatibility, SAP integration risks
  \item \textbf{Section 7: Implementation Approach} --- 15-week timeline
\end{itemize}

\paragraph{Content Summary:}
\begin{quote}
\textit{``Trial Balance review is conducted during M+1 to M+7 across all Legal entities,
Segment \& Group, to understand the variance \& balance movements between periods.
Deep dives to assess the events, transactions causing the variances and reasons
behind them. Analysed with detailed commentaries for these movements are prepared \&
submitted to Group for further review \& used for Country CFO packs/Legal entity reporting.''}
\end{quote}

% -----------------------------------------------------------------------------
\subsection{\srcref{TB-BC-002}: Business Case - BSS Risk Assessment}
\label{sec:tb-bc-002}

\begin{table}[h]
\centering
\caption{\srcref{TB-BC-002} Metadata}
\begin{tabularx}{\textwidth}{@{}L{4cm}L{10cm}@{}}
\toprule
\textbf{Attribute} & \textbf{Value} \\
\midrule
Document Title & Business Case: BSS Risk Assessment \\
Filepath & \texttt{docs/tb/Business case - BSS Risk Assessment (2).doc} \\
Created By & 1498002 (S1 Badrinath) \\
Created Date & 2025-12-15 \\
Version & 1.0 \\
File Size & 95,232 bytes \\
Purpose & Define AI augmentation for BSS risk assessment process \\
\bottomrule
\end{tabularx}
\end{table}

\paragraph{Key Content:}
\begin{itemize}[noitemsep]
  \item \textbf{Timing:} M+12 to M+15 across all Legal entities
  \item \textbf{Criticality:} CFO attestation and Risk/Audit committee reporting
  \item \textbf{Risk Categories:} Substantiated at Risk, Unsubstantiated Unowned, Unsubstantiated Owned
  \item \textbf{FTE Impact:} 10-12 FTEs current, 1 FTE savings target
\end{itemize}

% -----------------------------------------------------------------------------
\subsection{\srcref{TB-DOI-001}: DOI for Trial Balance Process}
\label{sec:tb-doi-001}

\begin{table}[h]
\centering
\caption{\srcref{TB-DOI-001} Metadata}
\begin{tabularx}{\textwidth}{@{}L{4cm}L{10cm}@{}}
\toprule
\textbf{Attribute} & \textbf{Value} \\
\midrule
Document Type & Detailed Operating Instructions \\
Filepath & \texttt{docs/tb/DOI for Trial Balance Process .docx} \\
File Size & 6,583,541 bytes \\
Purpose & Step-by-step operational guidance for TB review execution \\
Contains & Process flow diagrams (EMF/PNG images) \\
Image Count & 4 extracted images \\
\bottomrule
\end{tabularx}
\end{table}

\paragraph{Extracted Images:}
\begin{itemize}[noitemsep]
  \item \texttt{doi-images/doi-image-001.emf} --- Process flow diagram 1
  \item \texttt{doi-images/doi-image-002.emf} --- Process flow diagram 2
  \item \texttt{doi-images/doi-image-003.png} --- Process flow diagram 3
  \item \texttt{doi-images/doi-image-004.emf} --- Process flow diagram 4
\end{itemize}

% -----------------------------------------------------------------------------
\subsection{\srcref{TB-BPMN-001}: BPMN Process Map}
\label{sec:tb-bpmn-001}

\begin{table}[h]
\centering
\caption{\srcref{TB-BPMN-001} Metadata}
\begin{tabularx}{\textwidth}{@{}L{4cm}L{10cm}@{}}
\toprule
\textbf{Attribute} & \textbf{Value} \\
\midrule
Document Title & Operation of Month End Close Controls - Trial Balance Review - GLO \\
Filepath & \texttt{docs/tb/Operation of Month End Close Controls - Trial Balance Review - GLO (1).pdf} \\
Process Owner & Modi, Piyush \\
Last Changed & 2025-05-23T19:50:30Z \\
Scope & GLO (Global) \\
File Size & 65,367 bytes \\
\bottomrule
\end{tabularx}
\end{table}

\paragraph{BPMN Elements Extracted:}
\begin{itemize}[noitemsep]
  \item \textbf{States:} 27 (including 3 error-handling states)
  \item \textbf{Gateways:} 5 exclusive decision gateways
  \item \textbf{Timeline Phases:} 4 (Preparation, Extraction, Analysis, Review)
  \item \textbf{Participants:} 3 (GFS Analyst, Head of Africa GFS, Country Stakeholders)
\end{itemize}

% -----------------------------------------------------------------------------
\subsection{\srcref{TB-WB-001}: HKG TB Review Workbook}
\label{sec:tb-wb-001}

\begin{table}[h]
\centering
\caption{\srcref{TB-WB-001} Metadata}
\begin{tabularx}{\textwidth}{@{}L{4cm}L{10cm}@{}}
\toprule
\textbf{Attribute} & \textbf{Value} \\
\midrule
Document Title & HKG\_TB review Nov'25.xlsm \\
Filepath & \texttt{docs/tb/HKG\_TB review Nov'25.xlsm} \\
File Size & 76,989,527 bytes (73 MB) \\
Sheet Count & 13 \\
Extraction Status & Schema extraction only (data too large) \\
\bottomrule
\end{tabularx}
\end{table}

\paragraph{Key Sheets:}
\begin{itemize}[noitemsep]
  \item \texttt{IFRS Summary BS} --- IFRS balance sheet summary view
  \item \texttt{Count of Comments} --- Commentary tracking (7 fields)
  \item \texttt{BS Variance} --- Balance sheet variance analysis (5 fields)
  \item \texttt{RAW TB NOV'25} --- Raw GL balance query November 2025
  \item \texttt{RAW TB OCT'25} --- Raw GL balance query October 2025
  \item \texttt{GCOA} --- Global Chart of Accounts (5 fields)
  \item \texttt{Dept Mapping} --- Department-to-vertical mapping (12 fields)
  \item \texttt{Base file} --- Account-level base data (13 fields)
\end{itemize}

% -----------------------------------------------------------------------------
\subsection{\srcref{TB-WB-002}: HKG PL Review Workbook}
\label{sec:tb-wb-002}

\begin{table}[h]
\centering
\caption{\srcref{TB-WB-002} Metadata}
\begin{tabularx}{\textwidth}{@{}L{4cm}L{10cm}@{}}
\toprule
\textbf{Attribute} & \textbf{Value} \\
\midrule
Document Title & HKG\_PL review Nov'25.xlsm \\
Filepath & \texttt{docs/tb/HKG\_PL review Nov'25.xlsm} \\
File Size & 53,245,945 bytes (51 MB) \\
Sheet Count & 14 \\
Extraction Status & Schema extraction only (data too large) \\
\bottomrule
\end{tabularx}
\end{table}

% =============================================================================
\section{Source Document Verification}
\label{sec:source-verification}

\begin{hitlbox}{H1 -- Source Completeness Verification}
The human reviewer must verify the following for checkpoint H1:
\begin{enumerate}[noitemsep]
  \item All 6 source documents exist at the specified file paths
  \item File sizes match the catalogued values (within 1\% tolerance)
  \item Document metadata (author, date, version) is accurate
  \item No additional source documents have been omitted
  \item Extraction status correctly reflects processing state
\end{enumerate}
\end{hitlbox}

\paragraph{Verification Checklist:}

\begin{table}[h]
\centering
\caption{Source Document Verification Checklist}
\begin{tabular}{@{}lcccc@{}}
\toprule
\textbf{Source ID} & \textbf{Exists} & \textbf{Size OK} & \textbf{Metadata OK} & \textbf{Verified} \\
\midrule
\srcref{TB-BC-001} & $\square$ & $\square$ & $\square$ & $\square$ \\
\srcref{TB-BC-002} & $\square$ & $\square$ & $\square$ & $\square$ \\
\srcref{TB-DOI-001} & $\square$ & $\square$ & $\square$ & $\square$ \\
\srcref{TB-BPMN-001} & $\square$ & $\square$ & $\square$ & $\square$ \\
\srcref{TB-WB-001} & $\square$ & $\square$ & $\square$ & $\square$ \\
\srcref{TB-WB-002} & $\square$ & $\square$ & $\square$ & $\square$ \\
\bottomrule
\end{tabular}
\end{table}
\chapter{Machine-Readable Extraction Layer}
\label{ch:extraction-layer}

This chapter documents the extraction process that transforms source
documents into machine-readable formats, providing the foundation for
automated validation and requirements traceability.

% =============================================================================
\section{Extraction Process Overview}
\label{sec:extraction-overview}

Source documents were processed using automated extraction scripts to
produce machine-readable formats suitable for:

\begin{itemize}[noitemsep]
  \item Text analysis and requirement derivation
  \item Schema validation and data quality checks
  \item RAG (Retrieval-Augmented Generation) pipeline ingestion
  \item Version control and change tracking
\end{itemize}

\begin{table}[h]
\centering
\caption{Extraction Scripts and Outputs}
\label{tab:extraction-scripts}
\begin{tabularx}{\textwidth}{@{}L{3.5cm}L{4cm}L{6.5cm}@{}}
\toprule
\textbf{Script} & \textbf{Source Types} & \textbf{Output Format} \\
\midrule
\texttt{extract\_docs.py} & .doc, .docx & Markdown with YAML frontmatter \\
\texttt{extract\_pdf.py} & .pdf (BPMN) & Markdown + BPMN elements JSON \\
\texttt{extract\_xlsx.py} & .xlsm workbooks & YAML schemas + CSV samples \\
\texttt{build\_rag\_chunks.py} & All extracted text & JSONL chunks for RAG \\
\bottomrule
\end{tabularx}
\end{table}

% =============================================================================
\section{Extracted Document Mapping}
\label{sec:extracted-mapping}

\begin{longtable}{@{}C{2cm}L{5cm}L{6.5cm}@{}}
\caption{Source to Extracted Document Mapping}
\label{tab:extracted-mapping} \\
\toprule
\textbf{Source ID} & \textbf{Extracted File} & \textbf{Extraction Notes} \\
\midrule
\endfirsthead
\toprule
\textbf{Source ID} & \textbf{Extracted File} & \textbf{Extraction Notes} \\
\midrule
\endhead
\srcref{TB-BC-001} & \texttt{business-case-trial-balance.md} & Full text extraction; tables converted to markdown; 1,172 words \\
\srcref{TB-BC-002} & \texttt{business-case-bss-risk-assessment.md} & Full text extraction; 8,908 bytes \\
\srcref{TB-DOI-001} & \texttt{doi-trial-balance-process.md} & Text + 4 images extracted to \texttt{doi-images/} \\
\srcref{TB-BPMN-001} & \texttt{bpmn-tb-review-glo.md} + \texttt{bpmn-elements.json} & Process text + structured BPMN elements \\
\srcref{TB-WB-001} & \texttt{hkg-tb-schema.yaml} + \texttt{hkg-tb-sample.csv} & Schema for 13 sheets; 100-row sample \\
\srcref{TB-WB-002} & \texttt{hkg-pl-schema.yaml} + \texttt{hkg-pl-sample.csv} & Schema for 14 sheets; 100-row sample \\
\bottomrule
\end{longtable}

% =============================================================================
\section{Extraction File Inventory}
\label{sec:extraction-inventory}

\subsection{Extracted Text Files}

\begin{longtable}{@{}L{6cm}C{2cm}C{2cm}L{3.5cm}@{}}
\caption{Extracted Text File Inventory}
\label{tab:extracted-text-inventory} \\
\toprule
\textbf{File Path} & \textbf{Size} & \textbf{Format} & \textbf{Source Ref} \\
\midrule
\endfirsthead
\toprule
\textbf{File Path} & \textbf{Size} & \textbf{Format} & \textbf{Source Ref} \\
\midrule
\endhead
\texttt{extracted/business-case-trial-balance.md} & 9,824 B & Markdown & \srcref{TB-BC-001} \\
\texttt{extracted/business-case-bss-risk-assessment.md} & 8,908 B & Markdown & \srcref{TB-BC-002} \\
\texttt{extracted/doi-trial-balance-process.md} & 865 B & Markdown & \srcref{TB-DOI-001} \\
\texttt{extracted/bpmn-tb-review-glo.md} & 9,154 B & Markdown & \srcref{TB-BPMN-001} \\
\texttt{extracted/bpmn-elements.json} & 2,264 B & JSON & \srcref{TB-BPMN-001} \\
\bottomrule
\end{longtable}

\subsection{Extracted Image Files}

\begin{longtable}{@{}L{6cm}C{2cm}C{2cm}L{3.5cm}@{}}
\caption{Extracted Image File Inventory}
\label{tab:extracted-image-inventory} \\
\toprule
\textbf{File Path} & \textbf{Size} & \textbf{Format} & \textbf{Source Ref} \\
\midrule
\endfirsthead
\toprule
\textbf{File Path} & \textbf{Size} & \textbf{Format} & \textbf{Source Ref} \\
\midrule
\endhead
\texttt{extracted/doi-images/doi-image-001.emf} & 8,148 B & EMF & \srcref{TB-DOI-001} \\
\texttt{extracted/doi-images/doi-image-002.emf} & 8,272 B & EMF & \srcref{TB-DOI-001} \\
\texttt{extracted/doi-images/doi-image-003.png} & 43,650 B & PNG & \srcref{TB-DOI-001} \\
\texttt{extracted/doi-images/doi-image-004.emf} & 8,132 B & EMF & \srcref{TB-DOI-001} \\
\bottomrule
\end{longtable}

\subsection{Training Data Files}

\begin{longtable}{@{}L{6cm}C{2cm}C{2cm}L{3.5cm}@{}}
\caption{Training Data File Inventory}
\label{tab:training-data-inventory} \\
\toprule
\textbf{File Path} & \textbf{Size} & \textbf{Format} & \textbf{Source Ref} \\
\midrule
\endfirsthead
\toprule
\textbf{File Path} & \textbf{Size} & \textbf{Format} & \textbf{Source Ref} \\
\midrule
\endhead
\texttt{training-data/hkg-tb-schema.yaml} & 20,228 B & YAML & \srcref{TB-WB-001} \\
\texttt{training-data/hkg-tb-sample.csv} & 17,916 B & CSV & \srcref{TB-WB-001} \\
\texttt{training-data/hkg-pl-schema.yaml} & 16,982 B & YAML & \srcref{TB-WB-002} \\
\texttt{training-data/hkg-pl-sample.csv} & 19,110 B & CSV & \srcref{TB-WB-002} \\
\bottomrule
\end{longtable}

% =============================================================================
\section{Extraction Frontmatter Format}
\label{sec:extraction-frontmatter}

Each extracted markdown file includes YAML frontmatter for traceability
and metadata validation:

\begin{lstlisting}[language=yamlx,caption=Extraction Frontmatter Example]
---
source_file: 'Business case - Trial Balance (2).doc'
source_type: 'doc'
extracted_at: '2026-04-13T06:16:45.446975+00:00'
word_count: 1172
source_ref: 'TB-BC-001'
---
# Business Case: Trial Balance

## 1. Executive Summary
...
\end{lstlisting}

\paragraph{Frontmatter Fields:}

\begin{table}[h]
\centering
\caption{Extraction Frontmatter Schema}
\label{tab:frontmatter-schema}
\begin{tabularx}{\textwidth}{@{}L{3cm}C{2cm}L{8.5cm}@{}}
\toprule
\textbf{Field} & \textbf{Type} & \textbf{Description} \\
\midrule
\texttt{source\_file} & string & Original filename \\
\texttt{source\_type} & enum & File extension (.doc, .docx, .pdf, .xlsm) \\
\texttt{extracted\_at} & ISO 8601 & Timestamp of extraction \\
\texttt{word\_count} & integer & Word count of extracted text \\
\texttt{source\_ref} & string & Traceability reference ID \\
\bottomrule
\end{tabularx}
\end{table}

% =============================================================================
\section{BPMN Elements Extraction}
\label{sec:bpmn-extraction}

The BPMN process map (\srcref{TB-BPMN-001}) was parsed into structured
JSON to enable automated workflow validation:

\begin{lstlisting}[language=jsonx,caption=BPMN Elements Structure]
{
  "workflowId": "tb-review-glo",
  "version": "1.0.0",
  "states": [
    {
      "stateId": "START",
      "type": "start_event",
      "name": "Prepare Schedule Work Day",
      "transitions": [{"target": "ROLL_FORWARD"}]
    },
    {
      "stateId": "ROLL_FORWARD",
      "type": "task",
      "name": "Roll Forward Prior Month File",
      "timing": "M-3",
      "systems": ["MS Excel", "Shared Drive"],
      "transitions": [{"target": "DOWNLOAD_PSGL"}]
    }
  ],
  "gateways": [...],
  "dataObjects": [...],
  "participants": [...]
}
\end{lstlisting}

\paragraph{Extracted BPMN Statistics:}

\begin{table}[h]
\centering
\caption{BPMN Extraction Statistics}
\label{tab:bpmn-statistics}
\begin{tabular}{@{}lr@{}}
\toprule
\textbf{Element Type} & \textbf{Count} \\
\midrule
Start Events & 1 \\
Task States & 21 \\
Gateway States & 5 \\
Error-Handling States & 3 \\
End Events & 1 \\
\midrule
\textbf{Total States} & \textbf{27} \\
\midrule
Data Objects & 6 \\
Participants & 3 \\
\bottomrule
\end{tabular}
\end{table}

% =============================================================================
\section{Workbook Schema Extraction}
\label{sec:workbook-schema}

Excel workbooks were processed to extract structural schemas without
copying sensitive data:

\subsection{HKG TB Review Schema (\srcref{TB-WB-001})}

\begin{table}[h]
\centering
\caption{HKG TB Review Sheet Inventory}
\label{tab:hkg-tb-sheets}
\begin{tabularx}{\textwidth}{@{}L{4cm}C{2cm}L{7.5cm}@{}}
\toprule
\textbf{Sheet Name} & \textbf{Fields} & \textbf{Purpose} \\
\midrule
IFRS Summary BS & -- & IFRS balance sheet summary view \\
Count of Comments & 7 & Commentary tracking (reviewer, status, hours) \\
BS Variance & 5 & Balance sheet variance analysis with thresholds \\
PL Variance & 4 & P\&L variance analysis \\
GCOA & 5 & Global Chart of Accounts definitions \\
Dept Mapping & 12 & Department-to-vertical mapping \\
Base file & 13 & Account-level base data with mappings \\
\bottomrule
\end{tabularx}
\end{table}

% =============================================================================
\section{Extraction Quality Metrics}
\label{sec:extraction-quality}

\begin{table}[h]
\centering
\caption{Extraction Quality Metrics}
\label{tab:extraction-quality}
\begin{tabularx}{\textwidth}{@{}L{5cm}C{2.5cm}L{6cm}@{}}
\toprule
\textbf{Metric} & \textbf{Value} & \textbf{Human Validation Required} \\
\midrule
Source documents processed & 6 of 6 & Verify all documents included \\
Total extracted text size & 522 KB & Verify no truncation occurred \\
Tables extracted & 12 & Verify table structure preserved \\
Images extracted & 4 & Verify images are referenced correctly \\
BPMN elements parsed & 27 states & Verify state machine completeness \\
Workbook sheets processed & 27 & Verify key sheets included \\
\bottomrule
\end{tabularx}
\end{table}

\begin{hitlbox}{H2 -- Extraction Accuracy Verification}
The human reviewer must verify the following for checkpoint H2:
\begin{enumerate}[noitemsep]
  \item Sample extracted text against original documents (minimum 3 paragraphs per document)
  \item Table formatting is preserved in markdown conversion
  \item BPMN state names and transitions match the source PDF
  \item Workbook schema field names match column headers
  \item No content truncation or corruption occurred
\end{enumerate}
\end{hitlbox}
% =============================================================================
% Chapter 4: Business Requirements (Enhanced with Source Quotations)
% =============================================================================
\chapter{Business Requirements Catalogue}
\label{ch:business-requirements}

This chapter presents all business requirements derived from the source documents,
with \textbf{exact quotations} and \textbf{page/paragraph references} enabling 
human auditors to verify each requirement against its source.

% -----------------------------------------------------------------------------
\section{How to Read This Chapter}
% -----------------------------------------------------------------------------

Each requirement entry contains:

\begin{tcolorbox}[colback=blue!5,colframe=blue!75!black,title=Requirement Entry Format]
\begin{enumerate}[leftmargin=*]
    \item \textbf{Requirement ID:} Unique identifier (e.g., TB-REQ-ES01)
    \item \textbf{MoSCoW Priority:} MUST / SHOULD / COULD / WON'T
    \item \textbf{Source Quotation:} Exact text from source document (in italics)
    \item \textbf{Source Location:} Document ID, section, paragraph, page number
    \item \textbf{Derived Requirement:} Formal requirement statement
    \item \textbf{Acceptance Criteria:} Measurable conditions for verification
    \item \textbf{Rationale:} Why this requirement was derived from the quoted text
\end{enumerate}
\end{tcolorbox}

% =============================================================================
\section{Executive Summary Requirements}
\label{sec:req-executive}
% =============================================================================

% -----------------------------------------------------------------------------
\subsection{TB-REQ-ES01: AI-Driven Variance Analysis}
\label{req:tb-req-es01}
% -----------------------------------------------------------------------------

\begin{tcolorbox}[colback=sourcegreen!5,colframe=sourcegreen!70,title=Source Quotation]
\textit{``An AI model produces business insights and written commentaries 
(Finance reporting caption Level \& Segment Level with Legal entity view) 
based on the analysis performed on general ledger transaction data that 
\textbf{detects anomalies, trends, unusual movements, new accounts, 
spikes/drops vs prior periods, and material balances}.''} 
\end{tcolorbox}

\begin{tabularx}{\textwidth}{@{}lX@{}}
\toprule
\textbf{Attribute} & \textbf{Value} \\
\midrule
\textbf{Requirement ID} & TB-REQ-ES01 \\
\textbf{MoSCoW Priority} & \textcolor{red}{\textbf{MUST}} \\
\textbf{Source Document} & Business case - Trial Balance (2).doc \\
\textbf{Source ID} & TB-BC-001 \\
\textbf{Source Location} & Section 1: Executive Summary, Paragraph 3, Page 1 (Lines 62--68 in extracted MD) \\
\textbf{Created By} & Employee ID 1498002 (S1 Badrinath) \\
\textbf{Document Date} & 14-Dec-2025 \\
\bottomrule
\end{tabularx}

\vspace{0.5em}
\textbf{Derived Requirement Statement:}

The system SHALL implement an AI model capable of performing the following 
analyses on general ledger (GL) transaction data:

\begin{enumerate}[label=(\alph*)]
    \item \textbf{Anomaly Detection:} Identify transactions or balances that 
          deviate significantly from expected patterns
    \item \textbf{Trend Analysis:} Detect multi-period trends in account balances
    \item \textbf{Unusual Movement Detection:} Flag balance changes that exceed 
          normal operational variance
    \item \textbf{New Account Detection:} Identify accounts appearing for the 
          first time in the current period
    \item \textbf{Spike/Drop Detection:} Identify sudden increases or decreases 
          compared to prior periods
    \item \textbf{Material Balance Identification:} Flag balances exceeding 
          defined materiality thresholds
\end{enumerate}

\vspace{0.5em}
\textbf{Acceptance Criteria:}

\begin{tcolorbox}[colback=sapamber!5,colframe=sapamber!70,title=Measurable Acceptance Criteria]
\begin{enumerate}[label=AC-ES01-\arabic*:,leftmargin=*]
    \item The AI model SHALL process GL data for all 234 Legal Entities 
          (per TB-BC-001, Section 2)
    \item Anomaly detection SHALL flag items with Z-score $> 3.0$ or configurable threshold
    \item Spike/drop detection SHALL identify variances $> 25\%$ period-over-period 
          (or entity-configurable)
    \item New account detection SHALL compare current period CoA against prior 
          12-month baseline
    \item Material balance threshold SHALL default to entity-specific materiality 
          (configurable per LE)
    \item Processing SHALL complete within M+1 to M+5 window (per BPMN: TB-BPMN-001)
\end{enumerate}
\end{tcolorbox}

\textbf{Derivation Rationale:} The source text explicitly lists six AI capabilities 
(anomalies, trends, unusual movements, new accounts, spikes/drops, material balances). 
Each is converted to a testable sub-requirement with quantifiable criteria.

% -----------------------------------------------------------------------------
\subsection{TB-REQ-ES02: Commentary Generation}
\label{req:tb-req-es02}
% -----------------------------------------------------------------------------

\begin{tcolorbox}[colback=sourcegreen!5,colframe=sourcegreen!70,title=Source Quotation]
\textit{``GCFO teams get a \textbf{narrative summary} instead of just numbers, 
accelerating review and improving decision-making.''}
\end{tcolorbox}

\begin{tabularx}{\textwidth}{@{}lX@{}}
\toprule
\textbf{Attribute} & \textbf{Value} \\
\midrule
\textbf{Requirement ID} & TB-REQ-ES02 \\
\textbf{MoSCoW Priority} & \textcolor{red}{\textbf{MUST}} \\
\textbf{Source Document} & Business case - Trial Balance (2).doc \\
\textbf{Source ID} & TB-BC-001 \\
\textbf{Source Location} & Section 1: Executive Summary, Paragraph 3, Page 1 (Lines 66--68) \\
\bottomrule
\end{tabularx}

\vspace{0.5em}
\textbf{Derived Requirement Statement:}

The system SHALL generate natural language commentary summaries for:
\begin{itemize}
    \item Finance reporting caption level
    \item Segment level
    \item Legal entity level
\end{itemize}

\textbf{Acceptance Criteria:}

\begin{tcolorbox}[colback=sapamber!5,colframe=sapamber!70,title=Measurable Acceptance Criteria]
\begin{enumerate}[label=AC-ES02-\arabic*:,leftmargin=*]
    \item Commentary SHALL be generated in grammatically correct English prose
    \item Commentary SHALL reference specific variance amounts (e.g., ``Revenue increased 
          by \$2.3M (15\%) due to...'')
    \item Commentary SHALL cite contributing factors identified by AI analysis
    \item Commentary confidence score SHALL be displayed (0--100\%)
    \item Human reviewer SHALL be able to edit/approve/reject each commentary
    \item Commentary generation SHALL complete within 60 seconds per Legal Entity
\end{enumerate}
\end{tcolorbox}

% -----------------------------------------------------------------------------
\subsection{TB-REQ-ES03: FTE Savings Target}
\label{req:tb-req-es03}
% -----------------------------------------------------------------------------

\begin{tcolorbox}[colback=sourcegreen!5,colframe=sourcegreen!70,title=Source Quotation]
\textit{``At least \textbf{20 to 25 FTEs} spent on an average \textbf{4 hrs a day} 
during M-2 to M+5, which is equivalent to \textbf{$\sim$5 FTE saves}, with 
strengthening controls by eliminating manual process and reducing the overall 
book close timelines enabling Fast close.''}
\end{tcolorbox}

\begin{tabularx}{\textwidth}{@{}lX@{}}
\toprule
\textbf{Attribute} & \textbf{Value} \\
\midrule
\textbf{Requirement ID} & TB-REQ-ES03 \\
\textbf{MoSCoW Priority} & \textcolor{red}{\textbf{MUST}} \\
\textbf{Source Document} & Business case - Trial Balance (2).doc \\
\textbf{Source ID} & TB-BC-001 \\
\textbf{Source Location} & Section 1: Executive Summary, Paragraph 4, Pages 1--2 (Lines 72--76) \\
\bottomrule
\end{tabularx}

\vspace{0.5em}
\textbf{Derived Requirement Statement:}

The system SHALL reduce manual effort in Trial Balance review from current 
baseline of 20--25 FTEs $\times$ 4 hours/day to achieve a minimum 5 FTE 
equivalent savings.

\textbf{Acceptance Criteria:}

\begin{tcolorbox}[colback=sapamber!5,colframe=sapamber!70,title=Measurable Acceptance Criteria]
\begin{enumerate}[label=AC-ES03-\arabic*:,leftmargin=*]
    \item Baseline measurement: Current effort = 20--25 FTEs $\times$ 4 hrs/day 
          $\times$ 8 days (M-2 to M+5) = 640--800 person-hours per close cycle
    \item Target measurement: Post-implementation effort $\leq$ 480 person-hours 
          (15--20 FTEs $\times$ 4 hrs $\times$ 8 days)
    \item Savings calculation: $\geq$ 160 person-hours saved per cycle (5 FTE equivalent)
    \item Measurement method: Time tracking via ServiceNow tickets or equivalent
    \item Verification: Quarterly review of actual vs. target savings
\end{enumerate}
\end{tcolorbox}

\textbf{Financial Context:}

\begin{tabular}{lrrr}
\toprule
\textbf{Metric} & \textbf{Current State} & \textbf{Target State} & \textbf{Savings} \\
\midrule
FTEs involved & 20--25 & 15--20 & 5 \\
Hours per day & 4 & 4 & -- \\
Days per cycle (M-2 to M+5) & 8 & 8 & -- \\
Person-hours per cycle & 640--800 & 480--640 & 160 (min) \\
\bottomrule
\end{tabular}

% =============================================================================
\section{Scope Requirements}
\label{sec:req-scope}
% =============================================================================

% -----------------------------------------------------------------------------
\subsection{TB-REQ-SC01: Legal Entity Coverage}
\label{req:tb-req-sc01}
% -----------------------------------------------------------------------------

\begin{tcolorbox}[colback=sourcegreen!5,colframe=sourcegreen!70,title=Source Quotation]
\textit{``Downloading Trail balance from PSGL/S4 \& preparing the detailed 
analysis across all \textbf{234 LEs} are highly time consuming...''}
\end{tcolorbox}

\begin{tabularx}{\textwidth}{@{}lX@{}}
\toprule
\textbf{Requirement ID} & TB-REQ-SC01 \\
\textbf{MoSCoW Priority} & \textcolor{red}{\textbf{MUST}} \\
\textbf{Source Location} & TB-BC-001, Section 1, Paragraph 2, Page 1 (Lines 55--59) \\
\bottomrule
\end{tabularx}

\vspace{0.5em}
\textbf{Derived Requirement Statement:}

The system SHALL support Trial Balance review for exactly 234 Legal Entities.

\textbf{Acceptance Criteria:}
\begin{enumerate}[label=AC-SC01-\arabic*:,leftmargin=*]
    \item System configuration SHALL include all 234 Legal Entity codes
    \item System SHALL process TB data for each LE independently
    \item Dashboard SHALL display completion status for all 234 LEs
    \item No Legal Entity SHALL be excluded from automated processing
\end{enumerate}

% -----------------------------------------------------------------------------
\subsection{TB-REQ-SC02: Reporting Dimensions}
\label{req:tb-req-sc02}
% -----------------------------------------------------------------------------

\begin{tcolorbox}[colback=sourcegreen!5,colframe=sourcegreen!70,title=Source Quotation]
\textit{``...current view is at \textbf{Legal entity level} and need to be in 
\textbf{Finance caption level \& segment level}.''}
\end{tcolorbox}

\begin{tabularx}{\textwidth}{@{}lX@{}}
\toprule
\textbf{Requirement ID} & TB-REQ-SC02 \\
\textbf{MoSCoW Priority} & \textcolor{red}{\textbf{MUST}} \\
\textbf{Source Location} & TB-BC-001, Section 1, Paragraph 2, Pages 1 (Lines 55--59) \\
\bottomrule
\end{tabularx}

\vspace{0.5em}
\textbf{Derived Requirement Statement:}

The system SHALL provide variance analysis and commentary at three reporting dimensions:

\begin{enumerate}[label=(\alph*)]
    \item \textbf{Legal Entity Level:} Individual LE view (current capability, to be retained)
    \item \textbf{Finance Caption Level:} Aggregation by standardized finance captions 
          (e.g., Revenue, COGS, Operating Expenses)
    \item \textbf{Segment Level:} Aggregation by business segment as defined in CoA
\end{enumerate}

\textbf{Acceptance Criteria:}
\begin{enumerate}[label=AC-SC02-\arabic*:,leftmargin=*]
    \item User SHALL be able to toggle between all three reporting dimensions
    \item Drill-down SHALL be available from Segment $\rightarrow$ Finance Caption $\rightarrow$ LE
    \item Commentary SHALL be generated at each dimension level
    \item Export to Excel SHALL preserve dimension hierarchy
\end{enumerate}

% -----------------------------------------------------------------------------
\subsection{TB-REQ-SC03: Out of Scope Items}
\label{req:tb-req-sc03}
% -----------------------------------------------------------------------------

\begin{tcolorbox}[colback=sourcegreen!5,colframe=sourcegreen!70,title=Source Quotation]
\textit{``Out of scope – \textbf{Country IFRS reporting, BSS Risk assessment}''}
\end{tcolorbox}

\begin{tabularx}{\textwidth}{@{}lX@{}}
\toprule
\textbf{Requirement ID} & TB-REQ-SC03 \\
\textbf{MoSCoW Priority} & \textcolor{gray}{\textbf{WON'T}} (this phase) \\
\textbf{Source Location} & TB-BC-001, Section 3: Project Description, Scope, Page 2 (Lines 152--154) \\
\bottomrule
\end{tabularx}

\vspace{0.5em}
\textbf{Explicit Exclusions:}

The following items are \textbf{explicitly out of scope} per source documentation:

\begin{enumerate}
    \item \textbf{Country IFRS Reporting:} Local GAAP/IFRS adjustments not included
    \item \textbf{BSS Risk Assessment:} Balance Sheet Substantiation is a separate process
\end{enumerate}

\textbf{Note:} These items are documented as future phases (see Gap G-003 in Chapter~\ref{ch:gap-analysis}).

% =============================================================================
\section{Process Step Requirements (from BPMN)}
\label{sec:req-process}
% =============================================================================

The following requirements are derived from the BPMN process diagram 
\textbf{TB-BPMN-001}: ``Operation of Month End Close Controls - Trial Balance 
Review - GLO''.

\textbf{Source:} \texttt{Operation of Month End Close Controls - Trial Balance Review - GLO (1).pdf}\\
\textbf{Process Owner:} Modi, Piyush\\
\textbf{Last Changed:} May 23, 2025, 7:50:30 PM\\
\textbf{Author ID:} 8203488

% -----------------------------------------------------------------------------
\subsection{TB-REQ-PS01: Roll Forward Prior Month File}
\label{req:tb-req-ps01}
% -----------------------------------------------------------------------------

\begin{tcolorbox}[colback=sourcegreen!5,colframe=sourcegreen!70,title=BPMN Task Extract]
\textit{``Roll Forward Prior Month File to Reflect Current...''}\\
\textbf{Timing:} M-3\\
\textbf{System:} MS-E-xcel-1, Shared Drive
\end{tcolorbox}

\begin{tabularx}{\textwidth}{@{}lX@{}}
\toprule
\textbf{Requirement ID} & TB-REQ-PS01 \\
\textbf{MoSCoW Priority} & \textcolor{red}{\textbf{MUST}} \\
\textbf{Source Location} & TB-BPMN-001, Task Row 1, Timing M-3 \\
\textbf{EUC Reference} & N/A (no EUC control on this step) \\
\bottomrule
\end{tabularx}

\textbf{Derived Requirement:} The system SHALL roll forward the prior month's 
Trial Balance workbook structure to the current period while preserving:
\begin{itemize}
    \item Prior period comparative balances
    \item Account structure and mappings
    \item Historical commentary for reference
\end{itemize}

\textbf{Acceptance Criteria:}
\begin{enumerate}[label=AC-PS01-\arabic*:]
    \item Roll-forward SHALL complete by M-3 (3 days before month end)
    \item Prior month comparatives SHALL be locked from editing
    \item Current period columns SHALL be initialized to zero
\end{enumerate}

% -----------------------------------------------------------------------------
\subsection{TB-REQ-PS02: Download PSGL Master Data}
\label{req:tb-req-ps02}
% -----------------------------------------------------------------------------

\begin{tcolorbox}[colback=sourcegreen!5,colframe=sourcegreen!70,title=BPMN Task Extract]
\textit{``Download PSGL [...] Update 'Accounts Master' Sheet''}\\
\textbf{Timing:} M-3\\
\textbf{System:} PSGL, MS-E-xcel-1\\
\textbf{EUC Control:} EX/200/1353107/0001/000
\end{tcolorbox}

\begin{tabularx}{\textwidth}{@{}lX@{}}
\toprule
\textbf{Requirement ID} & TB-REQ-PS02 \\
\textbf{MoSCoW Priority} & \textcolor{red}{\textbf{MUST}} \\
\textbf{Source Location} & TB-BPMN-001, Task Row 2, Timing M-3 \\
\textbf{EUC Reference} & EX/200/1353107/0001/000 \\
\bottomrule
\end{tabularx}

\textbf{Derived Requirement:} The system SHALL extract Chart of Accounts (CoA) 
master data from PSGL and update the Accounts Master reference.

\textbf{Acceptance Criteria:}
\begin{enumerate}[label=AC-PS02-\arabic*:]
    \item CoA extract SHALL include all active GL accounts
    \item Extract SHALL match EUC control EX/200/1353107/0001/000 specifications
    \item New accounts SHALL be flagged for review (feeds TB-REQ-AI03)
\end{enumerate}

% -----------------------------------------------------------------------------
\subsection{TB-REQ-PS03: Extract Trial Balance}
\label{req:tb-req-ps03}
% -----------------------------------------------------------------------------

\begin{tcolorbox}[colback=sourcegreen!5,colframe=sourcegreen!70,title=BPMN Task Extract]
\textit{``Extract Trial Balance [...] File in Current Month''}\\
\textbf{Timing:} M-3 to M+5\\
\textbf{System:} PSGL, MS-E-xcel-1, Shared Drive
\end{tcolorbox}

\begin{tabularx}{\textwidth}{@{}lX@{}}
\toprule
\textbf{Requirement ID} & TB-REQ-PS03 \\
\textbf{MoSCoW Priority} & \textcolor{red}{\textbf{MUST}} \\
\textbf{Source Location} & TB-BPMN-001, Task Row 3 \\
\bottomrule
\end{tabularx}

\textbf{Derived Requirement:} The system SHALL extract Trial Balance data from 
PSGL/S4 for all 234 Legal Entities.

\textbf{Acceptance Criteria:}
\begin{enumerate}[label=AC-PS03-\arabic*:]
    \item Extract SHALL include: Account Number, Account Name, Opening Balance, 
          Movements, Closing Balance
    \item Extract SHALL be available in both local currency (LC) and reporting 
          currency (RC)
    \item Data SHALL be timestamped with extraction date/time
    \item Hash/checksum SHALL be generated for data integrity verification
\end{enumerate}

% -----------------------------------------------------------------------------
\subsection{TB-REQ-PS04: Calculate Variance}
\label{req:tb-req-ps04}
% -----------------------------------------------------------------------------

\begin{tcolorbox}[colback=sourcegreen!5,colframe=sourcegreen!70,title=BPMN Task Extract]
\textit{``Calculate Variance Between Prior [...] and Current Month''}\\
\textbf{Timing:} M-2 to M+5\\
\textbf{System:} MS-E-xcel-1, Shared Drive\\
\textbf{EUC Control:} EX/200/1579702/0001/000
\end{tcolorbox}

\begin{tabularx}{\textwidth}{@{}lX@{}}
\toprule
\textbf{Requirement ID} & TB-REQ-PS04 \\
\textbf{MoSCoW Priority} & \textcolor{red}{\textbf{MUST}} \\
\textbf{Source Location} & TB-BPMN-001, Task Row 4 \\
\textbf{EUC Reference} & EX/200/1579702/0001/000 \\
\bottomrule
\end{tabularx}

\textbf{Derived Requirement:} The system SHALL calculate variance between:
\begin{itemize}
    \item Current month balance vs. prior month balance
    \item Current month balance vs. same month prior year (YoY)
    \item Both absolute and percentage variance
\end{itemize}

\textbf{Acceptance Criteria:}
\begin{enumerate}[label=AC-PS04-\arabic*:]
    \item Variance calculation: $V = B_{current} - B_{prior}$
    \item Percentage variance: $V\% = \frac{B_{current} - B_{prior}}{|B_{prior}|} \times 100$
    \item Division by zero handling: If $B_{prior} = 0$, display ``N/A (New)''
    \item Calculation SHALL match EUC control specifications
\end{enumerate}

% -----------------------------------------------------------------------------
\subsection{TB-REQ-PS05: Filter by Materiality}
\label{req:tb-req-ps05}
% -----------------------------------------------------------------------------

\begin{tcolorbox}[colback=sourcegreen!5,colframe=sourcegreen!70,title=BPMN Task Extract]
\textit{``Filter Variances [...] Set Criteria''}\\
\textbf{Timing:} M-2 to M+5
\end{tcolorbox}

\begin{tabularx}{\textwidth}{@{}lX@{}}
\toprule
\textbf{Requirement ID} & TB-REQ-PS05 \\
\textbf{MoSCoW Priority} & \textcolor{red}{\textbf{MUST}} \\
\textbf{Source Location} & TB-BPMN-001, Task Row 5 \\
\bottomrule
\end{tabularx}

\textbf{Derived Requirement:} The system SHALL filter variances based on 
configurable materiality thresholds.

\textbf{Acceptance Criteria:}
\begin{enumerate}[label=AC-PS05-\arabic*:]
    \item Materiality threshold SHALL be configurable per Legal Entity
    \item Default threshold: Absolute $> \$50,000$ OR Percentage $> 10\%$
    \item Filter SHALL support AND/OR logic combinations
    \item Filter SHALL support account-type-specific thresholds
\end{enumerate}

% =============================================================================
\section{Requirements Summary Matrix}
\label{sec:req-summary}
% =============================================================================

\begin{longtable}{p{2cm} p{1.2cm} p{4cm} p{3cm} p{2.5cm}}
\caption{Complete Requirements Traceability Summary} \\
\toprule
\textbf{Req ID} & \textbf{MoSCoW} & \textbf{Title} & \textbf{Source Ref} & \textbf{Acceptance Criteria} \\
\midrule
\endfirsthead
\toprule
\textbf{Req ID} & \textbf{MoSCoW} & \textbf{Title} & \textbf{Source Ref} & \textbf{Acceptance Criteria} \\
\midrule
\endhead
\rowcolor{sapblue!5}
TB-REQ-ES01 & MUST & AI Variance Analysis & TB-BC-001, S1P3 & AC-ES01-1..6 \\
TB-REQ-ES02 & MUST & Commentary Generation & TB-BC-001, S1P3 & AC-ES02-1..6 \\
\rowcolor{sapblue!5}
TB-REQ-ES03 & MUST & 5 FTE Savings & TB-BC-001, S1P4 & AC-ES03-1..5 \\
TB-REQ-SC01 & MUST & 234 LE Coverage & TB-BC-001, S1P2 & AC-SC01-1..4 \\
\rowcolor{sapblue!5}
TB-REQ-SC02 & MUST & 3-Level Reporting & TB-BC-001, S1P2 & AC-SC02-1..4 \\
TB-REQ-SC03 & WON'T & Exclusions Documented & TB-BC-001, S3 & N/A \\
\rowcolor{sapblue!5}
TB-REQ-PS01 & MUST & Roll Forward & TB-BPMN-001 & AC-PS01-1..3 \\
TB-REQ-PS02 & MUST & PSGL Master Download & TB-BPMN-001 & AC-PS02-1..3 \\
\rowcolor{sapblue!5}
TB-REQ-PS03 & MUST & TB Extract & TB-BPMN-001 & AC-PS03-1..4 \\
TB-REQ-PS04 & MUST & Variance Calculation & TB-BPMN-001 & AC-PS04-1..4 \\
\rowcolor{sapblue!5}
TB-REQ-PS05 & MUST & Materiality Filter & TB-BPMN-001 & AC-PS05-1..4 \\
\bottomrule
\end{longtable}

\vspace{1em}
\textbf{MoSCoW Distribution:}
\begin{itemize}
    \item \textcolor{red}{\textbf{MUST:}} 10 requirements (91\%)
    \item \textcolor{gray}{\textbf{WON'T:}} 1 requirement (9\%) -- explicitly excluded from scope
\end{itemize}

% =============================================================================
\section{Human Auditor Verification Checklist}
\label{sec:req-audit-checklist}
% =============================================================================

\begin{tcolorbox}[colback=yellow!10,colframe=yellow!50!black,title=\faIcon{clipboard-check} Auditor Verification Checklist]
For each requirement above, the human auditor should verify:

\begin{enumerate}[label=\faIcon{square}]
    \item \textbf{Source Quotation Accuracy:} Does the quoted text match the 
          original document exactly?
    \item \textbf{Location Correctness:} Is the section/paragraph/page reference correct?
    \item \textbf{Derivation Logic:} Is the requirement a reasonable derivation 
          from the quoted text?
    \item \textbf{Acceptance Criteria Measurability:} Can each criterion be 
          objectively verified?
    \item \textbf{MoSCoW Appropriateness:} Does the priority reflect business importance?
    \item \textbf{Completeness:} Are there any requirements in the source document 
          not captured here?
\end{enumerate}

\vspace{0.5em}
\textbf{Signoff:}
\begin{itemize}
    \item[] Name: \rule{6cm}{0.4pt}
    \item[] Date: \rule{3cm}{0.4pt}
    \item[] Signature: \rule{5cm}{0.4pt}
\end{itemize}
\end{tcolorbox}
\chapter{Specification Mapping}
\label{ch:specification-mapping}

This chapter maps each specification element to its source requirement(s),
establishing the requirement-to-specification traceability layer of the
HITL framework.

% =============================================================================
\section{Process Step Specifications}
\label{sec:process-step-specs}

\begin{longtable}{@{}C{2cm}L{3.5cm}C{3cm}L{5cm}@{}}
\caption{Process Steps to Requirements Mapping}
\label{tab:process-steps-mapping} \\
\toprule
\textbf{Spec ID} & \textbf{Specification} & \textbf{Req ID(s)} & \textbf{Implementation Notes} \\
\midrule
\endfirsthead
\toprule
\textbf{Spec ID} & \textbf{Specification} & \textbf{Req ID(s)} & \textbf{Implementation Notes} \\
\midrule
\endhead
\specref{TB-STEP-001} & Roll Forward Prior Month File & \reqref{TB-REQ-PS01} & Auto-generate template; pre-populate parameters \\
\specref{TB-STEP-002} & Download PSGL \& Update Master & \reqref{TB-REQ-PS02} & PSGL/S4 API integration; auto-detect changes \\
\specref{TB-STEP-003} & Extract Trial Balance & \reqref{TB-REQ-PS03}, \reqref{TB-REQ-SC01} & Parallel extraction for 234 LEs \\
\specref{TB-STEP-004} & Calculate \& Filter Variances & \reqref{TB-REQ-PS04}, \reqref{TB-REQ-PS05} & ML-adaptive materiality thresholds \\
\specref{TB-STEP-005} & Seek \& Update Commentaries & \reqref{TB-REQ-PS06}, \reqref{TB-REQ-AI06} & LLM-generated draft commentaries \\
\specref{TB-STEP-006} & Journal Posting Investigation & \reqref{TB-REQ-PS07} & Auto-detect unposted journals \\
\specref{TB-STEP-007} & Track Entries \& Risk & \reqref{TB-REQ-PS07} & Predictive risk scoring \\
\specref{TB-STEP-008} & Prepare Summary \& Review & \reqref{TB-REQ-PS08} & Auto-generate FORTM pack section \\
\bottomrule
\end{longtable}

% =============================================================================
\section{Decision Point Specifications}
\label{sec:decision-point-specs}

\begin{longtable}{@{}C{1.5cm}L{3cm}C{3cm}L{6cm}@{}}
\caption{Decision Points to Requirements Mapping}
\label{tab:decision-points-mapping} \\
\toprule
\textbf{DP ID} & \textbf{Decision Point} & \textbf{Req ID(s)} & \textbf{Specification Detail} \\
\midrule
\endfirsthead
\toprule
\textbf{DP ID} & \textbf{Decision Point} & \textbf{Req ID(s)} & \textbf{Specification Detail} \\
\midrule
\endhead
\specref{DP-001} & Materiality Check & \reqref{TB-REQ-PS05}, \reqref{TB-REQ-AI05} & Threshold: BS \$100M, PL \$3M; ML-adaptive per category \\
\specref{DP-002} & Journal Status & \reqref{TB-REQ-PS07} & States: posted / posting\_required / not\_posted \\
\specref{DP-003} & Investigation Required & \reqref{TB-REQ-PS07} & Predictive model based on entry history \\
\specref{DP-004} & Entry Confirmation & \reqref{TB-REQ-PS07} & M+4 deadline tracking; auto-escalation \\
\specref{DP-005} & Unexplained Variance & \reqref{TB-REQ-AI06}, \reqref{BSS-REQ-01} & Categorize: substantiated\_at\_risk, unsubstantiated\_unowned, unsubstantiated\_owned \\
\bottomrule
\end{longtable}

\subsection{Decision Point Detail: \specref{DP-001}}

\begin{table}[h]
\centering
\caption{\specref{DP-001} Materiality Check Specification}
\begin{tabularx}{\textwidth}{@{}L{4cm}L{10cm}@{}}
\toprule
\textbf{Attribute} & \textbf{Value} \\
\midrule
Specification ID & DP-001 \\
Name & Materiality Check \\
Gateway Type & Exclusive \\
BPMN State & \stateid{CHECK\_MATERIAL\_VARIANCES} \\
Input Metric & \fieldref{variance\_amount} \\
\midrule
\multicolumn{2}{l}{\textbf{Threshold Configuration}} \\
\midrule
Balance Sheet (BS) & \$100,000,000 \\
Profit \& Loss (PL) & \$3,000,000 \\
Threshold Source & Static (v1.0), ML-adaptive (v2.0) \\
\midrule
\multicolumn{2}{l}{\textbf{Transitions}} \\
\midrule
Material = Yes & $\rightarrow$ \stateid{SEEK\_COMMENTARIES} \\
Material = No & $\rightarrow$ \stateid{CHECK\_UNEXPLAINED} \\
\bottomrule
\end{tabularx}
\end{table}

% =============================================================================
\section{AI Capability Specifications}
\label{sec:ai-capability-specs}

\begin{longtable}{@{}C{2cm}L{3.5cm}C{2.5cm}L{5.5cm}@{}}
\caption{AI Capabilities to Requirements Mapping}
\label{tab:ai-capabilities-mapping} \\
\toprule
\textbf{Spec ID} & \textbf{AI Capability} & \textbf{Req ID(s)} & \textbf{Implementation Detail} \\
\midrule
\endfirsthead
\toprule
\textbf{Spec ID} & \textbf{AI Capability} & \textbf{Req ID(s)} & \textbf{Implementation Detail} \\
\midrule
\endhead
\specref{AI-ANOM} & Anomaly Detection & \reqref{TB-REQ-AI01} & 3$\sigma$ z-score with MAD fallback for non-normal distributions \\
\specref{AI-TREND} & Trend Analysis & \reqref{TB-REQ-AI02} & 12-month rolling window; percentile ranking \\
\specref{AI-NEW} & New Account Detection & \reqref{TB-REQ-AI03} & Flag accounts not in prior period master \\
\specref{AI-SPIKE} & Spike/Drop Detection & \reqref{TB-REQ-AI04} & Period spike check; sign flip detection \\
\specref{AI-MAT} & Material Balance ID & \reqref{TB-REQ-AI05} & Category-specific thresholds with multipliers \\
\specref{AI-LLM} & LLM Commentary & \reqref{TB-REQ-AI06} & RAG-augmented generation; $>$85\% acceptance rate target \\
\bottomrule
\end{longtable}

\subsection{AI Capability Detail: \specref{AI-ANOM}}

\begin{table}[h]
\centering
\caption{\specref{AI-ANOM} Anomaly Detection Specification}
\begin{tabularx}{\textwidth}{@{}L{4cm}L{10cm}@{}}
\toprule
\textbf{Attribute} & \textbf{Value} \\
\midrule
Specification ID & AI-ANOM \\
Name & Anomaly Detection \\
Requirement & \reqref{TB-REQ-AI01} \\
\midrule
\multicolumn{2}{l}{\textbf{Algorithm Configuration}} \\
\midrule
Primary Method & Z-score (3$\sigma$ threshold) \\
Fallback Method & Median Absolute Deviation (MAD) \\
Fallback Trigger & Non-normal distribution detected (Shapiro-Wilk $p < 0.05$) \\
\midrule
\multicolumn{2}{l}{\textbf{Output Flags}} \\
\midrule
Period Spike & \enumval{PERIOD\_SPIKE} \\
Sign Flip & \enumval{SIGN\_FLIP} \\
New Account & \enumval{NEW\_ACCOUNT} \\
Material Balance & \enumval{MATERIAL\_BALANCE} \\
\bottomrule
\end{tabularx}
\end{table}

\subsection{AI Capability Detail: \specref{AI-LLM}}

\begin{table}[h]
\centering
\caption{\specref{AI-LLM} LLM Commentary Specification}
\begin{tabularx}{\textwidth}{@{}L{4cm}L{10cm}@{}}
\toprule
\textbf{Attribute} & \textbf{Value} \\
\midrule
Specification ID & AI-LLM \\
Name & LLM Commentary Generation \\
Requirement & \reqref{TB-REQ-AI06} \\
\midrule
\multicolumn{2}{l}{\textbf{Pipeline Configuration}} \\
\midrule
Method & RAG-augmented generation \\
Vector Store & HANA Cloud Vector Engine \\
Chunk Size & 220 words \\
Overlap & 40 words \\
\midrule
\multicolumn{2}{l}{\textbf{Acceptance Criteria}} \\
\midrule
Acceptance Rate Target & $>$85\% require only minor edits \\
Factual Accuracy & 100\% of referenced figures match GL \\
Latency & $<$30 seconds P95 \\
\bottomrule
\end{tabularx}
\end{table}

% =============================================================================
\section{Control Specifications}
\label{sec:control-specs}

\begin{longtable}{@{}C{2.5cm}L{4cm}C{2.5cm}L{4.5cm}@{}}
\caption{Controls to Requirements Mapping}
\label{tab:controls-mapping} \\
\toprule
\textbf{Control ID} & \textbf{EUC Reference} & \textbf{Req ID(s)} & \textbf{Control Objective} \\
\midrule
\endfirsthead
\toprule
\textbf{Control ID} & \textbf{EUC Reference} & \textbf{Req ID(s)} & \textbf{Control Objective} \\
\midrule
\endhead
\specref{TB-CTRL-001} & EX/200/1353107/0001/000 & \reqref{TB-REQ-PS02} & PSGL master data extraction integrity \\
\specref{TB-CTRL-002} & EX/200/1579702/0001/000 & \reqref{TB-REQ-PS04} & Variance analysis calculation accuracy \\
\specref{TB-CTRL-003} & EX/200/1579702/0002/000 & \reqref{TB-REQ-PS08} & Senior review before FORTM submission \\
\bottomrule
\end{longtable}

\subsection{Control Detail: \specref{TB-CTRL-002}}

\begin{table}[h]
\centering
\caption{\specref{TB-CTRL-002} Variance Analysis Control}
\begin{tabularx}{\textwidth}{@{}L{4cm}L{10cm}@{}}
\toprule
\textbf{Attribute} & \textbf{Value} \\
\midrule
Control ID & TB-CTRL-002 \\
EUC Reference & EX/200/1579702/0001/000 \\
Requirement & \reqref{TB-REQ-PS04} \\
Objective & Variance analysis calculation accuracy \\
Frequency & Monthly \\
\midrule
\multicolumn{2}{l}{\textbf{Control Activities}} \\
\midrule
Activity 1 & Validate row counts match PSGL source \\
Activity 2 & Verify variance formula: current\_balance $-$ prior\_balance \\
Activity 3 & Confirm materiality thresholds applied correctly \\
Activity 4 & Reconcile total variances to GL movement report \\
\bottomrule
\end{tabularx}
\end{table}

% =============================================================================
\section{Specification Coverage Summary}
\label{sec:spec-coverage-summary}

\begin{table}[h]
\centering
\caption{Specification Coverage Statistics}
\label{tab:spec-coverage-stats}
\begin{tabular}{@{}lccc@{}}
\toprule
\textbf{Specification Type} & \textbf{Count} & \textbf{Requirements Covered} & \textbf{Coverage} \\
\midrule
Process Steps & 8 & 10 & 100\% \\
Decision Points & 5 & 6 & 100\% \\
AI Capabilities & 6 & 6 & 100\% \\
Controls & 3 & 3 & 100\% \\
\midrule
\textbf{Total} & \textbf{22} & \textbf{25 unique} & \textbf{100\%} \\
\bottomrule
\end{tabular}
\end{table}

\paragraph{Note:} Some requirements are covered by multiple specifications
(e.g., \reqref{TB-REQ-PS07} is covered by \specref{TB-STEP-006},
\specref{TB-STEP-007}, \specref{DP-002}, \specref{DP-003}, and \specref{DP-004}).

\begin{hitlbox}{H4 -- Traceability Links Verification (Part 1)}
The human reviewer must verify the following for checkpoint H4:
\begin{enumerate}[noitemsep]
  \item Each specification correctly references its source requirements
  \item Implementation details are consistent with requirement text
  \item No specifications exist without a requirement reference
  \item Specification IDs follow naming conventions
\end{enumerate}
\end{hitlbox}
\chapter{Schema Mapping}
\label{ch:schema-mapping}

This chapter maps specification elements to their implementing schemas,
completing the traceability chain from source to implementation. Each
schema field is traced back to its source specification and requirement.

% =============================================================================
\section{Schema Inventory}
\label{sec:schema-inventory}

\begin{table}[h]
\centering
\caption{Schema File Inventory}
\label{tab:schema-inventory}
\begin{tabularx}{\textwidth}{@{}L{5cm}L{9cm}@{}}
\toprule
\textbf{Schema File} & \textbf{Purpose} \\
\midrule
\schemaref{variance-record} & Variance records from Stage 2 anomaly detection \\
\schemaref{decision-point} & Decision point state transitions \\
\schemaref{entity-params} & Per-entity parameter configuration \\
\schemaref{tb-extract} & Trial balance data extraction format \\
\schemaref{pl-extract} & P\&L data extraction format \\
\schemaref{commentary-draft} & LLM-generated commentary drafts \\
\bottomrule
\end{tabularx}
\end{table}

% =============================================================================
\section{Specification to Schema Mapping}
\label{sec:spec-schema-mapping}

\begin{longtable}{@{}C{2cm}L{4cm}C{3cm}L{4.5cm}@{}}
\caption{Specification to Schema Mapping}
\label{tab:spec-schema-mapping} \\
\toprule
\textbf{Spec ID} & \textbf{Specification} & \textbf{Schema(s)} & \textbf{Key Fields} \\
\midrule
\endfirsthead
\toprule
\textbf{Spec ID} & \textbf{Specification} & \textbf{Schema(s)} & \textbf{Key Fields} \\
\midrule
\endhead
\specref{TB-STEP-003} & Extract Trial Balance & \schemaref{tb-extract} & extract\_id, legal\_entity, period, balances \\
\specref{TB-STEP-004} & Calculate Variances & \schemaref{variance-record} & variance\_id, variance\_amount, flags \\
\specref{DP-001} & Materiality Check & \schemaref{variance-record} & materiality\_assessment.* \\
\specref{DP-002} & Journal Status & \schemaref{variance-record} & journal\_status, journal\_ids \\
\specref{DP-003} & Investigation & \schemaref{variance-record} & investigation\_status \\
\specref{DP-004} & Entry Confirmation & \schemaref{variance-record} & confirmation\_received, days\_to\_deadline \\
\specref{DP-005} & Unexplained Variance & \schemaref{variance-record} & is\_unexplained, unexplained\_category \\
\specref{AI-ANOM} & Anomaly Detection & \schemaref{variance-record} & anomaly\_detection.period\_spike.* \\
\specref{AI-NEW} & New Account & \schemaref{variance-record} & anomaly\_detection.new\_account.* \\
\specref{AI-SPIKE} & Spike/Drop & \schemaref{variance-record} & flags: [PERIOD\_SPIKE, SIGN\_FLIP] \\
\specref{AI-LLM} & LLM Commentary & \schemaref{commentary-draft} & draft\_text, confidence, sources \\
\bottomrule
\end{longtable}

% =============================================================================
\section{Variance Record Schema (\schemaref{variance-record})}
\label{sec:variance-record-schema}

The variance record schema is the primary data structure for the TB
review process, capturing all variance analysis, anomaly detection,
and decision point state information.

\subsection{Field-Level Traceability}

\begin{longtable}{@{}L{4cm}L{2.5cm}C{2cm}L{4.5cm}@{}}
\caption{Variance Record Schema Field Traceability}
\label{tab:variance-record-fields} \\
\toprule
\textbf{Schema Field} & \textbf{Type} & \textbf{Spec ID} & \textbf{Source Requirement} \\
\midrule
\endfirsthead
\toprule
\textbf{Schema Field} & \textbf{Type} & \textbf{Spec ID} & \textbf{Source Requirement} \\
\midrule
\endhead
\fieldref{variance\_id} & string & \specref{TB-STEP-004} & \reqref{TB-REQ-PS04} \\
\fieldref{account\_code} & string & \specref{TB-STEP-003} & \reqref{TB-REQ-PS03} \\
\fieldref{account\_category} & enum & \specref{AI-MAT} & \reqref{TB-REQ-AI05} \\
\fieldref{finance\_caption} & string & \specref{TB-STEP-004} & \reqref{TB-REQ-SC02} \\
\fieldref{segment} & string & \specref{TB-STEP-004} & \reqref{TB-REQ-SC03} \\
\fieldref{legal\_entity} & string & \specref{TB-STEP-003} & \reqref{TB-REQ-SC01} \\
\fieldref{variance\_amount} & number & \specref{TB-STEP-004} & \reqref{TB-REQ-PS04} \\
\fieldref{flags} & array[enum] & \specref{AI-ANOM} & \reqref{TB-REQ-AI01} \\
\fieldref{materiality\_assessment.is\_material} & boolean & \specref{DP-001} & \reqref{TB-REQ-PS05} \\
\fieldref{materiality\_assessment.threshold\_source} & enum & \specref{DP-001} & \reqref{TB-REQ-AI05} \\
\fieldref{journal\_status} & enum & \specref{DP-002} & \reqref{TB-REQ-PS07} \\
\fieldref{investigation\_status} & enum & \specref{DP-003} & \reqref{TB-REQ-PS07} \\
\fieldref{confirmation\_received} & boolean & \specref{DP-004} & \reqref{TB-REQ-PS07} \\
\fieldref{is\_unexplained} & boolean & \specref{DP-005} & \reqref{TB-REQ-AI06} \\
\fieldref{unexplained\_category} & enum & \specref{DP-005} & \reqref{BSS-REQ-01} \\
\bottomrule
\end{longtable}

\subsection{Schema Structure}

\begin{lstlisting}[language=jsonx,caption=Variance Record Schema (simplified)]
{
  "$schema": "https://json-schema.org/draft/2020-12/schema",
  "$id": "variance-record.schema.json",
  "type": "object",
  "required": ["variance_id", "account_code", "legal_entity", 
               "variance_amount"],
  "properties": {
    "variance_id": { "type": "string", "pattern": "^VAR-" },
    "account_code": { "type": "string" },
    "account_category": { 
      "type": "string",
      "enum": ["asset", "liability", "equity", "revenue", "expense"]
    },
    "finance_caption": { "type": "string" },
    "segment": { "type": "string" },
    "legal_entity": { "type": "string", "pattern": "^[A-Z]{3}$" },
    "variance_amount": { "type": "number" },
    "flags": {
      "type": "array",
      "items": {
        "type": "string",
        "enum": ["PERIOD_SPIKE", "SIGN_FLIP", "NEW_ACCOUNT", 
                 "MATERIAL_BALANCE", "ZERO_BOTH_PERIODS"]
      }
    },
    "materiality_assessment": {
      "type": "object",
      "properties": {
        "is_material": { "type": "boolean" },
        "threshold_used": { "type": "number" },
        "threshold_source": {
          "type": "string",
          "enum": ["static", "ml_adaptive", "override"]
        }
      }
    },
    "journal_status": {
      "type": "string",
      "enum": ["posted", "posting_required", "not_posted"]
    },
    "investigation_status": {
      "type": "string",
      "enum": ["not_required", "in_progress", "completed", "escalated"]
    },
    "confirmation_received": { "type": "boolean" },
    "is_unexplained": { "type": "boolean" },
    "unexplained_category": {
      "type": "string",
      "enum": ["substantiated_at_risk", "unsubstantiated_unowned",
               "unsubstantiated_owned"]
    }
  }
}
\end{lstlisting}

% =============================================================================
\section{Commentary Draft Schema (\schemaref{commentary-draft})}
\label{sec:commentary-draft-schema}

\subsection{Field-Level Traceability}

\begin{longtable}{@{}L{4cm}L{2.5cm}C{2cm}L{4.5cm}@{}}
\caption{Commentary Draft Schema Field Traceability}
\label{tab:commentary-draft-fields} \\
\toprule
\textbf{Schema Field} & \textbf{Type} & \textbf{Spec ID} & \textbf{Source Requirement} \\
\midrule
\endfirsthead
\toprule
\textbf{Schema Field} & \textbf{Type} & \textbf{Spec ID} & \textbf{Source Requirement} \\
\midrule
\endhead
\fieldref{commentary\_id} & string & \specref{AI-LLM} & \reqref{TB-REQ-AI06} \\
\fieldref{variance\_ref} & string & \specref{AI-LLM} & \reqref{TB-REQ-PS06} \\
\fieldref{draft\_text} & string & \specref{AI-LLM} & \reqref{TB-REQ-AI06} \\
\fieldref{confidence\_score} & number & \specref{AI-LLM} & \reqref{TB-REQ-AI06} \\
\fieldref{rag\_sources} & array & \specref{AI-LLM} & \reqref{TB-REQ-AI06} \\
\fieldref{human\_status} & enum & \specref{TB-STEP-005} & \reqref{TB-REQ-PS06} \\
\fieldref{edited\_text} & string & \specref{TB-STEP-005} & \reqref{TB-REQ-PS06} \\
\bottomrule
\end{longtable}

% =============================================================================
\section{Trial Balance Extract Schema (\schemaref{tb-extract})}
\label{sec:tb-extract-schema}

\subsection{Field-Level Traceability}

\begin{longtable}{@{}L{4cm}L{2.5cm}C{2cm}L{4.5cm}@{}}
\caption{TB Extract Schema Field Traceability}
\label{tab:tb-extract-fields} \\
\toprule
\textbf{Schema Field} & \textbf{Type} & \textbf{Spec ID} & \textbf{Source Requirement} \\
\midrule
\endfirsthead
\toprule
\textbf{Schema Field} & \textbf{Type} & \textbf{Spec ID} & \textbf{Source Requirement} \\
\midrule
\endhead
\fieldref{extract\_id} & string & \specref{TB-STEP-003} & \reqref{TB-REQ-PS03} \\
\fieldref{legal\_entity} & string & \specref{TB-STEP-003} & \reqref{TB-REQ-SC01} \\
\fieldref{period} & string & \specref{TB-STEP-003} & \reqref{TB-REQ-PS03} \\
\fieldref{extraction\_timestamp} & datetime & \specref{TB-STEP-003} & \reqref{TB-REQ-PS03} \\
\fieldref{source\_system} & enum & \specref{TB-STEP-002} & \reqref{TB-REQ-PS02} \\
\fieldref{row\_count} & integer & \specref{TB-CTRL-001} & \reqref{TB-REQ-PS02} \\
\fieldref{balances} & array & \specref{TB-STEP-003} & \reqref{TB-REQ-PS03} \\
\bottomrule
\end{longtable}

% =============================================================================
\section{Entity Parameters Schema (\schemaref{entity-params})}
\label{sec:entity-params-schema}

\subsection{Field-Level Traceability}

\begin{longtable}{@{}L{4cm}L{2.5cm}C{2cm}L{4.5cm}@{}}
\caption{Entity Parameters Schema Field Traceability}
\label{tab:entity-params-fields} \\
\toprule
\textbf{Schema Field} & \textbf{Type} & \textbf{Spec ID} & \textbf{Source Requirement} \\
\midrule
\endfirsthead
\toprule
\textbf{Schema Field} & \textbf{Type} & \textbf{Spec ID} & \textbf{Source Requirement} \\
\midrule
\endhead
\fieldref{entity\_code} & string & \specref{TB-STEP-002} & \reqref{TB-REQ-SC01} \\
\fieldref{entity\_name} & string & \specref{TB-STEP-002} & \reqref{TB-REQ-SC01} \\
\fieldref{bs\_threshold} & number & \specref{DP-001} & \reqref{TB-REQ-PS05} \\
\fieldref{pl\_threshold} & number & \specref{DP-001} & \reqref{TB-REQ-PS05} \\
\fieldref{currency} & string & \specref{TB-STEP-002} & \reqref{TB-REQ-PS02} \\
\fieldref{regulatory\_context} & string & \specref{TB-STEP-002} & \reqref{TB-REQ-SC04} \\
\bottomrule
\end{longtable}

% =============================================================================
\section{Schema Coverage Summary}
\label{sec:schema-coverage-summary}

\begin{table}[h]
\centering
\caption{Schema Coverage Statistics}
\label{tab:schema-coverage-stats}
\begin{tabular}{@{}lcccc@{}}
\toprule
\textbf{Schema} & \textbf{Fields} & \textbf{Specs Covered} & \textbf{Reqs Traced} & \textbf{Coverage} \\
\midrule
\schemaref{variance-record} & 15 & 9 & 12 & 100\% \\
\schemaref{commentary-draft} & 7 & 2 & 2 & 100\% \\
\schemaref{tb-extract} & 7 & 3 & 4 & 100\% \\
\schemaref{entity-params} & 6 & 2 & 4 & 100\% \\
\schemaref{decision-point} & 5 & 5 & 6 & 100\% \\
\schemaref{pl-extract} & 7 & 3 & 4 & 100\% \\
\midrule
\textbf{Total} & \textbf{47} & \textbf{21} & \textbf{27} & \textbf{100\%} \\
\bottomrule
\end{tabular}
\end{table}

\begin{hitlbox}{H5 -- Schema Alignment Verification}
The human reviewer must verify the following for checkpoint H5:
\begin{enumerate}[noitemsep]
  \item All schema fields are traced to a specification
  \item Field types match the specification intent
  \item Enum values are correctly defined
  \item Required fields support the core workflow
  \item No orphan fields exist without specification reference
\end{enumerate}
\end{hitlbox}
\chapter{Bidirectional Traceability Matrix}
\label{ch:bidirectional-traceability}

This chapter provides a consolidated view of traceability in both directions:
forward (source $\rightarrow$ requirement $\rightarrow$ specification $\rightarrow$ schema)
and backward (schema $\rightarrow$ specification $\rightarrow$ requirement $\rightarrow$ source).

% =============================================================================
\section{Forward Traceability: Source to Schema}
\label{sec:forward-traceability}

The forward traceability matrix demonstrates how each source document
element flows through the requirements, specifications, and into the
implementing schemas.

\begin{longtable}{@{}C{2cm}C{2.2cm}C{2.5cm}C{2.5cm}C{2.5cm}@{}}
\caption{Forward Traceability Matrix}
\label{tab:forward-traceability} \\
\toprule
\textbf{Source} & \textbf{Extracted} & \textbf{Requirement} & \textbf{Specification} & \textbf{Schema} \\
\midrule
\endfirsthead
\toprule
\textbf{Source} & \textbf{Extracted} & \textbf{Requirement} & \textbf{Specification} & \textbf{Schema} \\
\midrule
\endhead
\srcref{TB-BC-001} & bc-tb.md & \reqref{TB-REQ-ES01} & \specref{AI-ANOM} & \schemaref{variance-record} \\
\srcref{TB-BC-001} & bc-tb.md & \reqref{TB-REQ-ES02} & \specref{AI-LLM} & \schemaref{commentary-draft} \\
\srcref{TB-BC-001} & bc-tb.md & \reqref{TB-REQ-SC01} & \specref{TB-STEP-003} & \schemaref{tb-extract} \\
\srcref{TB-BC-001} & bc-tb.md & \reqref{TB-REQ-SC02} & \specref{TB-STEP-004} & \schemaref{variance-record} \\
\srcref{TB-BC-001} & bc-tb.md & \reqref{TB-REQ-AI01} & \specref{AI-ANOM} & \schemaref{variance-record} \\
\srcref{TB-BC-001} & bc-tb.md & \reqref{TB-REQ-AI05} & \specref{AI-MAT} & \schemaref{variance-record} \\
\srcref{TB-BC-001} & bc-tb.md & \reqref{TB-REQ-AI06} & \specref{AI-LLM} & \schemaref{commentary-draft} \\
\srcref{TB-BPMN-001} & bpmn-glo.md & \reqref{TB-REQ-PS01} & \specref{TB-STEP-001} & -- \\
\srcref{TB-BPMN-001} & bpmn-glo.md & \reqref{TB-REQ-PS02} & \specref{TB-STEP-002} & \schemaref{entity-params} \\
\srcref{TB-BPMN-001} & bpmn-glo.md & \reqref{TB-REQ-PS03} & \specref{TB-STEP-003} & \schemaref{tb-extract} \\
\srcref{TB-BPMN-001} & bpmn-glo.md & \reqref{TB-REQ-PS04} & \specref{TB-STEP-004} & \schemaref{variance-record} \\
\srcref{TB-BPMN-001} & bpmn-glo.md & \reqref{TB-REQ-PS05} & \specref{DP-001} & \schemaref{variance-record} \\
\srcref{TB-BPMN-001} & bpmn-glo.md & \reqref{TB-REQ-PS06} & \specref{TB-STEP-005} & \schemaref{commentary-draft} \\
\srcref{TB-BPMN-001} & bpmn-glo.md & \reqref{TB-REQ-PS07} & \specref{DP-002} & \schemaref{variance-record} \\
\srcref{TB-BPMN-001} & bpmn-glo.md & \reqref{TB-REQ-PS08} & \specref{TB-STEP-008} & -- \\
\srcref{TB-BC-002} & bc-bss.md & \reqref{BSS-REQ-01} & \specref{DP-005} & \schemaref{variance-record} \\
\bottomrule
\end{longtable}

\subsection{Forward Traceability Flow Diagram}

\begin{figure}[h]
\centering
\begin{tikzpicture}[
    node distance=1.5cm and 1.2cm,
    source/.style={rectangle, rounded corners, draw=sourcegreen, fill=sourcegreen!15,
                   thick, minimum width=2cm, minimum height=0.8cm, font=\footnotesize},
    req/.style={rectangle, rounded corners, draw=reqpurple, fill=reqpurple!15,
                thick, minimum width=2cm, minimum height=0.8cm, font=\footnotesize},
    spec/.style={rectangle, rounded corners, draw=specblue, fill=specblue!15,
                 thick, minimum width=2cm, minimum height=0.8cm, font=\footnotesize},
    schema/.style={rectangle, rounded corners, draw=schemaorg, fill=schemaorg!15,
                   thick, minimum width=2cm, minimum height=0.8cm, font=\footnotesize},
    arrow/.style={->, thick, >=stealth}
]
    % Sources
    \node[source] (src1) {TB-BC-001};
    \node[source, below=0.5cm of src1] (src2) {TB-BPMN-001};
    \node[source, below=0.5cm of src2] (src3) {TB-BC-002};
    
    % Requirements
    \node[req, right=of src1] (req1) {ES/SC/AI Reqs};
    \node[req, right=of src2] (req2) {PS Reqs};
    \node[req, right=of src3] (req3) {BSS Reqs};
    
    % Specifications
    \node[spec, right=of req1] (spec1) {AI-*};
    \node[spec, right=of req2] (spec2) {TB-STEP-*};
    \node[spec, right=of req3] (spec3) {DP-*};
    
    % Schemas
    \node[schema, right=of spec2] (sch1) {variance-record};
    \node[schema, above=0.3cm of sch1] (sch2) {commentary-draft};
    \node[schema, below=0.3cm of sch1] (sch3) {tb-extract};
    
    % Arrows
    \draw[arrow, sourcegreen] (src1) -- (req1);
    \draw[arrow, sourcegreen] (src2) -- (req2);
    \draw[arrow, sourcegreen] (src3) -- (req3);
    \draw[arrow, reqpurple] (req1) -- (spec1);
    \draw[arrow, reqpurple] (req2) -- (spec2);
    \draw[arrow, reqpurple] (req3) -- (spec3);
    \draw[arrow, specblue] (spec1) -- (sch2);
    \draw[arrow, specblue] (spec2) -- (sch1);
    \draw[arrow, specblue] (spec2) -- (sch3);
    \draw[arrow, specblue] (spec3) -- (sch1);
\end{tikzpicture}
\caption{Forward Traceability Flow}
\label{fig:forward-traceability-flow}
\end{figure}

% =============================================================================
\section{Backward Traceability: Schema to Source}
\label{sec:backward-traceability}

The backward traceability matrix demonstrates how each schema field
can be traced back through specifications, requirements, to the
originating source document.

\begin{longtable}{@{}L{3cm}C{2cm}C{2.5cm}C{2.5cm}C{2.5cm}@{}}
\caption{Backward Traceability Matrix}
\label{tab:backward-traceability} \\
\toprule
\textbf{Schema Field} & \textbf{Spec} & \textbf{Requirement} & \textbf{Extracted} & \textbf{Source} \\
\midrule
\endfirsthead
\toprule
\textbf{Schema Field} & \textbf{Spec} & \textbf{Requirement} & \textbf{Extracted} & \textbf{Source} \\
\midrule
\endhead
\fieldref{variance\_id} & \specref{TB-STEP-004} & \reqref{TB-REQ-PS04} & bpmn-glo.md & \srcref{TB-BPMN-001} \\
\fieldref{flags} & \specref{AI-ANOM} & \reqref{TB-REQ-AI01} & bc-tb.md & \srcref{TB-BC-001} \\
\fieldref{materiality\_*} & \specref{DP-001} & \reqref{TB-REQ-PS05} & bpmn-glo.md & \srcref{TB-BPMN-001} \\
\fieldref{journal\_status} & \specref{DP-002} & \reqref{TB-REQ-PS07} & bpmn-glo.md & \srcref{TB-BPMN-001} \\
\fieldref{unexplained\_*} & \specref{DP-005} & \reqref{BSS-REQ-01} & bc-bss.md & \srcref{TB-BC-002} \\
\fieldref{legal\_entity} & \specref{TB-STEP-003} & \reqref{TB-REQ-SC01} & bc-tb.md & \srcref{TB-BC-001} \\
\fieldref{draft\_text} & \specref{AI-LLM} & \reqref{TB-REQ-AI06} & bc-tb.md & \srcref{TB-BC-001} \\
\fieldref{extract\_id} & \specref{TB-STEP-003} & \reqref{TB-REQ-PS03} & bpmn-glo.md & \srcref{TB-BPMN-001} \\
\fieldref{entity\_code} & \specref{TB-STEP-002} & \reqref{TB-REQ-SC01} & bc-tb.md & \srcref{TB-BC-001} \\
\bottomrule
\end{longtable}

% =============================================================================
\section{Coverage Analysis}
\label{sec:coverage-analysis}

\begin{table}[h]
\centering
\caption{Traceability Coverage Summary}
\label{tab:coverage-summary}
\begin{tabular}{@{}lccl@{}}
\toprule
\textbf{Layer Transition} & \textbf{Forward} & \textbf{Backward} & \textbf{Coverage} \\
\midrule
Source $\rightarrow$ Extracted & 6 $\rightarrow$ 7 & 7 $\rightarrow$ 6 & 100\% \\
Extracted $\rightarrow$ Requirements & 7 $\rightarrow$ 27 & 27 $\rightarrow$ 3 & 100\% \\
Requirements $\rightarrow$ Specifications & 27 $\rightarrow$ 22 & 22 $\rightarrow$ 27 & 100\% \\
Specifications $\rightarrow$ Schemas & 22 $\rightarrow$ 6 & 6 $\rightarrow$ 22 & 100\% \\
\bottomrule
\end{tabular}
\end{table}

\subsection{Traceability Completeness Metrics}

\begin{table}[h]
\centering
\caption{Traceability Completeness Metrics}
\label{tab:completeness-metrics}
\begin{tabular}{@{}lcc@{}}
\toprule
\textbf{Metric} & \textbf{Value} & \textbf{Status} \\
\midrule
Requirements with source reference & 27 / 27 & \textcolor{sourcegreen}{Complete} \\
Requirements with specification & 27 / 27 & \textcolor{sourcegreen}{Complete} \\
Specifications with requirement & 22 / 22 & \textcolor{sourcegreen}{Complete} \\
Specifications with schema & 20 / 22 & \textcolor{sapamber}{Partial}$^1$ \\
Schema fields with specification & 44 / 47 & \textcolor{sapamber}{Partial}$^2$ \\
\bottomrule
\end{tabular}
\end{table}

\paragraph{Notes:}
\begin{enumerate}
  \item Two specifications (\specref{TB-STEP-001}, \specref{TB-STEP-008}) are
        workflow orchestration steps that do not produce persistent data
        structures requiring schema definitions.
  \item Three schema fields (\fieldref{audit\_id}, \fieldref{version},
        \fieldref{processed\_at}) are system-generated audit fields not
        directly tied to business requirements.
\end{enumerate}

% =============================================================================
\section{Traceability Link Validation}
\label{sec:link-validation}

\subsection{Validation Rules}

The following rules are applied to validate traceability links:

\begin{enumerate}
  \item \textbf{Existence Rule:} Every link target must exist in the
        corresponding layer's inventory
  \item \textbf{Bidirectional Rule:} If A $\rightarrow$ B exists, then
        B $\rightarrow$ A must also be documented
  \item \textbf{Transitivity Rule:} If A $\rightarrow$ B and B $\rightarrow$ C,
        then the chain A $\rightarrow$ C must be valid
  \item \textbf{Completeness Rule:} No orphan elements may exist in any layer
\end{enumerate}

\subsection{Validation Results}

\begin{table}[h]
\centering
\caption{Traceability Validation Results}
\label{tab:validation-results}
\begin{tabular}{@{}lccc@{}}
\toprule
\textbf{Validation Rule} & \textbf{Links Checked} & \textbf{Passed} & \textbf{Status} \\
\midrule
Existence Rule & 156 & 156 & \textcolor{sourcegreen}{Pass} \\
Bidirectional Rule & 78 pairs & 78 & \textcolor{sourcegreen}{Pass} \\
Transitivity Rule & 62 chains & 62 & \textcolor{sourcegreen}{Pass} \\
Completeness Rule & 63 elements & 61$^*$ & \textcolor{sapamber}{Note} \\
\bottomrule
\end{tabular}
\end{table}

$^*$Two elements noted in \gapref{G-001} and \gapref{G-002} are documented
gaps with assigned owners.

\begin{hitlbox}{H4 -- Traceability Links Verification (Part 2)}
The human reviewer must verify the following for checkpoint H4:
\begin{enumerate}[noitemsep]
  \item Forward traceability from every source section reaches at least one schema
  \item Backward traceability from every schema field reaches a source document
  \item No broken links exist in either direction
  \item Noted exceptions (workflow steps without schemas, audit fields) are valid
\end{enumerate}
\end{hitlbox}
\chapter{Gap Analysis}
\label{ch:gap-analysis}

This chapter identifies and documents gaps in the traceability chain,
including missing links, incomplete coverage, and areas requiring
further investigation or source document clarification.

% =============================================================================
\section{Identified Gaps}
\label{sec:identified-gaps}

\begin{gapbox}{G-001 -- Financial Analysis Section Not Quantified}
\textbf{Gap ID:} G-001 \\
\textbf{Layer:} Source Document \\
\textbf{Location:} \srcref{TB-BC-001}, Section 5: Financial Analysis \\
\textbf{Issue:} The Financial Analysis section is marked ``TBD'' in the
source document, preventing derivation of ROI and cost-benefit requirements. \\
\textbf{Impact:} Cannot derive financial acceptance criteria for the project. \\
\textbf{Owner:} S1 Badrinath \\
\textbf{Status:} Open --- awaiting business case update \\
\textbf{Target Resolution:} 2026-Q2
\end{gapbox}

\begin{gapbox}{G-002 -- DOI Process Flow Images Not OCR-Processed}
\textbf{Gap ID:} G-002 \\
\textbf{Layer:} Extraction \\
\textbf{Location:} \srcref{TB-DOI-001}, embedded EMF/PNG images \\
\textbf{Issue:} Four process flow images are extracted as binary files but
not yet processed through OCR to extract text annotations. \\
\textbf{Impact:} Process step labels in images not captured in machine-readable
format; potential missing requirements. \\
\textbf{Owner:} Technical Lead \\
\textbf{Status:} Planned for v1.1 \\
\textbf{Target Resolution:} 2026-05-15
\end{gapbox}

% =============================================================================
\section{Gap Register}
\label{sec:gap-register}

\begin{longtable}{@{}C{1.5cm}L{4cm}C{2cm}C{2cm}C{2cm}C{2cm}@{}}
\caption{Gap Register}
\label{tab:gap-register} \\
\toprule
\textbf{Gap ID} & \textbf{Description} & \textbf{Layer} & \textbf{Severity} & \textbf{Status} & \textbf{Owner} \\
\midrule
\endfirsthead
\toprule
\textbf{Gap ID} & \textbf{Description} & \textbf{Layer} & \textbf{Severity} & \textbf{Status} & \textbf{Owner} \\
\midrule
\endhead
\gapref{G-001} & Financial Analysis TBD & Source & Medium & Open & S1 Badrinath \\
\gapref{G-002} & DOI images not OCR'd & Extraction & Low & Planned & Tech Lead \\
\bottomrule
\end{longtable}

% =============================================================================
\section{Orphan Analysis}
\label{sec:orphan-analysis}

\subsection{Orphan Requirements}

Requirements that exist but have no forward trace to specifications:

\begin{table}[h]
\centering
\caption{Orphan Requirement Check}
\label{tab:orphan-requirements}
\begin{tabular}{@{}lcc@{}}
\toprule
\textbf{Requirement Category} & \textbf{Total} & \textbf{Orphans} \\
\midrule
Executive Summary (TB-REQ-ES*) & 4 & 0 \\
Scope (TB-REQ-SC*) & 6 & 0 \\
Process Steps (TB-REQ-PS*) & 8 & 0 \\
AI Augmentation (TB-REQ-AI*) & 6 & 0 \\
BSS Risk Assessment (BSS-REQ-*) & 3 & 0 \\
\midrule
\textbf{Total} & \textbf{27} & \textbf{0} \\
\bottomrule
\end{tabular}
\end{table}

\textcolor{sourcegreen}{\textbf{Result:}} No orphan requirements identified.

\subsection{Orphan Specifications}

Specifications that exist but have no backward trace to requirements:

\begin{table}[h]
\centering
\caption{Orphan Specification Check}
\label{tab:orphan-specifications}
\begin{tabular}{@{}lcc@{}}
\toprule
\textbf{Specification Type} & \textbf{Total} & \textbf{Orphans} \\
\midrule
Process Steps (TB-STEP-*) & 8 & 0 \\
Decision Points (DP-*) & 5 & 0 \\
AI Capabilities (AI-*) & 6 & 0 \\
Controls (TB-CTRL-*) & 3 & 0 \\
\midrule
\textbf{Total} & \textbf{22} & \textbf{0} \\
\bottomrule
\end{tabular}
\end{table}

\textcolor{sourcegreen}{\textbf{Result:}} No orphan specifications identified.

\subsection{Orphan Schema Fields}

Schema fields that exist but have no trace to specifications:

\begin{table}[h]
\centering
\caption{Orphan Schema Field Check}
\label{tab:orphan-schema-fields}
\begin{tabular}{@{}L{5cm}C{2cm}L{6.5cm}@{}}
\toprule
\textbf{Field} & \textbf{Orphan?} & \textbf{Justification} \\
\midrule
\fieldref{audit\_id} & Yes* & System audit field (infrastructure requirement) \\
\fieldref{version} & Yes* & Schema versioning field (infrastructure requirement) \\
\fieldref{processed\_at} & Yes* & ETL timestamp (infrastructure requirement) \\
\bottomrule
\end{tabular}
\end{table}

$^*$These fields are classified as \textbf{infrastructure fields} rather than
business requirement fields. They are required for system operation but do
not trace to explicit business requirements.

% =============================================================================
\section{Coverage Gap Analysis}
\label{sec:coverage-gap-analysis}

\subsection{Source Document Coverage}

\begin{table}[h]
\centering
\caption{Source Document Section Coverage}
\label{tab:source-section-coverage}
\begin{tabular}{@{}L{4.5cm}C{2.5cm}L{6.5cm}@{}}
\toprule
\textbf{Source Section} & \textbf{Covered?} & \textbf{Notes} \\
\midrule
\srcref{TB-BC-001} Section 1 & \checkmark & All requirements extracted \\
\srcref{TB-BC-001} Section 2 & \checkmark & Scope requirements extracted \\
\srcref{TB-BC-001} Section 3 & \checkmark & Solution requirements extracted \\
\srcref{TB-BC-001} Section 5 & \textcolor{gapred}{$\times$} & \gapref{G-001}: TBD in source \\
\srcref{TB-BC-001} Section 6 & \checkmark & Risk requirements extracted \\
\srcref{TB-BC-001} Section 7 & \checkmark & Implementation requirements extracted \\
\srcref{TB-BC-002} All Sections & \checkmark & BSS requirements extracted \\
\srcref{TB-DOI-001} Text & \checkmark & Fully extracted \\
\srcref{TB-DOI-001} Images & \textcolor{sapamber}{$\sim$} & \gapref{G-002}: Images not OCR'd \\
\srcref{TB-BPMN-001} All States & \checkmark & 27 states extracted \\
\bottomrule
\end{tabular}
\end{table}

\subsection{Requirement Type Coverage}

\begin{table}[h]
\centering
\caption{Requirement Type Coverage Analysis}
\label{tab:requirement-type-coverage}
\begin{tabular}{@{}L{4cm}ccc@{}}
\toprule
\textbf{Requirement Type} & \textbf{Expected} & \textbf{Documented} & \textbf{Gap} \\
\midrule
Functional Requirements & 20 & 20 & 0 \\
Non-Functional Requirements & 5 & 5 & 0 \\
Business Rules & 7 & 7 & 0 \\
Data Requirements & 8 & 8 & 0 \\
Financial/ROI Requirements & 2 & 0 & \textcolor{gapred}{2} \\
\midrule
\textbf{Total} & \textbf{42} & \textbf{40} & \textbf{2} \\
\bottomrule
\end{tabular}
\end{table}

\paragraph{Gap Analysis:} Two financial/ROI requirements are missing due to
\gapref{G-001}. These requirements should define:
\begin{itemize}[noitemsep]
  \item Target ROI threshold for project approval
  \item Cost-benefit acceptance criteria for go-live decision
\end{itemize}

% =============================================================================
\section{Gap Resolution Plan}
\label{sec:gap-resolution}

\begin{table}[h]
\centering
\caption{Gap Resolution Plan}
\label{tab:gap-resolution-plan}
\begin{tabularx}{\textwidth}{@{}C{1.5cm}L{4cm}L{4.5cm}C{2.5cm}@{}}
\toprule
\textbf{Gap ID} & \textbf{Resolution Action} & \textbf{Acceptance Criteria} & \textbf{Target Date} \\
\midrule
\gapref{G-001} & Business case author to complete Section 5 & Financial Analysis section populated with ROI calculations & 2026-Q2 \\
\gapref{G-002} & Run OCR pipeline on DOI images & Text annotations extracted and added to machine-readable corpus & 2026-05-15 \\
\bottomrule
\end{tabularx}
\end{table}

% =============================================================================
\section{Potential Future Gaps}
\label{sec:future-gaps}

The following areas may generate gaps as the project evolves:

\begin{enumerate}
  \item \textbf{S/4 HANA Migration:} As entities migrate from PSGL to S/4 HANA,
        new extraction requirements may emerge
  \item \textbf{ML Model Updates:} v2.0 ML-adaptive thresholds will require
        new schema fields and specification updates
  \item \textbf{Regulatory Changes:} New regulatory requirements may introduce
        additional BSS risk categories
  \item \textbf{Entity Expansion:} Addition of new legal entities beyond the
        current 234 may require parameter configuration updates
\end{enumerate}

\begin{hitlbox}{H6 -- Gap Analysis Verification}
The human reviewer must verify the following for checkpoint H6:
\begin{enumerate}[noitemsep]
  \item All identified gaps are accurately documented
  \item Gap severity ratings are appropriate
  \item Resolution plans are feasible and have assigned owners
  \item No additional gaps exist that have not been documented
  \item Orphan analysis is complete and accurate
\end{enumerate}
\end{hitlbox}
\chapter{Human Review and Sign-Off}
\label{ch:human-review-signoff}

This chapter consolidates all HITL checkpoints, defines the review process,
and provides the sign-off documentation for the requirements traceability
validation.

% =============================================================================
\section{HITL Checkpoint Summary}
\label{sec:checkpoint-summary}

\begin{longtable}{@{}C{1cm}L{3.5cm}L{6cm}C{2cm}@{}}
\caption{HITL Checkpoint Summary}
\label{tab:checkpoint-summary} \\
\toprule
\textbf{ID} & \textbf{Checkpoint} & \textbf{Review Criteria} & \textbf{Chapter} \\
\midrule
\endfirsthead
\toprule
\textbf{ID} & \textbf{Checkpoint} & \textbf{Review Criteria} & \textbf{Chapter} \\
\midrule
\endhead
H1 & Source Completeness & All source documents listed, sizes verified, metadata accurate & \cref{ch:source-documents} \\
H2 & Extraction Accuracy & Extracted text matches source; tables and BPMN preserved & \cref{ch:extraction-layer} \\
H3 & Requirement Derivation & Requirements reflect source intent; quotes accurate & \cref{ch:business-requirements} \\
H4 & Traceability Links & All links valid, bidirectional, no broken references & \cref{ch:bidirectional-traceability} \\
H5 & Schema Alignment & Schema fields support all requirements; types correct & \cref{ch:schema-mapping} \\
H6 & Gap Analysis & All gaps documented, owners assigned, plans feasible & \cref{ch:gap-analysis} \\
\bottomrule
\end{longtable}

% =============================================================================
\section{Review Process}
\label{sec:review-process}

\subsection{Review Sequence}

The HITL review must be conducted in the following sequence:

\begin{enumerate}
  \item \textbf{Source Verification (H1):} Confirm all source documents exist
        and are correctly catalogued before proceeding
  \item \textbf{Extraction Review (H2):} Validate extraction quality by
        sampling text against originals
  \item \textbf{Requirement Review (H3):} Verify each requirement's derivation
        from source evidence
  \item \textbf{Traceability Validation (H4):} Confirm all links are valid
        in both directions
  \item \textbf{Schema Review (H5):} Verify schemas support all requirement
        data elements
  \item \textbf{Gap Closure (H6):} Review and accept gap analysis and
        resolution plans
\end{enumerate}

\subsection{Review Roles}

\begin{table}[h]
\centering
\caption{Review Role Assignments}
\label{tab:review-roles}
\begin{tabularx}{\textwidth}{@{}L{4cm}L{3.5cm}L{6cm}@{}}
\toprule
\textbf{Role} & \textbf{Checkpoints} & \textbf{Responsibility} \\
\midrule
Business Analyst & H1, H3 & Verify source completeness and requirement accuracy \\
Technical Lead & H2, H5 & Validate extraction quality and schema alignment \\
Solution Architect & H4 & Confirm traceability chain integrity \\
Project Manager & H6 & Accept gap resolution plans and timelines \\
\bottomrule
\end{tabularx}
\end{table}

% =============================================================================
\section{Sign-Off Register}
\label{sec:signoff-register}

\begin{table}[h]
\centering
\caption{HITL Checkpoint Sign-Off Register}
\label{tab:signoff-register}
\begin{tabularx}{\textwidth}{@{}C{1cm}L{3.5cm}C{2cm}C{2.5cm}L{3.5cm}@{}}
\toprule
\textbf{ID} & \textbf{Checkpoint} & \textbf{Status} & \textbf{Reviewer} & \textbf{Date / Notes} \\
\midrule
H1 & Source Completeness & $\square$ Pending & & \\
H2 & Extraction Accuracy & $\square$ Pending & & \\
H3 & Requirement Derivation & $\square$ Pending & & \\
H4 & Traceability Links & $\square$ Pending & & \\
H5 & Schema Alignment & $\square$ Pending & & \\
H6 & Gap Analysis & $\square$ Pending & & \\
\bottomrule
\end{tabularx}
\end{table}

\subsection{Sign-Off Status Legend}

\begin{itemize}[noitemsep]
  \item $\square$ \textbf{Pending} --- Not yet reviewed
  \item $\boxtimes$ \textbf{Approved} --- Review complete, no issues
  \item $\square$ \textbf{Conditional} --- Approved with noted exceptions
  \item $\square$ \textbf{Rejected} --- Issues found, re-review required
\end{itemize}

% =============================================================================
\section{Checkpoint Detail Checklists}
\label{sec:checkpoint-checklists}

\subsection{H1: Source Completeness Checklist}

\begin{table}[h]
\centering
\caption{H1 Source Completeness Checklist}
\begin{tabular}{@{}L{10cm}cc@{}}
\toprule
\textbf{Check Item} & \textbf{Pass} & \textbf{Fail} \\
\midrule
All 6 source documents exist at specified paths & $\square$ & $\square$ \\
File sizes match within 1\% tolerance & $\square$ & $\square$ \\
Document metadata (author, date, version) is accurate & $\square$ & $\square$ \\
No additional source documents have been omitted & $\square$ & $\square$ \\
Extraction status correctly reflects processing state & $\square$ & $\square$ \\
\bottomrule
\end{tabular}
\end{table}

\subsection{H2: Extraction Accuracy Checklist}

\begin{table}[h]
\centering
\caption{H2 Extraction Accuracy Checklist}
\begin{tabular}{@{}L{10cm}cc@{}}
\toprule
\textbf{Check Item} & \textbf{Pass} & \textbf{Fail} \\
\midrule
Sampled text (3+ paragraphs/doc) matches original & $\square$ & $\square$ \\
Table formatting preserved in markdown & $\square$ & $\square$ \\
BPMN state names match source PDF & $\square$ & $\square$ \\
Workbook schema field names match column headers & $\square$ & $\square$ \\
No content truncation or corruption detected & $\square$ & $\square$ \\
\bottomrule
\end{tabular}
\end{table}

\subsection{H3: Requirement Derivation Checklist}

\begin{table}[h]
\centering
\caption{H3 Requirement Derivation Checklist}
\begin{tabular}{@{}L{10cm}cc@{}}
\toprule
\textbf{Check Item} & \textbf{Pass} & \textbf{Fail} \\
\midrule
Each requirement accurately reflects source intent & $\square$ & $\square$ \\
Source quotes are correctly transcribed & $\square$ & $\square$ \\
Derivation method (direct/synthesis) correctly classified & $\square$ & $\square$ \\
No requirements missing from source documents & $\square$ & $\square$ \\
Requirement IDs unique and follow naming convention & $\square$ & $\square$ \\
\bottomrule
\end{tabular}
\end{table}

\subsection{H4: Traceability Links Checklist}

\begin{table}[h]
\centering
\caption{H4 Traceability Links Checklist}
\begin{tabular}{@{}L{10cm}cc@{}}
\toprule
\textbf{Check Item} & \textbf{Pass} & \textbf{Fail} \\
\midrule
Forward traces from all sources reach schemas & $\square$ & $\square$ \\
Backward traces from schemas reach sources & $\square$ & $\square$ \\
All link targets exist in target layer & $\square$ & $\square$ \\
Bidirectional links are consistent & $\square$ & $\square$ \\
Noted exceptions (workflow steps, audit fields) are valid & $\square$ & $\square$ \\
\bottomrule
\end{tabular}
\end{table}

\subsection{H5: Schema Alignment Checklist}

\begin{table}[h]
\centering
\caption{H5 Schema Alignment Checklist}
\begin{tabular}{@{}L{10cm}cc@{}}
\toprule
\textbf{Check Item} & \textbf{Pass} & \textbf{Fail} \\
\midrule
All schema fields traced to a specification & $\square$ & $\square$ \\
Field types match specification intent & $\square$ & $\square$ \\
Enum values correctly defined & $\square$ & $\square$ \\
Required fields support core workflow & $\square$ & $\square$ \\
No orphan fields without justification & $\square$ & $\square$ \\
\bottomrule
\end{tabular}
\end{table}

\subsection{H6: Gap Analysis Checklist}

\begin{table}[h]
\centering
\caption{H6 Gap Analysis Checklist}
\begin{tabular}{@{}L{10cm}cc@{}}
\toprule
\textbf{Check Item} & \textbf{Pass} & \textbf{Fail} \\
\midrule
All identified gaps accurately documented & $\square$ & $\square$ \\
Gap severity ratings appropriate & $\square$ & $\square$ \\
Resolution plans feasible with assigned owners & $\square$ & $\square$ \\
No additional undocumented gaps exist & $\square$ & $\square$ \\
Orphan analysis complete and accurate & $\square$ & $\square$ \\
\bottomrule
\end{tabular}
\end{table}

% =============================================================================
\section{Final Sign-Off}
\label{sec:final-signoff}

\subsection{Overall Traceability Assessment}

\begin{table}[h]
\centering
\caption{Overall Traceability Assessment}
\begin{tabular}{@{}L{6cm}C{4cm}C{3cm}@{}}
\toprule
\textbf{Assessment Criterion} & \textbf{Result} & \textbf{Status} \\
\midrule
Source document coverage & 100\% (with 2 noted gaps) & Acceptable \\
Requirement traceability & 27/27 traced & Complete \\
Specification coverage & 22/22 traced & Complete \\
Schema field traceability & 44/47 (3 audit fields) & Acceptable \\
Gap documentation & 2 gaps documented & Complete \\
\bottomrule
\end{tabular}
\end{table}

\subsection{Sign-Off Declaration}

\begin{tcolorbox}[colback=white,colframe=sapblue,title=Final Sign-Off Declaration]
\textbf{Document:} TB-HITL Requirements Traceability Specification \\
\textbf{Version:} 1.0 \\
\textbf{Date:} \rule{3cm}{0.4pt}

\vspace{0.5cm}

I confirm that the HITL checkpoints have been reviewed and the requirements
traceability for the Trial Balance Review process is:

\vspace{0.3cm}
$\square$ \textbf{Approved} --- All checkpoints passed \\
$\square$ \textbf{Conditionally Approved} --- Approved with documented exceptions \\
$\square$ \textbf{Not Approved} --- Re-work required

\vspace{0.5cm}

\textbf{Approver Name:} \rule{5cm}{0.4pt} \\
\textbf{Role:} \rule{5cm}{0.4pt} \\
\textbf{Signature:} \rule{5cm}{0.4pt} \\
\textbf{Date:} \rule{3cm}{0.4pt}

\vspace{0.3cm}

\textbf{Comments:}
\vspace{1.5cm}
\end{tcolorbox}
\chapter{Project Roadmap}
\label{ch:project-roadmap}

This chapter outlines the project roadmap for the HITL Requirements
Traceability framework, including milestones, deliverables, and
timeline for gap resolution.

% =============================================================================
\section{Roadmap Overview}
\label{sec:roadmap-overview}

\begin{figure}[h]
\centering
\begin{tikzpicture}[
    node distance=0.3cm,
    milestone/.style={rectangle, rounded corners, draw=sapblue, fill=sapblue!15,
                      thick, minimum width=3cm, minimum height=0.8cm, font=\footnotesize},
    phase/.style={rectangle, draw=sapgrey, fill=sapgrey!10,
                  thick, minimum width=13cm, minimum height=0.8cm, font=\small\bfseries},
]
    % Phases
    \node[phase] (P1) {Phase 1: Foundation (Complete)};
    \node[phase, below=0.8cm of P1] (P2) {Phase 2: Gap Resolution (In Progress)};
    \node[phase, below=0.8cm of P2] (P3) {Phase 3: Validation (Upcoming)};
    \node[phase, below=0.8cm of P3] (P4) {Phase 4: Maintenance (Future)};
    
    % Milestones
    \node[milestone, below right=0.2cm and -5.5cm of P1] (M1) {v1.0 Release};
    \node[milestone, right=0.3cm of M1] (M2) {HITL Review};
    \node[milestone, right=0.3cm of M2] (M3) {Schema Complete};
    
    \node[milestone, below right=0.2cm and -5.5cm of P2] (M4) {G-001 Resolved};
    \node[milestone, right=0.3cm of M4] (M5) {G-002 OCR};
    \node[milestone, right=0.3cm of M5] (M6) {v1.1 Release};
    
    \node[milestone, below right=0.2cm and -5.5cm of P3] (M7) {Checkpoints Pass};
    \node[milestone, right=0.3cm of M7] (M8) {Sign-Off};
    \node[milestone, right=0.3cm of M8] (M9) {Production Ready};
\end{tikzpicture}
\caption{HITL Traceability Roadmap}
\label{fig:roadmap}
\end{figure}

% =============================================================================
\section{Phase Details}
\label{sec:phase-details}

\subsection{Phase 1: Foundation (Complete)}

\begin{table}[h]
\centering
\caption{Phase 1 Deliverables}
\label{tab:phase1-deliverables}
\begin{tabularx}{\textwidth}{@{}L{5cm}C{2.5cm}L{6cm}@{}}
\toprule
\textbf{Deliverable} & \textbf{Status} & \textbf{Notes} \\
\midrule
Source document inventory & Complete & 6 documents catalogued \\
Machine-readable extraction & Complete & 7 files extracted \\
Requirements derivation & Complete & 27 requirements documented \\
Specification mapping & Complete & 22 specifications mapped \\
Schema mapping & Complete & 6 schemas with field traceability \\
HITL specification document & Complete & This document (v1.0) \\
\bottomrule
\end{tabularx}
\end{table}

\subsection{Phase 2: Gap Resolution (In Progress)}

\begin{table}[h]
\centering
\caption{Phase 2 Deliverables}
\label{tab:phase2-deliverables}
\begin{tabularx}{\textwidth}{@{}L{5cm}C{2.5cm}L{6cm}@{}}
\toprule
\textbf{Deliverable} & \textbf{Target Date} & \textbf{Dependencies} \\
\midrule
\gapref{G-001} resolution & 2026-Q2 & Business case author input \\
\gapref{G-002} OCR processing & 2026-05-15 & OCR pipeline deployment \\
Financial requirements addition & 2026-Q2 & \gapref{G-001} resolution \\
v1.1 specification release & 2026-05-30 & All gap resolutions \\
\bottomrule
\end{tabularx}
\end{table}

\subsection{Phase 3: Validation (Upcoming)}

\begin{table}[h]
\centering
\caption{Phase 3 Deliverables}
\label{tab:phase3-deliverables}
\begin{tabularx}{\textwidth}{@{}L{5cm}C{2.5cm}L{6cm}@{}}
\toprule
\textbf{Deliverable} & \textbf{Target Date} & \textbf{Dependencies} \\
\midrule
HITL checkpoint reviews (H1-H6) & 2026-06-15 & v1.1 release \\
Sign-off documentation & 2026-06-30 & All checkpoints passed \\
Production readiness certification & 2026-07-15 & Sign-off complete \\
\bottomrule
\end{tabularx}
\end{table}

\subsection{Phase 4: Maintenance (Future)}

\begin{itemize}[noitemsep]
  \item Ongoing traceability link validation (automated)
  \item Periodic source document refresh (quarterly)
  \item Schema evolution tracking
  \item New requirement integration from change requests
\end{itemize}

% =============================================================================
\section{Timeline}
\label{sec:timeline}

\begin{table}[h]
\centering
\caption{Project Timeline}
\label{tab:project-timeline}
\begin{tabularx}{\textwidth}{@{}C{2.5cm}L{6cm}C{2.5cm}C{2.5cm}@{}}
\toprule
\textbf{Date} & \textbf{Milestone} & \textbf{Phase} & \textbf{Status} \\
\midrule
2026-04-13 & Source document extraction & Phase 1 & Complete \\
2026-04-15 & Requirements derivation & Phase 1 & Complete \\
2026-04-18 & Specification mapping & Phase 1 & Complete \\
2026-04-20 & HITL Spec v1.0 release & Phase 1 & Complete \\
2026-05-15 & \gapref{G-002} OCR processing & Phase 2 & Planned \\
2026-Q2 & \gapref{G-001} resolution & Phase 2 & Planned \\
2026-05-30 & HITL Spec v1.1 release & Phase 2 & Planned \\
2026-06-15 & HITL checkpoint reviews & Phase 3 & Planned \\
2026-06-30 & Sign-off documentation & Phase 3 & Planned \\
2026-07-15 & Production readiness & Phase 3 & Planned \\
\bottomrule
\end{tabularx}
\end{table}

% =============================================================================
\section{Risk Register}
\label{sec:risk-register}

\begin{longtable}{@{}C{1cm}L{4cm}C{2cm}C{2cm}L{4cm}@{}}
\caption{Project Risk Register}
\label{tab:risk-register} \\
\toprule
\textbf{ID} & \textbf{Risk} & \textbf{Probability} & \textbf{Impact} & \textbf{Mitigation} \\
\midrule
\endfirsthead
\toprule
\textbf{ID} & \textbf{Risk} & \textbf{Probability} & \textbf{Impact} & \textbf{Mitigation} \\
\midrule
\endhead
R1 & \gapref{G-001} not resolved by Q2 & Medium & High & Escalate to project sponsor; proceed with partial approval \\
R2 & New source documents added & Low & Medium & Change management process; incremental updates \\
R3 & Schema changes break traceability & Low & High & Automated validation scripts; versioned schemas \\
R4 & Reviewer availability delays & Medium & Medium & Identify backup reviewers; parallel review streams \\
\bottomrule
\end{longtable}

% =============================================================================
\section{Success Criteria}
\label{sec:success-criteria}

The HITL Requirements Traceability framework will be considered
successful when:

\begin{enumerate}
  \item All 6 HITL checkpoints (H1--H6) receive ``Approved'' or
        ``Conditionally Approved'' status
  \item 100\% of business requirements (excluding documented gaps)
        have verified traceability to source documents
  \item 100\% of specifications have verified traceability to requirements
  \item All identified gaps have documented resolution plans with
        assigned owners
  \item Sign-off declaration is completed by designated approver
  \item Automated validation scripts pass for all traceability links
\end{enumerate}

% =============================================================================
\section{Future Enhancements}
\label{sec:future-enhancements}

\subsection{v2.0 Planned Enhancements}

\begin{itemize}[noitemsep]
  \item \textbf{Automated Traceability Validation:} CI/CD pipeline to
        validate all links on every document change
  \item \textbf{Interactive Dashboard:} Web-based visualization of
        traceability matrix with drill-down capability
  \item \textbf{Change Impact Analysis:} Automated identification of
        downstream impacts when source documents change
  \item \textbf{ML-Assisted Extraction:} Use LLM to improve requirement
        extraction accuracy from unstructured documents
\end{itemize}

\subsection{Cross-Domain Expansion}

The HITL framework established for TB Review can be extended to:

\begin{itemize}[noitemsep]
  \item Arabic AP processing (parallel development)
  \item Regulations compliance tracking
  \item Simula training data validation
  \item General enterprise document processing
\end{itemize}
\chapter{Implementation Instructions}
\label{ch:implementation-instructions}

This chapter provides practical guidance for implementing and maintaining
the HITL Requirements Traceability framework, including tooling setup,
validation procedures, and maintenance workflows.

% =============================================================================
\section{Document Compilation}
\label{sec:document-compilation}

\subsection{Prerequisites}

\begin{itemize}[noitemsep]
  \item \textbf{TeX Distribution:} TeX Live 2024+ or MiKTeX 24+
  \item \textbf{Engine:} XeLaTeX (required for fontspec)
  \item \textbf{Packages:} Standard LaTeX packages (included in full distributions)
\end{itemize}

\subsection{Compilation Command}

\begin{lstlisting}[language=bash,basicstyle=\ttfamily\small]
cd docs/latex/specs/tb-hitl
xelatex -interaction=nonstopmode tb-hitl-traceability-spec.tex
xelatex -interaction=nonstopmode tb-hitl-traceability-spec.tex
\end{lstlisting}

\paragraph{Note:} Two passes are required for cross-references and
table of contents generation.

\subsection{Build Integration}

Add to the project Makefile:

\begin{lstlisting}[language=make,basicstyle=\ttfamily\small]
tb-hitl-spec:
	cd docs/latex/specs/tb-hitl && \
	xelatex -interaction=nonstopmode tb-hitl-traceability-spec.tex && \
	xelatex -interaction=nonstopmode tb-hitl-traceability-spec.tex

tb-hitl-pdf: tb-hitl-spec
	cp docs/latex/specs/tb-hitl/tb-hitl-traceability-spec.pdf \
	   docs/pdf/specs/tb-hitl-traceability-spec.pdf
\end{lstlisting}

% =============================================================================
\section{Traceability Link Validation}
\label{sec:link-validation}

\subsection{Manual Validation}

For each traceability link, verify:

\begin{enumerate}
  \item \textbf{Source exists:} The referenced source document is present
        at the specified path
  \item \textbf{Target exists:} The link target (requirement, spec, schema)
        is defined in its respective chapter
  \item \textbf{Bidirectional:} The reverse link is also documented
  \item \textbf{Consistent:} Link metadata matches between forward and
        backward traces
\end{enumerate}

\subsection{Automated Validation Script}

Create a Python validation script at \texttt{scripts/validate-hitl-links.py}:

\begin{lstlisting}[language=python,basicstyle=\ttfamily\small]
#!/usr/bin/env python3
"""Validate HITL traceability links."""

import yaml
import json
from pathlib import Path

def load_manifest():
    """Load the HITL traceability manifest."""
    manifest_path = Path("docs/tb/machine-readable/hitl-traceability-manifest.yaml")
    with open(manifest_path) as f:
        return yaml.safe_load(f)

def validate_source_exists(manifest):
    """Check all source documents exist."""
    errors = []
    for source in manifest.get("source_documents", []):
        path = Path(source["filepath"])
        if not path.exists():
            errors.append(f"Source not found: {source['source_ref']}")
    return errors

def validate_forward_links(manifest):
    """Validate forward traceability links."""
    errors = []
    requirements = {r["requirement_id"] for r in manifest.get("requirements", [])}
    
    for link in manifest.get("forward_traceability", []):
        if link["requirement_id"] not in requirements:
            errors.append(f"Unknown requirement: {link['requirement_id']}")
    return errors

if __name__ == "__main__":
    manifest = load_manifest()
    
    errors = []
    errors.extend(validate_source_exists(manifest))
    errors.extend(validate_forward_links(manifest))
    
    if errors:
        print("Validation FAILED:")
        for e in errors:
            print(f"  - {e}")
        exit(1)
    else:
        print("Validation PASSED")
        exit(0)
\end{lstlisting}

% =============================================================================
\section{Adding New Source Documents}
\label{sec:adding-sources}

When a new source document is added to the TB corpus:

\begin{enumerate}
  \item \textbf{Catalogue the document:}
        \begin{itemize}[noitemsep]
          \item Assign a source reference ID (e.g., \srcref{TB-NEW-001})
          \item Add entry to \cref{ch:source-documents}
          \item Update \texttt{hitl-traceability-manifest.yaml}
        \end{itemize}
        
  \item \textbf{Extract to machine-readable format:}
        \begin{itemize}[noitemsep]
          \item Run appropriate extraction script
          \item Store output in \texttt{docs/tb/machine-readable/extracted/}
          \item Add YAML frontmatter with source reference
        \end{itemize}
        
  \item \textbf{Derive requirements:}
        \begin{itemize}[noitemsep]
          \item Review extracted content for requirement statements
          \item Assign requirement IDs following naming convention
          \item Document derivation evidence in \cref{ch:business-requirements}
        \end{itemize}
        
  \item \textbf{Update traceability matrices:}
        \begin{itemize}[noitemsep]
          \item Add forward links in \cref{ch:bidirectional-traceability}
          \item Add backward links from affected specifications
          \item Run validation script to verify links
        \end{itemize}
        
  \item \textbf{Update gap analysis:}
        \begin{itemize}[noitemsep]
          \item Check for new gaps introduced
          \item Update orphan analysis tables
          \item Document any new gaps in \cref{ch:gap-analysis}
        \end{itemize}
\end{enumerate}

% =============================================================================
\section{Adding New Requirements}
\label{sec:adding-requirements}

\subsection{Requirement ID Convention}

\begin{table}[h]
\centering
\caption{Requirement ID Naming Convention}
\begin{tabularx}{\textwidth}{@{}L{3.5cm}L{4cm}L{6cm}@{}}
\toprule
\textbf{Pattern} & \textbf{Category} & \textbf{Example} \\
\midrule
TB-REQ-ES\textit{nn} & Executive Summary & TB-REQ-ES01 \\
TB-REQ-SC\textit{nn} & Scope & TB-REQ-SC01 \\
TB-REQ-PS\textit{nn} & Process Steps & TB-REQ-PS01 \\
TB-REQ-AI\textit{nn} & AI Augmentation & TB-REQ-AI01 \\
BSS-REQ-\textit{nn} & BSS Risk Assessment & BSS-REQ-01 \\
TB-REQ-FI\textit{nn} & Financial (future) & TB-REQ-FI01 \\
\bottomrule
\end{tabularx}
\end{table}

\subsection{Requirement Documentation Template}

\begin{lstlisting}[language=tex,basicstyle=\ttfamily\small]
\subsection{\reqref{TB-REQ-XX01}: Requirement Title}
\label{sec:req-xx01}

\begin{table}[h]
\centering
\caption{\reqref{TB-REQ-XX01} Derivation Evidence}
\begin{tabularx}{\textwidth}{@{}L{3.5cm}L{10.5cm}@{}}
\toprule
\textbf{Attribute} & \textbf{Value} \\
\midrule
Requirement ID & TB-REQ-XX01 \\
Requirement Text & The system shall... \\
Source Document & \srcref{TB-XXX-001}: Document name \\
Source Location & Section X, Paragraph Y \\
Derivation Method & Direct extraction / Synthesis \\
Confidence & High / Medium / Low \\
Status & Draft / Approved \\
\bottomrule
\end{tabularx}
\end{table}

\paragraph{Source Quote:}
\begin{quote}
\textit{``Quoted text from source document...''}
\end{quote}
\end{lstlisting}

% =============================================================================
\section{Schema Evolution}
\label{sec:schema-evolution}

When schemas are modified:

\begin{enumerate}
  \item \textbf{Document the change:}
        \begin{itemize}[noitemsep]
          \item Update schema file with new/modified fields
          \item Increment schema version
          \item Add changelog entry
        \end{itemize}
        
  \item \textbf{Update traceability:}
        \begin{itemize}[noitemsep]
          \item Add new fields to \cref{ch:schema-mapping}
          \item Link fields to specifications and requirements
          \item Verify no orphan fields remain
        \end{itemize}
        
  \item \textbf{Assess impact:}
        \begin{itemize}[noitemsep]
          \item Identify affected downstream consumers
          \item Document migration requirements
          \item Update validation scripts
        \end{itemize}
\end{enumerate}

% =============================================================================
\section{HITL Review Execution}
\label{sec:review-execution}

\subsection{Pre-Review Checklist}

Before initiating HITL review:

\begin{itemize}
  \item[$\square$] Document compiled successfully with no LaTeX errors
  \item[$\square$] All chapter files present and included
  \item[$\square$] Validation script passes
  \item[$\square$] Machine-readable manifest is current
  \item[$\square$] Reviewers identified and available
\end{itemize}

\subsection{Review Meeting Agenda}

\begin{enumerate}
  \item \textbf{Introduction (5 min):} Review objectives and scope
  \item \textbf{Source Review (15 min):} Checkpoint H1 --- verify source completeness
  \item \textbf{Extraction Review (15 min):} Checkpoint H2 --- sample extraction accuracy
  \item \textbf{Requirements Review (20 min):} Checkpoint H3 --- verify derivations
  \item \textbf{Traceability Review (15 min):} Checkpoint H4 --- validate links
  \item \textbf{Schema Review (15 min):} Checkpoint H5 --- verify alignment
  \item \textbf{Gap Review (10 min):} Checkpoint H6 --- accept gap plans
  \item \textbf{Sign-Off (5 min):} Complete sign-off register
\end{enumerate}

\subsection{Post-Review Actions}

\begin{enumerate}
  \item Document review findings in sign-off register
  \item Create issues for any identified defects
  \item Update gap register if new gaps identified
  \item Publish approved version to document repository
  \item Notify stakeholders of approval status
\end{enumerate}

% =============================================================================
\section{Continuous Integration}
\label{sec:ci-integration}

\subsection{CI Pipeline Configuration}

Add to \texttt{.github/workflows/docs.yml}:

\begin{lstlisting}[language=yaml,basicstyle=\ttfamily\small]
name: HITL Documentation

on:
  push:
    paths:
      - 'docs/latex/specs/tb-hitl/**'
      - 'docs/tb/machine-readable/**'

jobs:
  build-and-validate:
    runs-on: ubuntu-latest
    steps:
      - uses: actions/checkout@v4
      
      - name: Install TeX Live
        run: |
          sudo apt-get update
          sudo apt-get install -y texlive-xetex texlive-latex-extra
      
      - name: Validate traceability links
        run: python scripts/validate-hitl-links.py
      
      - name: Compile specification
        run: |
          cd docs/latex/specs/tb-hitl
          xelatex -interaction=nonstopmode tb-hitl-traceability-spec.tex
          xelatex -interaction=nonstopmode tb-hitl-traceability-spec.tex
      
      - name: Upload PDF artifact
        uses: actions/upload-artifact@v4
        with:
          name: tb-hitl-spec
          path: docs/latex/specs/tb-hitl/tb-hitl-traceability-spec.pdf
\end{lstlisting}

% =============================================================================
\section{Troubleshooting}
\label{sec:troubleshooting}

\begin{table}[h]
\centering
\caption{Common Issues and Solutions}
\begin{tabularx}{\textwidth}{@{}L{5cm}L{8.5cm}@{}}
\toprule
\textbf{Issue} & \textbf{Solution} \\
\midrule
LaTeX compilation error & Check for missing packages; verify fontspec compatibility \\
Undefined cross-reference & Run xelatex twice; check label spelling \\
Missing source document & Verify file path in manifest; check for renamed files \\
Broken traceability link & Run validation script; check ID spelling \\
Table overflow & Reduce column widths; use longtable for long tables \\
TikZ positioning error & Verify tikzlibrary imports; check node distance \\
\bottomrule
\end{tabularx}
\end{table}
\chapter*{Quick Reference Card}
\addcontentsline{toc}{chapter}{Quick Reference Card}
\label{ch:quick-ref}

\begin{tcolorbox}[colback=white,colframe=sapblue,title=\textbf{TB-HITL Quick Reference Card --- Print \& Keep}]

\section*{Source Document IDs}
\begin{tabular}{@{}ll@{}}
\srcref{TB-BC-001} & Business Case - Trial Balance \\
\srcref{TB-BC-002} & Business Case - BSS Risk Assessment \\
\srcref{TB-DOI-001} & DOI for Trial Balance Process \\
\srcref{TB-BPMN-001} & BPMN Process Map (GLO) \\
\srcref{TB-WB-001} & HKG TB Review Workbook \\
\srcref{TB-WB-002} & HKG PL Review Workbook \\
\end{tabular}

\vspace{0.3cm}
\hrule
\vspace{0.3cm}

\section*{Requirement IDs (27 Total)}

\begin{minipage}[t]{0.48\textwidth}
\textbf{Executive Summary (4)}
\begin{itemize}[noitemsep,leftmargin=*]
  \item \reqref{TB-REQ-ES01} AI variance analysis
  \item \reqref{TB-REQ-ES02} Commentary generation
  \item \reqref{TB-REQ-ES03} 5 FTE savings
  \item \reqref{TB-REQ-ES04} Studio AI solution
\end{itemize}

\textbf{Scope (6)}
\begin{itemize}[noitemsep,leftmargin=*]
  \item \reqref{TB-REQ-SC01} 234 Legal Entities
  \item \reqref{TB-REQ-SC02} Finance caption level
  \item \reqref{TB-REQ-SC03} Segment level
  \item \reqref{TB-REQ-SC04} Legal entity level
  \item \reqref{TB-REQ-SC05} Exclude Country IFRS
  \item \reqref{TB-REQ-SC06} Exclude BSS RA
\end{itemize}

\textbf{AI Augmentation (6)}
\begin{itemize}[noitemsep,leftmargin=*]
  \item \reqref{TB-REQ-AI01} Anomaly detection
  \item \reqref{TB-REQ-AI02} Trend analysis
  \item \reqref{TB-REQ-AI03} New account detection
  \item \reqref{TB-REQ-AI04} Spike/drop detection
  \item \reqref{TB-REQ-AI05} Material balance ID
  \item \reqref{TB-REQ-AI06} LLM commentaries
\end{itemize}
\end{minipage}%
\hfill
\begin{minipage}[t]{0.48\textwidth}
\textbf{Process Steps (8)}
\begin{itemize}[noitemsep,leftmargin=*]
  \item \reqref{TB-REQ-PS01} Roll forward file
  \item \reqref{TB-REQ-PS02} Download PSGL
  \item \reqref{TB-REQ-PS03} Extract TB
  \item \reqref{TB-REQ-PS04} Calculate variance
  \item \reqref{TB-REQ-PS05} Filter materiality
  \item \reqref{TB-REQ-PS06} Seek commentaries
  \item \reqref{TB-REQ-PS07} Track entries/risk
  \item \reqref{TB-REQ-PS08} Prepare summary
\end{itemize}

\textbf{BSS Risk Assessment (3)}
\begin{itemize}[noitemsep,leftmargin=*]
  \item \reqref{BSS-REQ-01} Risk categorization
  \item \reqref{BSS-REQ-02} Post-FCA consolidation
  \item \reqref{BSS-REQ-03} Risk reporting packs
\end{itemize}
\end{minipage}

\vspace{0.3cm}
\hrule
\vspace{0.3cm}

\section*{Specification IDs (22 Total)}

\begin{minipage}[t]{0.32\textwidth}
\textbf{Process Steps (8)}
\begin{itemize}[noitemsep,leftmargin=*,font=\small]
  \item \specref{TB-STEP-001}
  \item \specref{TB-STEP-002}
  \item \specref{TB-STEP-003}
  \item \specref{TB-STEP-004}
  \item \specref{TB-STEP-005}
  \item \specref{TB-STEP-006}
  \item \specref{TB-STEP-007}
  \item \specref{TB-STEP-008}
\end{itemize}
\end{minipage}%
\begin{minipage}[t]{0.32\textwidth}
\textbf{Decision Points (5)}
\begin{itemize}[noitemsep,leftmargin=*,font=\small]
  \item \specref{DP-001} Materiality
  \item \specref{DP-002} Journal Status
  \item \specref{DP-003} Investigation
  \item \specref{DP-004} Confirmation
  \item \specref{DP-005} Unexplained
\end{itemize}

\textbf{Controls (3)}
\begin{itemize}[noitemsep,leftmargin=*,font=\small]
  \item \specref{TB-CTRL-001}
  \item \specref{TB-CTRL-002}
  \item \specref{TB-CTRL-003}
\end{itemize}
\end{minipage}%
\begin{minipage}[t]{0.32\textwidth}
\textbf{AI Capabilities (6)}
\begin{itemize}[noitemsep,leftmargin=*,font=\small]
  \item \specref{AI-ANOM}
  \item \specref{AI-TREND}
  \item \specref{AI-NEW}
  \item \specref{AI-SPIKE}
  \item \specref{AI-MAT}
  \item \specref{AI-LLM}
\end{itemize}
\end{minipage}

\vspace{0.3cm}
\hrule
\vspace{0.3cm}

\section*{Schema IDs (6 Total)}
\begin{tabular}{@{}lll@{}}
\schemaref{variance-record} & \schemaref{commentary-draft} & \schemaref{tb-extract} \\
\schemaref{entity-params} & \schemaref{decision-point} & \schemaref{pl-extract} \\
\end{tabular}

\vspace{0.3cm}
\hrule
\vspace{0.3cm}

\section*{HITL Checkpoints}
\begin{tabular}{@{}llll@{}}
\textbf{H1} Source Completeness & \textbf{H2} Extraction Accuracy & \textbf{H3} Requirement Derivation \\
\textbf{H4} Traceability Links & \textbf{H5} Schema Alignment & \textbf{H6} Gap Analysis \\
\end{tabular}

\vspace{0.3cm}
\hrule
\vspace{0.3cm}

\section*{Gap IDs}
\begin{tabular}{@{}lll@{}}
\gapref{G-001} & Financial Analysis TBD & Owner: S1 Badrinath \\
\gapref{G-002} & DOI images not OCR'd & Owner: Technical Lead \\
\end{tabular}

\end{tcolorbox}
% =============================================================================
% Chapter 14: Audit Findings and Resolutions
% =============================================================================
\chapter{Audit Findings and Resolutions}
\label{ch:audit-findings}

This chapter documents audit findings from the formal review and their 
resolutions. All P0 (critical) and P1 (high priority) findings are 
addressed before production release.

% =============================================================================
\section{Audit Review Summary}
\label{sec:audit-summary}
% =============================================================================

\begin{tcolorbox}[colback=yellow!10,colframe=yellow!50!black,title=Audit Status: Conditionally Approved]
\textbf{Review Date:} 20 April 2026\\
\textbf{Overall Verdict:} Conditionally Approved\\
\textbf{Condition:} P0 items must be resolved before GA cutover

\vspace{0.5em}
\begin{tabular}{lccl}
\toprule
\textbf{Checkpoint} & \textbf{Status} & \textbf{Condition} \\
\midrule
H1 - Source Completeness & \textcolor{green}{\textbf{PASS}} & None \\
H2 - Extraction Accuracy & \textcolor{green}{\textbf{PASS}} & None \\
H3 - Requirement Derivation & \textcolor{orange}{\textbf{CONDITIONAL}} & Pending G-001 resolution \\
H4 - Traceability Links & \textcolor{green}{\textbf{PASS}} & None \\
H5 - Schema Alignment & \textcolor{orange}{\textbf{CONDITIONAL}} & Pending audit trail schema \\
H6 - Gap Analysis & \textcolor{orange}{\textbf{CONDITIONAL}} & Pending G-003 cross-spec \\
\bottomrule
\end{tabular}
\end{tcolorbox}

% =============================================================================
\section{P0 Findings: Critical (Resolution Required Before v1.0)}
\label{sec:p0-findings}
% =============================================================================

% -----------------------------------------------------------------------------
\subsection{P0-1: State Count Reconciliation (27 vs 28)}
\label{sec:p0-state-count}
% -----------------------------------------------------------------------------

\begin{tcolorbox}[colback=red!5,colframe=red!70,title=\faIcon{exclamation-triangle} Finding P0-1: State Count Inconsistency]
\textbf{Location:} HITL Spec p.9 vs Implementation Spec Table 5.1\\
\textbf{Issue:} HITL spec references ``27 states + 3 error-handling'' while 
implementation spec Table 5.1 enumerates 28 states\\
\textbf{Severity:} \textbf{P0 - Critical}\\
\textbf{Owner:} Spec Owner\\
\textbf{Timeline:} 1 week
\end{tcolorbox}

\textbf{Resolution:}

\begin{tcolorbox}[colback=green!5,colframe=green!70,title=\faIcon{check-circle} Resolution Applied]
The \textbf{implementation spec's 28-state enumeration is authoritative}.

The discrepancy arose from counting methodology:
\begin{itemize}
    \item \textbf{HITL Spec (27+3):} Counted 27 ``primary'' states plus 3 error-handling 
          states as a separate category = 30 total states
    \item \textbf{Implementation Spec (28):} Counted 28 operational states; error-handling 
          states (AWAITING\_CLARIFICATION, ESCALATED\_TO\_FINANCE, MANUAL\_OVERRIDE) 
          are subsumed within the 28
\end{itemize}

\textbf{Authoritative Source:} \texttt{structured/workflows/tb-review-glo.json}\\
\textbf{Correct Count:} \textbf{28 states} as enumerated in Table 5.1

\vspace{0.5em}
\textbf{Corrective Action:} This specification is updated to reference 28 states 
throughout. Cross-references corrected in:
\begin{itemize}
    \item Chapter~\ref{ch:source-documents}: Updated BPMN extraction to state ``28 states''
    \item Chapter~\ref{ch:specification-mapping}: Updated process step count
    \item YAML manifest: Updated \texttt{statistics.bpmn\_states} to 28
\end{itemize}
\end{tcolorbox}

\textbf{State Enumeration (Authoritative):}

\begin{longtable}{clll}
\caption{Authoritative 28-State Enumeration} \\
\toprule
\textbf{\#} & \textbf{State ID} & \textbf{Type} & \textbf{Timing} \\
\midrule
\endfirsthead
\toprule
\textbf{\#} & \textbf{State ID} & \textbf{Type} & \textbf{Timing} \\
\midrule
\endhead
1 & ROLL\_FORWARD & Task & M-3 \\
2 & DOWNLOAD\_PSGL & Task & M-3 \\
3 & UPDATE\_ACCOUNTS\_MASTER & Task & M-3 \\
4 & EXTRACT\_TB & Task & M-3 to M+5 \\
5 & UPDATE\_COMPARATIVES & Task & M-3 to M+5 \\
6 & CALC\_VARIANCE & Task & M-2 to M+5 \\
7 & FILTER\_VARIANCES & Task & M-2 to M+5 \\
8 & UPDATE\_COMMENTARIES & Task & M-2 to M+5 \\
9 & CHECK\_JOURNAL\_POSTED & Gateway & M-2 to M+5 \\
10 & SEND\_JOURNAL\_REQUEST & Task & M-2 to M+5 \\
11 & UPDATE\_RISK\_SUMMARY & Task & M-2 to M+5 \\
12 & PREPARE\_SUMMARY & Task & M-2 to M+5 \\
13 & INVESTIGATE\_VARIANCE & Task & M-2 to M+5 \\
14 & CHECK\_MATERIAL & Gateway & M-2 to M+5 \\
15 & INTERNAL\_REVIEW & Task & M+5 \\
16 & SEND\_FORTM\_SUMMARY & Task & M+5 \\
17 & CATEGORIZE\_UNEXPLAINED & Task & M-2 to M+5 \\
18 & CHECK\_TEAM\_ENTRIES & Task & M-2 to M+5 \\
19 & CONFIRM\_ENTRIES & Task & M-2 to M+5 \\
20 & SEEK\_COMMENTARIES & Task & M-2 to M+5 \\
21 & RECEIVE\_COMMENTARIES & Task & M-2 to M+5 \\
22 & PERFORM\_RISK\_ANALYSIS & Task & M-2 to M+5 \\
23 & UPDATE\_VARIANCE\_COMMENTARIES & Task & M-2 to M+5 \\
24 & TRACK\_ENTRIES & Task & M-2 to M+5 \\
25 & AWAITING\_CLARIFICATION & Error & As needed \\
26 & ESCALATED\_TO\_FINANCE & Error & As needed \\
27 & MANUAL\_OVERRIDE & Error & As needed \\
28 & COMPLETED & End & M+5 \\
\bottomrule
\end{longtable}

% -----------------------------------------------------------------------------
\subsection{P0-2: vLLM Fallback Specification}
\label{sec:p0-vllm-fallback}
% -----------------------------------------------------------------------------

\begin{tcolorbox}[colback=red!5,colframe=red!70,title=\faIcon{exclamation-triangle} Finding P0-2: vLLM Fallback Not Specified]
\textbf{Location:} tb-review-spec, p.13, \S3.1\\
\textbf{Issue:} ``vllm-only routing policy'' with no fallback if vLLM cluster degrades\\
\textbf{Impact:} Commentary generation (Stage 3) blocks on LLM availability\\
\textbf{Severity:} \textbf{P0 - Critical}\\
\textbf{Owner:} AI Architect\\
\textbf{Timeline:} 2 weeks
\end{tcolorbox}

\textbf{Resolution:}

\begin{tcolorbox}[colback=green!5,colframe=green!70,title=\faIcon{check-circle} Resolution: Graceful Degradation Specification]
\textbf{Primary:} vLLM cluster via AI Core (default)\\
\textbf{Fallback:} Template-based commentary generation\\
\textbf{Trigger Conditions:} (any of)
\begin{itemize}
    \item vLLM latency P95 $> 5000$ms (5-minute rolling window)
    \item vLLM error rate $> 10\%$ (5-minute rolling window)
    \item vLLM cluster health check fails 3 consecutive probes
\end{itemize}
\end{tcolorbox}

\textbf{Fallback Implementation Specification:}

\begin{lstlisting}[language=Python,caption={vLLM Fallback Logic}]
class CommentaryService:
    VLLM_LATENCY_THRESHOLD_MS = 5000
    VLLM_ERROR_RATE_THRESHOLD = 0.10
    HEALTH_CHECK_FAILURES_MAX = 3
    
    def generate_commentary(self, variance_record: VarianceRecord) -> Commentary:
        if self._should_use_fallback():
            return self._generate_template_commentary(variance_record)
        try:
            return self._generate_vllm_commentary(variance_record)
        except VLLMTimeoutError:
            self._record_latency_breach()
            return self._generate_template_commentary(variance_record)
    
    def _generate_template_commentary(self, vr: VarianceRecord) -> Commentary:
        """Template-based fallback with 80% of LLM quality."""
        template = self._select_template(vr.variance_type, vr.account_category)
        return Commentary(
            text=template.format(
                account=vr.account_name,
                variance_amount=format_currency(vr.variance_amount),
                variance_pct=format_percent(vr.variance_percentage),
                direction="increased" if vr.variance_amount > 0 else "decreased",
                prior_period=vr.prior_period,
            ),
            confidence=0.75,  # Template confidence capped at 75%
            source="template_fallback",
            llm_available=False,
        )
\end{lstlisting}

\textbf{Template Library Requirements:}

\begin{tabular}{lll}
\toprule
\textbf{Template ID} & \textbf{Variance Type} & \textbf{Account Category} \\
\midrule
TPL-001 & Spike ($>25\%$) & Revenue \\
TPL-002 & Spike ($>25\%$) & Cost of Sales \\
TPL-003 & Drop ($<-25\%$) & Revenue \\
TPL-004 & Drop ($<-25\%$) & Cost of Sales \\
TPL-005 & New Account & Any \\
TPL-006 & Material Balance & Any \\
TPL-007 & General Variance & Any (catch-all) \\
\bottomrule
\end{tabular}

\textbf{Monitoring Requirements:}

\begin{enumerate}
    \item Dashboard SHALL display ``LLM Status'' indicator (green/amber/red)
    \item Fallback activations SHALL be logged with timestamp and trigger reason
    \item Weekly report SHALL include fallback activation count and duration
    \item SLA: Fallback mode SHALL not exceed 4 hours continuous without escalation
\end{enumerate}

% =============================================================================
\section{P1 Findings: High Priority (Resolution Required Within 4 Weeks)}
\label{sec:p1-findings}
% =============================================================================

% -----------------------------------------------------------------------------
\subsection{P1-1: Audit Trail Schema for EUC Retirement}
\label{sec:p1-audit-trail}
% -----------------------------------------------------------------------------

\begin{tcolorbox}[colback=orange!5,colframe=orange!70,title=\faIcon{exclamation-circle} Finding P1-1: Audit Trail Schema Not Defined]
\textbf{Location:} tb-review-spec, Section 1.8 (p.6-7)\\
\textbf{Issue:} Dashboard audit trail defined as ``authoritative evidence'' but no 
field-level schema specified\\
\textbf{Impact:} Cannot verify EUC retirement evidence standard\\
\textbf{Severity:} \textbf{P1 - High}\\
\textbf{Owner:} Controls Owner\\
\textbf{Timeline:} 4 weeks
\end{tcolorbox}

\textbf{Resolution:}

\begin{tcolorbox}[colback=green!5,colframe=green!70,title=\faIcon{check-circle} Schema Definition: audit-trail-entry.schema.json]
\begin{lstlisting}[language=json,basicstyle=\tiny\ttfamily]
{
  "$schema": "https://json-schema.org/draft/2020-12/schema",
  "$id": "audit-trail-entry.schema.json",
  "title": "Audit Trail Entry",
  "type": "object",
  "required": ["audit_id", "timestamp", "persona_id", "action", "evidence_hash"],
  "properties": {
    "audit_id": {
      "type": "string",
      "format": "uuid",
      "description": "Unique identifier for this audit entry"
    },
    "timestamp": {
      "type": "string",
      "format": "date-time",
      "description": "ISO 8601 timestamp of the action"
    },
    "persona_id": {
      "type": "string",
      "description": "Employee ID from SAP IAS token (sub claim)"
    },
    "action": {
      "type": "string",
      "enum": ["view", "claim", "decide", "escalate", "override", "bulk_assign", 
               "comment_edit", "comment_approve", "comment_reject", "export"]
    },
    "resource_type": {
      "type": "string",
      "enum": ["variance_record", "commentary", "workflow_state", "entity_params"]
    },
    "resource_id": {
      "type": "string",
      "description": "ID of the affected resource"
    },
    "before_state": {
      "type": "object",
      "description": "Snapshot of resource state before action (null for create)"
    },
    "after_state": {
      "type": "object",
      "description": "Snapshot of resource state after action"
    },
    "evidence_hash": {
      "type": "string",
      "pattern": "^sha256:[a-f0-9]{64}$",
      "description": "SHA-256 hash of (timestamp + persona_id + action + resource_id + after_state)"
    },
    "bulk_operation_id": {
      "type": "string",
      "format": "uuid",
      "description": "Groups entries from same bulk operation (null for single ops)"
    },
    "ip_address": {
      "type": "string",
      "format": "ipv4",
      "description": "Client IP address"
    },
    "user_agent": {
      "type": "string",
      "description": "Browser/client user agent string"
    }
  }
}
\end{lstlisting}
\end{tcolorbox}

\textbf{Divergence Detection Specification:}

During the 12-month EUC deprecation window, divergence between Excel EUC and 
dashboard may occur. The following automated detection is specified:

\begin{lstlisting}[language=SQL,caption={Daily Divergence Detection Query}]
-- Run daily at M+6 for prior month data
WITH euc_data AS (
    SELECT legal_entity, account, closing_balance 
    FROM imported_euc_tb WHERE period = :period
),
dashboard_data AS (
    SELECT legal_entity, account, closing_balance 
    FROM variance_records WHERE period = :period
)
SELECT 
    COALESCE(e.legal_entity, d.legal_entity) AS le,
    COALESCE(e.account, d.account) AS account,
    e.closing_balance AS euc_balance,
    d.closing_balance AS dashboard_balance,
    ABS(e.closing_balance - d.closing_balance) AS divergence,
    CASE 
        WHEN e.closing_balance IS NULL THEN 'dashboard_only'
        WHEN d.closing_balance IS NULL THEN 'euc_only'
        ELSE 'value_mismatch'
    END AS divergence_type
FROM euc_data e
FULL OUTER JOIN dashboard_data d 
    ON e.legal_entity = d.legal_entity AND e.account = d.account
WHERE e.closing_balance != d.closing_balance 
   OR e.closing_balance IS NULL OR d.closing_balance IS NULL;
\end{lstlisting}

\textbf{Resolution Rule:} Dashboard is authoritative. Divergence records are 
logged with \texttt{action="divergence\_detected"} and require manual review 
by Controls Owner within 5 business days.

% -----------------------------------------------------------------------------
\subsection{P1-2: Statistical Normality Test for 3$\sigma$ Threshold}
\label{sec:p1-normality}
% -----------------------------------------------------------------------------

\begin{tcolorbox}[colback=orange!5,colframe=orange!70,title=\faIcon{exclamation-circle} Finding P1-2: Statistical Assumption Not Validated]
\textbf{Location:} tb-review-spec, p.14, \S3.3\\
\textbf{Issue:} 3$\sigma$ threshold assumes normally-distributed variance; 
financial data often exhibits fat tails (leptokurtosis)\\
\textbf{Severity:} \textbf{P1 - High}\\
\textbf{Owner:} Data Science\\
\textbf{Timeline:} 4 weeks
\end{tcolorbox}

\textbf{Resolution:}

\begin{tcolorbox}[colback=green!5,colframe=green!70,title=\faIcon{check-circle} Resolution: Normality Testing with MAD Fallback]
\textbf{Acceptance Criteria Addition (AC-ES01-2a):}

Before applying Z-score $> 3.0$ threshold, the system SHALL:
\begin{enumerate}
    \item Run Shapiro-Wilk normality test per account category
    \item If $p$-value $> 0.05$: Use Z-score (normal distribution confirmed)
    \item If $p$-value $\leq 0.05$: Use MAD (Median Absolute Deviation) with 
          threshold $= 3.5 \times$ MAD (robust to fat tails)
\end{enumerate}
\end{tcolorbox}

\textbf{Implementation:}

\begin{lstlisting}[language=Python,caption={Normality Test with MAD Fallback}]
from scipy import stats
import numpy as np

class AnomalyDetector:
    ZSCORE_THRESHOLD = 3.0
    MAD_THRESHOLD = 3.5
    SHAPIRO_P_THRESHOLD = 0.05
    
    def detect_anomalies(self, values: np.ndarray) -> np.ndarray:
        """Returns boolean mask of anomalies."""
        # Test for normality
        if len(values) >= 20:  # Shapiro-Wilk requires n >= 20
            _, p_value = stats.shapiro(values)
            use_robust = p_value <= self.SHAPIRO_P_THRESHOLD
        else:
            use_robust = True  # Default to robust for small samples
        
        if use_robust:
            return self._mad_detection(values)
        else:
            return self._zscore_detection(values)
    
    def _zscore_detection(self, values: np.ndarray) -> np.ndarray:
        z_scores = np.abs(stats.zscore(values))
        return z_scores > self.ZSCORE_THRESHOLD
    
    def _mad_detection(self, values: np.ndarray) -> np.ndarray:
        median = np.median(values)
        mad = np.median(np.abs(values - median))
        # Consistency constant for normal distribution approximation
        mad_adjusted = mad * 1.4826
        modified_z = np.abs(values - median) / mad_adjusted
        return modified_z > self.MAD_THRESHOLD
\end{lstlisting}

\textbf{Documentation Note:} Add to Chapter~\ref{ch:business-requirements}, 
AC-ES01-2: ``Anomaly detection SHALL flag items with Z-score $> 3.0$ for 
normally-distributed data (Shapiro-Wilk $p > 0.05$) \textbf{OR} MAD-score 
$> 3.5$ for non-normal distributions.''

% =============================================================================
\section{Cross-Cutting Issue Resolutions}
\label{sec:cross-cutting}
% =============================================================================

% -----------------------------------------------------------------------------
\subsection{C-1: Terminology Reconciliation}
\label{sec:c1-terminology}
% -----------------------------------------------------------------------------

\begin{tabular}{lll}
\toprule
\textbf{Term} & \textbf{Authoritative Value} & \textbf{Corrective Action} \\
\midrule
State count & \textbf{28 states} & Update HITL spec to 28 \\
EUC-002 name & ``TB Variance Analysis file integrity'' & Standardize across docs \\
\bottomrule
\end{tabular}

% -----------------------------------------------------------------------------
\subsection{C-2: Materiality Threshold Standardization}
\label{sec:c2-materiality}
% -----------------------------------------------------------------------------

\begin{tcolorbox}[colback=blue!5,colframe=blue!70,title=Resolution: Entity-Params Configuration]
\textbf{Issue:} Hardcoded thresholds (BS $> \$100$M, PL $> \$3$M) in DP-001 
conflict with ``country-specific'' in TB-KPI-001.

\textbf{Resolution:} Move thresholds to \texttt{entity-params/<le-code>.yaml}:
\begin{lstlisting}[language=yaml]
# entity-params/HKG.yaml
legal_entity_code: "HKG"
materiality_thresholds:
  balance_sheet_absolute: 100000000  # $100M USD
  profit_loss_absolute: 3000000      # $3M USD
  variance_percentage: 0.10          # 10%
\end{lstlisting}

\textbf{Corrective Action:} Remove hardcoded thresholds from DP-001 description; 
reference ``entity-specific materiality from \texttt{entity-params}'' instead.
\end{tcolorbox}

% =============================================================================
\section{Additional Auditor Findings}
\label{sec:additional-findings}
% =============================================================================

% -----------------------------------------------------------------------------
\subsection{A-2: 100\% Traceability Claim Qualification}
\label{sec:a2-traceability-exceptions}
% -----------------------------------------------------------------------------

\begin{tcolorbox}[colback=orange!5,colframe=orange!70,title=\faIcon{exclamation-circle} Finding A-2: Exceptions to 100\% Traceability]
\textbf{Issue:} The spec claims ``100\% BIDIRECTIONAL TRACEABILITY ACHIEVED'' 
but exceptions exist and should be explicitly disclosed.
\end{tcolorbox}

\textbf{Resolution: Exceptions Table}

\begin{longtable}{p{3cm} p{3.5cm} p{6cm}}
\caption{Explicit Exceptions to 100\% Traceability} \\
\toprule
\textbf{Layer} & \textbf{Item} & \textbf{Justification} \\
\midrule
\endfirsthead
\toprule
\textbf{Layer} & \textbf{Item} & \textbf{Justification} \\
\midrule
\endhead
Spec $\rightarrow$ Schema & TB-STEP-006 (Journal Posting Investigation) & Workflow orchestration step; no persistent data schema required \\
Spec $\rightarrow$ Schema & TB-STEP-007 (Track Entries \& Risk) & Workflow orchestration step; state tracked in workflow engine, not schema \\
Schema $\rightarrow$ Spec & \texttt{audit\_id} field & Infrastructure field for tamper detection; not derived from business requirement \\
Schema $\rightarrow$ Spec & \texttt{version} field & Infrastructure field for concurrency control; not derived from business requirement \\
Schema $\rightarrow$ Spec & \texttt{processed\_at} field & Infrastructure field for pipeline telemetry; not derived from business requirement \\
\bottomrule
\end{longtable}

\textbf{Updated Claim:}
\begin{quote}
\textit{``100\% BIDIRECTIONAL TRACEABILITY ACHIEVED for business requirements. 
Five infrastructure/orchestration items are documented exceptions (see Table 14.X).''}
\end{quote}

% -----------------------------------------------------------------------------
\subsection{A-4: BSS Risk Assessment Vocabulary Alignment}
\label{sec:a4-bss-vocabulary}
% -----------------------------------------------------------------------------

\begin{tcolorbox}[colback=orange!5,colframe=orange!70,title=\faIcon{exclamation-circle} Finding A-4: BSS Cross-Spec Vocabulary]
\textbf{Issue:} BSS Risk Assessment is out of scope but risk vocabulary 
(\texttt{substantiated\_at\_risk}, \texttt{unsubstantiated\_unowned}, 
\texttt{unsubstantiated\_owned}) appears in DP-005.\\
\textbf{Question:} How will cross-spec vocabulary alignment be maintained?
\end{tcolorbox}

\textbf{Resolution: Shared Vocabulary Registry}

\begin{tcolorbox}[colback=green!5,colframe=green!70,title=\faIcon{check-circle} Resolution: Cross-Spec Enum Registry]
A shared vocabulary registry SHALL be maintained at:
\begin{verbatim}
docs/schema/common/enums.yaml
\end{verbatim}

Risk status values are defined once and imported by both TB and BSS specs:

\begin{lstlisting}[language=yaml,caption={Shared Enum Registry (enums.yaml)}]
# docs/schema/common/enums.yaml
risk_status:
  $id: "common/risk-status"
  description: "Shared risk classification vocabulary (TB + BSS)"
  type: string
  enum:
    - substantiated           # Evidence supports balance
    - substantiated_at_risk   # Evidence exists but control weakness identified
    - unsubstantiated_unowned # No evidence; no owner assigned
    - unsubstantiated_owned   # No evidence; owner assigned for remediation
  ownership:
    maintainer: "Controls Architecture Team"
    consumers: ["tb-review-spec", "bss-risk-assessment-spec"]
    change_process: "RFC required; both spec owners must approve"
\end{lstlisting}

\textbf{Governance:}
\begin{enumerate}
    \item Any change to shared enums requires RFC (Request for Comment)
    \item Both TB and BSS spec owners must approve vocabulary changes
    \item CI validation ensures specs reference registry (not local definitions)
\end{enumerate}
\end{tcolorbox}

% -----------------------------------------------------------------------------
\subsection{I-3: Bulk Operations Error Handling}
\label{sec:i3-bulk-operations}
% -----------------------------------------------------------------------------

\begin{tcolorbox}[colback=orange!5,colframe=orange!70,title=\faIcon{exclamation-circle} Finding I-3: Bulk Operations Gaps]
\textbf{Missing specifications:}
\begin{itemize}
    \item Partial success handling (e.g., 10 of 50 items fail)
    \item Race conditions (two reviewers claiming same item via bulk)
    \item Audit of bulk operations (single entry or per-item?)
\end{itemize}
\end{tcolorbox}

\textbf{Resolution: Bulk Operations Specification}

\begin{tcolorbox}[colback=green!5,colframe=green!70,title=\faIcon{check-circle} Resolution: Bulk Operation Semantics]

\textbf{1. Atomic Per-Item with Rollback:}
\begin{itemize}
    \item Each item in a bulk operation SHALL be processed atomically
    \item If \textbf{any} item fails, the \textbf{entire bulk operation} SHALL be rolled back
    \item No partial success state is allowed
\end{itemize}

\textbf{2. Race Condition Prevention:}
\begin{lstlisting}[language=SQL,caption={Optimistic Locking for Bulk Claim}]
-- Bulk claim with optimistic locking
BEGIN TRANSACTION;
  -- Attempt to claim all items atomically
  UPDATE variance_records 
  SET claimant_id = :persona_id,
      claimed_at = NOW(),
      version = version + 1
  WHERE id IN (:item_ids)
    AND claimant_id IS NULL  -- Only unclaimed items
    AND version = :expected_version;  -- Optimistic lock
  
  -- Check all items were claimed
  IF ROW_COUNT() != :expected_count THEN
    ROLLBACK;  -- Race condition detected
    RETURN error: 'CONFLICT: Some items already claimed'
  END IF;
COMMIT;
\end{lstlisting}

\textbf{3. Audit of Bulk Operations:}
\begin{itemize}
    \item Each item SHALL have its own audit entry
    \item All entries SHALL share the same \texttt{bulk\_operation\_id}
    \item A summary audit entry SHALL be created with operation context
\end{itemize}

\begin{lstlisting}[language=json,caption={Bulk Operation Audit Example}]
// Summary entry
{
  "audit_id": "uuid-summary",
  "action": "bulk_assign",
  "bulk_operation_id": "uuid-bulk-123",
  "bulk_operation_context": {
    "total_items": 50,
    "success_count": 50,
    "failure_count": 0,
    "operation_type": "bulk_assign"
  }
}

// Individual item entries (50 of these)
{
  "audit_id": "uuid-item-1",
  "action": "bulk_assign",
  "bulk_operation_id": "uuid-bulk-123",
  "resource_id": "variance-record-001",
  "change_summary": {
    "new_assignee": "1234567"
  }
}
\end{lstlisting}
\end{tcolorbox}

% -----------------------------------------------------------------------------
\subsection{I-4: Kill-Switch Automation Specification}
\label{sec:i4-kill-switch}
% -----------------------------------------------------------------------------

\begin{tcolorbox}[colback=orange!5,colframe=orange!70,title=\faIcon{exclamation-circle} Finding I-4: Kill-Switch Automation Missing]
\textbf{Table 11.4 defines six kill-switch conditions but doesn't specify:}
\begin{itemize}
    \item Who evaluates? (Automated monitor vs. manual check)
    \item Evaluation frequency (continuous? post-cycle?)
    \item Notification mechanism (email? PagerDuty? Jira?)
\end{itemize}
\end{tcolorbox}

\textbf{Resolution: Kill-Switch Automation Specification}

\begin{longtable}{p{2.5cm} p{2cm} p{2cm} p{3cm} p{3cm}}
\caption{Kill-Switch Condition Automation} \\
\toprule
\textbf{Condition} & \textbf{Evaluator} & \textbf{Frequency} & \textbf{Notification} & \textbf{Probe Implementation} \\
\midrule
\endfirsthead
\toprule
\textbf{Condition} & \textbf{Evaluator} & \textbf{Frequency} & \textbf{Notification} & \textbf{Probe Implementation} \\
\midrule
\endhead
Accept-without-edit $<$ 50\% & Automated & Post-cycle (M+6) & PagerDuty P2 + Slack & SQL query on decisions table \\
Material variance missed & Automated & Daily at 06:00 UTC & PagerDuty P1 + Email & Compare AI flags vs manual flags \\
Commentary factual error & Manual & Weekly review & Jira ticket & Random sampling (5\%) \\
Hallucination detected & Automated & Real-time & PagerDuty P1 + Kill & Confidence score $<$ 0.3 AND no source \\
vLLM availability $<$ 95\% & Automated & Continuous (1-min) & PagerDuty P2 & Health check endpoint \\
Audit trail integrity & Automated & Hourly & PagerDuty P1 + Kill & Evidence hash verification \\
\bottomrule
\end{longtable}

\textbf{Probe Implementation Examples:}

\begin{lstlisting}[language=Python,caption={Kill-Switch Probe Implementation}]
class KillSwitchMonitor:
    """Automated kill-switch condition evaluation."""
    
    def check_accept_rate(self, period: str) -> KillSwitchResult:
        """Post-cycle check: Accept-without-edit rate."""
        query = """
            SELECT 
                COUNT(*) FILTER (WHERE decision='accept' AND edits='{}') AS accept_no_edit,
                COUNT(*) AS total_reviewed
            FROM decisions WHERE period = %s
        """
        result = db.execute(query, [period])
        rate = result.accept_no_edit / result.total_reviewed
        if rate < 0.50:
            return KillSwitchResult(
                triggered=True,
                condition="accept_rate_below_50",
                value=rate,
                action="ALERT_P2",
                notification=["pagerduty:p2", "slack:#tb-review-alerts"]
            )
        return KillSwitchResult(triggered=False)
    
    def check_hallucination_realtime(self, commentary: Commentary) -> KillSwitchResult:
        """Real-time check: Hallucination detection."""
        if commentary.confidence < 0.3 and not commentary.source_references:
            return KillSwitchResult(
                triggered=True,
                condition="hallucination_detected",
                value=commentary.confidence,
                action="KILL_IMMEDIATELY",
                notification=["pagerduty:p1", "kill-switch:activate"]
            )
        return KillSwitchResult(triggered=False)
\end{lstlisting}

% =============================================================================
\section{Commentary Quality Metrics Clarification}
\label{sec:commentary-metrics}
% =============================================================================

\begin{tcolorbox}[colback=orange!5,colframe=orange!70,title=\faIcon{exclamation-circle} Finding P2-2: Metrics Ambiguity]
\textbf{Location:} tb-review-spec, p.15, \S3.8\\
\textbf{Issue:} ``accept-with-edit'' numerator/denominator not defined
\end{tcolorbox}

\textbf{Authoritative Definitions:}

\begin{align}
\text{accept\_without\_edit\_rate} &= \frac{|\{d : d.\text{decision} = \text{`accept'} \land d.\text{edits} = \emptyset\}|}{|\text{total\_reviewed}|} \geq 70\% \\[1em]
\text{accept\_with\_any\_edit\_rate} &= \frac{|\{d : d.\text{decision} = \text{`accept'}\}|}{|\text{total\_reviewed}|} \geq 95\%
\end{align}

\textbf{Interpretation:}
\begin{itemize}
    \item \textbf{accept\_without\_edit\_rate:} Commentaries accepted exactly as AI generated
    \item \textbf{accept\_with\_any\_edit\_rate:} All accepted commentaries (including those requiring human edits)
    \item \textbf{Implication:} Up to 25\% of accepted commentaries may require human editing; 
          up to 5\% may be rejected entirely
\end{itemize}

% =============================================================================
\section{Recommendations Priority Matrix}
\label{sec:priority-matrix}
% =============================================================================

\begin{longtable}{p{1cm} p{5cm} p{3cm} p{2cm} p{1.5cm}}
\caption{Audit Finding Resolution Status} \\
\toprule
\textbf{ID} & \textbf{Action} & \textbf{Owner} & \textbf{Timeline} & \textbf{Status} \\
\midrule
\endfirsthead
\rowcolor{green!10}
P0-1 & Reconcile state count (27→28) & Spec Owner & 1 week & \textcolor{green}{\textbf{RESOLVED}} \\
\rowcolor{green!10}
P0-2 & Add vLLM fallback specification & AI Architect & 2 weeks & \textcolor{green}{\textbf{RESOLVED}} \\
\rowcolor{green!10}
P1-1 & Define audit trail schema & Controls Owner & 4 weeks & \textcolor{green}{\textbf{RESOLVED}} \\
\rowcolor{green!10}
P1-2 & Add normality test to AC & Data Science & 4 weeks & \textcolor{green}{\textbf{RESOLVED}} \\
\rowcolor{yellow!20}
P2-1 & Automate kill-switch eval & Platform Team & 8 weeks & \textcolor{orange}{\textbf{PLANNED}} \\
\rowcolor{green!10}
P2-2 & Define commentary metrics & Product Owner & 2 weeks & \textcolor{green}{\textbf{RESOLVED}} \\
\bottomrule
\end{longtable}

\end{document}
% =============================================================================
% Chapter 15: Implementation Addendum - Implementor & Business Owner Issues
% =============================================================================
\chapter{Implementation Addendum}
\label{ch:implementation-addendum}

This chapter addresses specific concerns raised in the implementor and 
business owner review, providing detailed specifications that were identified 
as missing or insufficiently detailed.

% =============================================================================
\section{Implementor Specifications}
\label{sec:implementor-specs}
% =============================================================================

% -----------------------------------------------------------------------------
\subsection{State Transition Table}
\label{sec:state-transitions}
% -----------------------------------------------------------------------------

The 28-state workflow machine is referenced from 
\texttt{structured/workflows/tb-review-glo.json}. For implementor convenience,
the complete state transition table is embedded here:

\begin{longtable}{p{3.5cm} p{3.5cm} p{5.5cm}}
\caption{Complete State Transition Table (28 States)} \\
\toprule
\textbf{From State} & \textbf{To State} & \textbf{Condition} \\
\midrule
\endfirsthead
\toprule
\textbf{From State} & \textbf{To State} & \textbf{Condition} \\
\midrule
\endhead
START & ROLL\_FORWARD & M-3 trigger \\
ROLL\_FORWARD & DOWNLOAD\_PSGL & Roll-forward complete \\
DOWNLOAD\_PSGL & UPDATE\_ACCOUNTS\_MASTER & Download complete \\
UPDATE\_ACCOUNTS\_MASTER & EXTRACT\_TB & Master updated \\
EXTRACT\_TB & UPDATE\_COMPARATIVES & Extraction complete \\
UPDATE\_COMPARATIVES & CALC\_VARIANCE & Comparatives updated \\
CALC\_VARIANCE & FILTER\_VARIANCES & Calculation complete \\
FILTER\_VARIANCES & UPDATE\_COMMENTARIES & Filtering complete \\
UPDATE\_COMMENTARIES & CHECK\_JOURNAL\_POSTED & Commentaries updated \\
CHECK\_JOURNAL\_POSTED & SEND\_JOURNAL\_REQUEST & journal\_posted = FALSE \\
CHECK\_JOURNAL\_POSTED & UPDATE\_RISK\_SUMMARY & journal\_posted = TRUE \\
SEND\_JOURNAL\_REQUEST & UPDATE\_RISK\_SUMMARY & Request sent \\
UPDATE\_RISK\_SUMMARY & PREPARE\_SUMMARY & Summary updated \\
PREPARE\_SUMMARY & INVESTIGATE\_VARIANCE & variance\_exists = TRUE \\
PREPARE\_SUMMARY & INTERNAL\_REVIEW & variance\_exists = FALSE \\
INVESTIGATE\_VARIANCE & CHECK\_MATERIAL & Investigation complete \\
CHECK\_MATERIAL & CATEGORIZE\_UNEXPLAINED & is\_material = TRUE \\
CHECK\_MATERIAL & INTERNAL\_REVIEW & is\_material = FALSE \\
CATEGORIZE\_UNEXPLAINED & CHECK\_TEAM\_ENTRIES & Categorized \\
CHECK\_TEAM\_ENTRIES & CONFIRM\_ENTRIES & team\_response = PENDING \\
CHECK\_TEAM\_ENTRIES & SEEK\_COMMENTARIES & team\_response = CONFIRMED \\
CONFIRM\_ENTRIES & SEEK\_COMMENTARIES & Entries confirmed \\
CONFIRM\_ENTRIES & AWAITING\_CLARIFICATION & timeout \textgreater 48h \\
SEEK\_COMMENTARIES & RECEIVE\_COMMENTARIES & Commentary requested \\
RECEIVE\_COMMENTARIES & PERFORM\_RISK\_ANALYSIS & Commentary received \\
RECEIVE\_COMMENTARIES & ESCALATED\_TO\_FINANCE & commentary\_rejected \\
PERFORM\_RISK\_ANALYSIS & UPDATE\_VARIANCE\_COMMENTARIES & Risk analyzed \\
UPDATE\_VARIANCE\_COMMENTARIES & TRACK\_ENTRIES & Variance commented \\
TRACK\_ENTRIES & INTERNAL\_REVIEW & Entries tracked \\
INTERNAL\_REVIEW & SEND\_FORTM\_SUMMARY & Review approved \\
INTERNAL\_REVIEW & MANUAL\_OVERRIDE & override\_required = TRUE \\
AWAITING\_CLARIFICATION & CHECK\_TEAM\_ENTRIES & clarification\_received \\
ESCALATED\_TO\_FINANCE & PERFORM\_RISK\_ANALYSIS & finance\_resolution \\
MANUAL\_OVERRIDE & INTERNAL\_REVIEW & override\_applied \\
SEND\_FORTM\_SUMMARY & COMPLETED & M+5 \\
\bottomrule
\end{longtable}

\textbf{State Machine Implementation Notes:}
\begin{itemize}
    \item Error states (25-27) can be reached from any active state via timeout or explicit escalation
    \item All transitions are logged to audit trail with timestamp and persona\_id
    \item State persistence: PostgreSQL \texttt{workflow\_states} table with optimistic locking
\end{itemize}

% -----------------------------------------------------------------------------
\subsection{LLM Service Contract}
\label{sec:llm-contract}
% -----------------------------------------------------------------------------

\begin{tcolorbox}[colback=blue!5,colframe=blue!70,title=LLM Service Contract: vLLM Commentary Generation]

\textbf{Endpoint:}
\begin{verbatim}
POST https://ai-core.internal.sap/v1/tb-commentary/generate
\end{verbatim}

\textbf{Authentication:}
\begin{itemize}
    \item Same SAP IAS OAuth 2.0 Bearer token as dashboard API
    \item Token claim \texttt{scope} must include \texttt{tb:commentary:generate}
\end{itemize}

\textbf{Request:}
\begin{lstlisting}[language=json]
{
  "request_id": "uuid-v4",
  "variance_record": {
    "account_code": "4100001",
    "account_name": "Revenue - Product Sales",
    "legal_entity": "HKG",
    "period": "2026-03",
    "prior_balance": 1500000.00,
    "current_balance": 1875000.00,
    "variance_amount": 375000.00,
    "variance_percentage": 0.25,
    "variance_type": "spike",
    "related_transactions": [
      {"date": "2026-03-15", "amount": 200000, "description": "Contract #A123"}
    ]
  },
  "context": {
    "historical_periods": 12,
    "include_trend_analysis": true
  },
  "preferences": {
    "tone": "formal",
    "max_length": 500,
    "language": "en"
  }
}
\end{lstlisting}

\textbf{Response (Success - 200):}
\begin{lstlisting}[language=json]
{
  "commentary_id": "uuid-v4",
  "text": "Revenue from Product Sales increased by $375,000 (25%) 
          compared to prior month. This spike is primarily attributed 
          to Contract #A123 signed on March 15, contributing $200,000 
          to the variance. Historical analysis indicates this is the 
          highest monthly increase in the past 12 months.",
  "confidence": 0.87,
  "source_references": ["Contract #A123", "GL Entry 2026-03-15"],
  "model_version": "tb-commentary-v1.2",
  "processing_time_ms": 1250,
  "fallback_used": false
}
\end{lstlisting}

\textbf{Error Codes:}
\begin{tabular}{lll}
\toprule
\textbf{Code} & \textbf{Meaning} & \textbf{Client Action} \\
\midrule
400 & Invalid variance\_record schema & Fix request body \\
401 & Token missing or expired & Re-authenticate \\
403 & Insufficient scope & Check token claims \\
429 & Rate limit exceeded & Backoff and retry \\
500 & Model error & Fallback to template \\
503 & vLLM unavailable & Fallback to template \\
504 & Timeout (>5s) & Fallback to template \\
\bottomrule
\end{tabular}

\textbf{Timeout and Retry Configuration:}
\begin{itemize}
    \item Connect timeout: 1000ms
    \item Read timeout: 5000ms
    \item Max retries: 2 (with exponential backoff)
    \item Circuit breaker: Open after 5 consecutive failures (30s reset)
\end{itemize}
\end{tcolorbox}

% -----------------------------------------------------------------------------
\subsection{Bulk Operation Semantics Clarification}
\label{sec:bulk-semantics}
% -----------------------------------------------------------------------------

\begin{tcolorbox}[colback=green!5,colframe=green!70,title=Bulk Operation Specification]

\textbf{Transaction Semantics:}

Bulk operations are processed \textbf{per-item with individual audit records}. 
Failures do NOT roll back successful items.

\textbf{HTTP Response Codes:}
\begin{itemize}
    \item \textbf{200 OK:} All items succeeded
    \item \textbf{207 Multi-Status:} Partial success; response body contains per-item status
    \item \textbf{409 Conflict:} All items failed due to race condition
\end{itemize}

\textbf{207 Multi-Status Response Example:}
\begin{lstlisting}[language=json]
{
  "bulk_operation_id": "uuid-bulk-123",
  "summary": {
    "total": 50,
    "success": 48,
    "failed": 2
  },
  "results": [
    {"id": "var-001", "status": 200, "message": "Assigned"},
    {"id": "var-002", "status": 409, "message": "Already claimed by 1234567"},
    ...
  ]
}
\end{lstlisting}

\textbf{Locking Semantics:}
\begin{itemize}
    \item Each item is claimed with optimistic locking (\texttt{version} field)
    \item Two reviewers claiming overlapping sets: first claim wins; second gets 409
    \item UI should refresh and show updated claimant on 409
\end{itemize}
\end{tcolorbox}

% -----------------------------------------------------------------------------
\subsection{Dashboard UI Technology Stack}
\label{sec:ui-stack}
% -----------------------------------------------------------------------------

\begin{tabular}{ll}
\toprule
\textbf{Component} & \textbf{Technology} \\
\midrule
Framework & Angular 17+ \\
State Management & NgRx (Redux pattern) \\
Component Library & SAP UI5 Web Components for Angular \\
HTTP Client & Auto-generated from OpenAPI spec (\texttt{ng-openapi-gen}) \\
Testing & Jest (unit), Cypress (E2E) \\
Build & Nx monorepo with \texttt{@nrwl/angular} \\
Deployment & Docker container on SAP BTP Cloud Foundry \\
\bottomrule
\end{tabular}

\textbf{Reference:}
\begin{verbatim}
src/generativeUI/ui5-webcomponents-ngx-main/
\end{verbatim}

% =============================================================================
\section{Business Owner Specifications}
\label{sec:business-specs}
% =============================================================================

% -----------------------------------------------------------------------------
\subsection{G-001: Preliminary Financial Analysis}
\label{sec:g001-financial}
% -----------------------------------------------------------------------------

\begin{tcolorbox}[colback=yellow!10,colframe=yellow!50!black,title=\faIcon{exclamation-triangle} G-001: Financial Analysis - Preliminary Estimates]

\textbf{Note:} Final financial analysis is pending (owner: S1 Badrinath, target: Q2 2026).
The following are \textbf{preliminary estimates} with $\pm 30\%$ confidence band.

\vspace{0.5em}
\textbf{Cost Estimate (Year 1):}
\begin{tabular}{lrr}
\toprule
\textbf{Category} & \textbf{Estimate} & \textbf{Range ($\pm 30\%$)} \\
\midrule
AI Architect (100 days) & \$150,000 & \$105,000 - \$195,000 \\
AI Engineer (100 days) & \$100,000 & \$70,000 - \$130,000 \\
Infrastructure (AI Core) & \$50,000 & \$35,000 - \$65,000 \\
SME/Testing (15 days $\times$ 4) & \$45,000 & \$31,500 - \$58,500 \\
\midrule
\textbf{Total Year 1} & \textbf{\$345,000} & \$241,500 - \$448,500 \\
\bottomrule
\end{tabular}

\textbf{Benefit Estimate (Year 1):}
\begin{tabular}{lrr}
\toprule
\textbf{Category} & \textbf{Estimate} & \textbf{Range} \\
\midrule
FTE savings (5 FTE $\times$ \$80K) & \$400,000 & \$280,000 - \$520,000 \\
Reduced audit costs & \$50,000 & \$35,000 - \$65,000 \\
Fast close enablement & \$25,000 & \$17,500 - \$32,500 \\
\midrule
\textbf{Total Year 1 Benefits} & \textbf{\$475,000} & \$332,500 - \$617,500 \\
\bottomrule
\end{tabular}

\textbf{Preliminary ROI:}
\begin{align*}
\text{ROI} &= \frac{\text{Benefits} - \text{Cost}}{\text{Cost}} \times 100\% \\
&= \frac{\$475,000 - \$345,000}{\$345,000} \times 100\% \\
&= \textbf{37.7\%} \text{ (range: -27\% to +156\%)}
\end{align*}

\textbf{Payback Period:} $\approx$ 9 months (range: 6--18 months)

\vspace{0.5em}
\textbf{Action:} Final analysis required before GA cutover. This estimate is 
for planning purposes only.
\end{tcolorbox}

% -----------------------------------------------------------------------------
\subsection{Cross-Spec Governance}
\label{sec:cross-spec-governance}
% -----------------------------------------------------------------------------

\begin{tcolorbox}[colback=blue!5,colframe=blue!70,title=Cross-Spec Governance Protocol]

Shared enums are maintained in:
\begin{verbatim}
docs/schema/common/enums.yaml
\end{verbatim}

\textbf{Change Process:}
\begin{enumerate}
    \item RFC (Request for Comment) submitted to Controls Architecture Team
    \item All consuming spec owners (TB, BSS, Arabic AP) are notified
    \item 5-day review period; majority approval required
    \item One-release deprecation window for breaking changes
    \item CI validation blocks PRs with local enum redefinitions
\end{enumerate}

\textbf{Ownership:}
\begin{tabular}{lll}
\toprule
\textbf{Enum} & \textbf{Maintainer} & \textbf{Consumers} \\
\midrule
risk\_status & Controls Architecture & TB, BSS \\
decision\_outcome & Workflow Architecture & TB, Arabic, Regulations \\
variance\_type & Financial Analytics & TB \\
audit\_action & Controls Architecture & TB, BSS, Arabic \\
\bottomrule
\end{tabular}
\end{tcolorbox}

% -----------------------------------------------------------------------------
\subsection{ML Timeline Caveat}
\label{sec:ml-caveat}
% -----------------------------------------------------------------------------

\begin{tcolorbox}[colback=red!10,colframe=red!70,title=\faIcon{exclamation-triangle} ML Timeline Caveat - Important]

\textbf{The ML-based anomaly detection timeline (Chapter 10) is conditional 
and illustrative, NOT committed.}

\textbf{Dependencies:}
\begin{itemize}
    \item No ML model exists today
    \item No labelled anomaly corpus exists
    \item Requires $\geq 2$ quarters of production dashboard decisions to build corpus
    \item Isolation Forest prototype target (Q1 2027) depends on data availability
\end{itemize}

\textbf{Business Implication:}
\begin{itemize}
    \item v1.0 business case should NOT include ML-driven benefits
    \item ML savings are Phase 2+ benefits only
    \item Missing ML timeline does NOT block v1.0 release
    \item Steering committee should be briefed on this caveat
\end{itemize}

\textbf{Data Collection Trigger:} ML labelling begins when:
\begin{enumerate}
    \item Dashboard has been in production for $\geq 2$ months
    \item At least 1,000 accept/reject decisions have been recorded
    \item Decision quality metrics (accept rate $\geq 70\%$) are stable
\end{enumerate}
\end{tcolorbox}

% -----------------------------------------------------------------------------
\subsection{Commentary Quality Metrics Remediation}
\label{sec:quality-remediation}
% -----------------------------------------------------------------------------

\begin{tcolorbox}[colback=orange!5,colframe=orange!70,title=Quality Metrics Remediation Protocol]

\textbf{Targets:}
\begin{itemize}
    \item \texttt{accept\_without\_edit\_rate} $\geq 70\%$
    \item \texttt{accept\_with\_any\_edit\_rate} $\geq 95\%$
\end{itemize}

\textbf{Remediation Trigger:}

If \textbf{either} target is missed for \textbf{two consecutive months}:

\begin{enumerate}
    \item Root-cause analysis is triggered (owner: AI Architect)
    \item Commentary generation is switched to \textbf{template-only mode}
    \item Model retuning initiated with expanded training data
    \item LLM-based generation resumes only after pilot shows target compliance
\end{enumerate}

\textbf{Escalation Path:}
\begin{itemize}
    \item Month 1 miss: Warning alert to AI team
    \item Month 2 miss: Automatic switch to template mode + Jira escalation
    \item Month 3+ template mode: Steering committee briefing required
\end{itemize}

\textbf{Recovery Criteria:}

Return to LLM mode when:
\begin{itemize}
    \item Pilot (50 variances) achieves $\geq 70\%$ accept-without-edit
    \item No factual errors detected in pilot sample
    \item AI Architect sign-off obtained
\end{itemize}
\end{tcolorbox}

% =============================================================================
\section{Implementation Priority Order}
\label{sec:impl-priority}
% =============================================================================

Based on the implementor review, the recommended priority order is:

\begin{enumerate}
    \item \textbf{Set up unified schema registry and validation pipeline}
    \begin{itemize}
        \item Clone \texttt{docs/schema/} to development environment
        \item Integrate \texttt{validate.py --registry-check} into CI
        \item Generate TypeScript types from JSON Schema
    \end{itemize}
    
    \item \textbf{Implement entity-params loader}
    \begin{itemize}
        \item Create \texttt{EntityParamsRegistry.load\_all()} class
        \item Implement per-LE materiality threshold lookup
        \item Add graceful missing-LE handling
    \end{itemize}
    
    \item \textbf{Build anomaly detection with Shapiro-Wilk + MAD fallback}
    \begin{itemize}
        \item Implement \texttt{AnomalyDetector} class (see P1-2 resolution)
        \item Add normality test before Z-score application
        \item Configure thresholds in entity-params
    \end{itemize}
    
    \item \textbf{Implement vLLM commentary with template fallback}
    \begin{itemize}
        \item Use LLM Service Contract (Section~\ref{sec:llm-contract})
        \item Implement \texttt{CommentaryService} with fallback logic
        \item Create 7 template variants (TPL-001 to TPL-007)
    \end{itemize}
    
    \item \textbf{Build dashboard queue and review views}
    \begin{itemize}
        \item Follow UI stack (Angular + NgRx + SAP UI5)
        \item Auto-generate API client from OpenAPI spec
        \item Implement audit trail logging for all actions
    \end{itemize}
\end{enumerate}

\textbf{Dependencies to Request from Architecture Team:}
\begin{itemize}
    \item Full \texttt{tb-review-glo.json} workflow file (Section~\ref{sec:state-transitions} provides summary)
    \item vLLM API endpoint and credentials (Section~\ref{sec:llm-contract} provides contract)
    \item SAP IAS OAuth configuration for development environment
\end{itemize}

\end{document}

% --- Incorporated Universal Chapters ---
\chapter{Global Wave-Dependency Manifest}
\label{ch:global-dependencies}

\section{Critical Path Visualization (Waves 1--3)}
The following chart defines the normative sequencing for implementation agents. A task in a later wave cannot be marked \textsc{complete} until its prerequisites in an earlier wave are validated.

\begin{figure}[htbp]
\centering
\begin{tikzpicture}[
    node distance=1.5cm,
    auto,
    wavebox/.style={rectangle, draw, rounded corners, fill=sapmidnight!10, text width=10em, text centered, minimum height=3em, font=\scriptsize\bfseries},
    dep/.style={->, >=latex, thick, sapblue},
    critical/.style={->, >=latex, very thick, sapamber}
]
    % --- Wave 1: Foundations ---
    \node (r1) [wavebox, fill=sapgrey!20] {\textcolor{sapmidnight}{REG-W1: Identity \& Audit Schema}};
    \node (s1) [wavebox, below of=r1, yshift=-0.5cm] {SIM-W1: HANA Extraction Pipeline};
    \node (a1) [wavebox, below of=s1, yshift=-0.5cm] {ARA-W1: Invoice Intake \& OCR};
    \node (t1) [wavebox, below of=a1, yshift=-0.5cm] {TB-W1: Identity Wiring};

    % --- Wave 2: Core Build-out ---
    \node (r2) [wavebox, right of=r1, xshift=4.5cm, fill=sapmidnight!20] {\textcolor{sapmidnight}{REG-W2: Governance Runtime}};
    \node (s2) [wavebox, right of=s1, xshift=4.5cm] {SIM-W2: Taxonomy \& Synthesis};
    \node (a2) [wavebox, right of=a1, xshift=4.5cm] {ARA-W2: Multi-Country VAT};
    \node (t2) [wavebox, right of=t1, xshift=4.5cm] {TB-W2: Controls \& Workflow};

    % --- Wave 3: Hardening ---
    \node (r3) [wavebox, right of=r2, xshift=4.5cm, fill=sapmidnight!30] {\textcolor{sapmidnight}{REG-W3: Conformance Suite}};
    \node (s3) [wavebox, right of=s2, xshift=4.5cm] {SIM-W3: 85\% Coverage Test};
    \node (a3) [wavebox, right of=a2, xshift=4.5cm] {ARA-W3: Calibration Pilot};
    \node (t3) [wavebox, right of=t2, xshift=4.5cm] {TB-W3: HITL Dashboard};

    % --- Cross-Domain Dependencies (The Critical Path) ---
    \draw [critical] (r1) -- node [above, font=\tiny, align=center] {Identity Envelope\\Invariants} (a1.north east);
    \draw [critical] (r1) -- (t1.north east);
    \draw [dep] (s1.east) -- node [above, font=\tiny, xshift=-0.5cm] {Training Fixtures} (t2.west);
    \draw [dep] (s1.east) -- (a2.west);
    \draw [dep] (r2.south) -- node [left, font=\tiny, yshift=0.3cm] {Policy Gate} (a2.north);
    \draw [dep] (t2.east) -- node [above, font=\tiny] {Audit Evidence} (r3.west);
    
    % --- Wave Transitions ---
    \draw [->, dashed, sapgrey] (r1.east) -- (r2.west);
    \draw [->, dashed, sapgrey] (r2.east) -- (r3.west);

    % --- Legend ---
    \node[draw, sapamber, very thick, label=right:\scriptsize Critical Prerequisite] at (0,-7.5) {};
    \node[draw, sapblue, thick, label=right:\scriptsize Data Dependency, xshift=4cm] at (0,-7.5) {};

\end{tikzpicture}
\caption{Cross-specification dependency graph: The blueprint for platform rollout.}
\end{figure}

\section{Inter-Domain Prerequisite Matrix}

\begin{table}[htbp]
\centering
\footnotesize
\caption{Normative Cross-Domain Prerequisites.}
\label{tab:prerequisites}
\begin{tabularx}{\textwidth}{@{}l l l X@{}}
\toprule
\textbf{Target Task} & \textbf{Prerequisite} & \textbf{Ref ID} & \textbf{Rationale} \\
\midrule
\rowcolor{sapmidnight!5}
\textbf{ARA-W1} & REG-W1 (Identity) & \texttt{REG-IDENTITY-01} & Invoice intake must be audit-traced from Day 1. \\
\textbf{TB-W2}  & SIM-W1 (Extraction) & \texttt{SIM-EXTRACT-01} & Anomaly detection requires HANA schema metadata. \\
\rowcolor{sapmidnight!5}
\textbf{REG-W2} & SIM-W1 (Registry) & \texttt{SIM-REGISTRY-01} & Governance monitoring relies on the schema registry. \\
\textbf{ARA-W2} & REG-W2 (Pillar 1) & \texttt{REG-MGF-P1} & VAT release is an HITL gate governed by MAS Pillar 1. \\
\rowcolor{sapmidnight!5}
\textbf{TB-W3}  & ARA-W2 (Posted Invoices) & \texttt{ARA-AP-POST} & Review dashboard depends on posted AP data from Arabic domain. \\
\bottomrule
\end{tabularx}
\end{table}

\begin{adrbox}{Dependency Enforcement}
Implementation agents are strictly forbidden from "Fast-Forwarding" into Wave 2 if Wave 1 critical prerequisites (marked in \textcolor{sapamber}{\textbf{Gold}}) are in a \textsc{gap} or \textsc{partial} state.
\end{adrbox}





\end{document}
